\chapter{Zhu's bijection via the tower \texorpdfstring{$A_n(V)$}{A\_n(V)}}\label{ch:L13}
\begingroup\setlength{\emergencystretch}{3em}

\providecommand{\lthirteenvac}{\mathbf{1}}
\providecommand{\lthirteenVh}{\widehat{V}}
\providecommand{\lthirteenUE}{\mathcal{U}}
\providecommand{\lthirteenMbar}{\overline{M}}
\providecommand{\lthirteenIrr}{\mathrm{Irr}}
\providecommand{\lthirteenO}{O}

This chapter documents the release of the LeanVOA library named in the front matter; it was first written against an earlier state of the library (``the snapshot'' below), and checks that were
run at an earlier state of the library and not re-run for this edition are marked as such. Every Lean quotation below is taken mechanically from the files of the release and is machine-checked by
\texttt{scripts/textbook/sync.py} to occur verbatim (up to whitespace and comments) in the source. The chapter covers
the $298$ files \leanpath{VOA/\allowbreak Modules/\allowbreak ZhuVerma*.lean} (about ninety thousand lines of Lean)
assigned to it in \texttt{Notes/Textbook/assignment.txt}. They contain one headline: the \emph{hard direction of
Zhu's correspondence} between irreducible modules of a vertex operator algebra $V$ and irreducible modules of its
Zhu algebra $A(V)$. The easy direction, the algebra $A(V)$ itself, the bundled module types \lean{AMod} and
\lean{LWModule}, and the definition of the library's higher Zhu spaces $A_n(V)$ were presented in
Chapter~\ref{ch:L12}; we use them freely and point back to them.

We do not tour $298$ files. Instead we (i) state exactly what the headline theorem says, (ii) present every
definition on which the statement and its proof depend, in the order in which the proof uses them, (iii) explain the
proof on paper, and (iv) describe the remaining files by role in one short section
(Section~\ref{l13:sec:files}). The ``critical path'' was not taken from the library's own documentation: we computed it
by a machine walk through the constants used by the proof term of the headline (Section~\ref{l13:sec:files}); it
touches $101$ of the $298$ files.

\emph{Release revision.} The chapter was revised for the first public release, applying the release audit \texttt{Notes/\allowbreak Textbook/\allowbreak referee/\allowbreak
RELEASE\_AUDIT\_ZHU.md} (2026-10-27). Two results proved outside this chapter's files since the snapshot change what can be
said about the full correspondence, and are presented in Section~\ref{l13:sec:scope}: the bottom of every irreducible module
is a simple $A(V)$-module, for every $V$ (\leanpath{VOA/\allowbreak Huang/\allowbreak LWBottomRecovery.lean}, Chapter~\ref{ch:L21}); and
the combined statement \lean{zhuCorrespondence_of_C2Cofinite}, a genuine bijection between $\Phi(V)$ and isomorphism classes of
simple $A(V)$-modules for $C_2$-cofinite $V$ (\leanpath{VOA/\allowbreak Modularity/\allowbreak ZhuModularityCanonical.lean},
Chapter~\ref{ch:L15}). Their quotations were copied at that revision.

\emph{Second release revision (2026-10-28).} At this revision one more consumer of this chapter's machinery is
recorded (finding F2 of \texttt{Notes/\allowbreak Textbook/\allowbreak referee/\allowbreak RELEASE\_AUDIT\_NEW\_DECLS.md}). The
generalised Verma module induced from the regular $A(V)$-module, together with the datum \lean{Pr} and bottom recovery, also
proves Zhu's semisimplicity theorem ``rational $\Rightarrow A(V)$ semisimple''
(\lean{isSemisimpleRing_}\allowbreak\lean{A_of_}\allowbreak\lean{admissiblyRational}, \leanpath{VOA/\allowbreak Huang/\allowbreak
ZhuSemisimple.lean}; presented in Chapter~\ref{ch:L12}, Theorem~\ref{l12:thm:zhuss}, and Chapter~\ref{ch:L15},
Theorem~\ref{l15:thm:zhuss}). So for a $C_2$-cofinite rational $V$ (rational in the library's sense \lean{AdmissiblyRational})
the simple $A(V)$-modules of Theorem~\ref{l13:thm:zhucorr} are exactly the simple modules of a semisimple algebra, and Zhu's
modularity theorem holds for $V$ (Chapter~\ref{ch:L15}). No statement or quotation of this chapter changed.

\paragraph{The findings of this chapter, in brief.} The reader who only wants the verdict should read this paragraph
and the warning boxes.
\begin{itemize}
\item \textbf{What is proved.} The headline \lean{zhuBijective} states that a certain map
  $[M]\mapsto[M_h]$ is a bijection. Its domain is the set of isomorphism classes of \emph{irreducible ordinary}
  $V$-modules (with their lowest weight); its codomain is the set of isomorphism classes of $A(V)$-modules
  \emph{that are isomorphic to the bottom of some irreducible ordinary module}. Because the codomain is defined as the
  image, surjectivity is a tautology (the library says so itself). The genuine content is \emph{injectivity}:
  \begin{quote}
  two irreducible ordinary $V$-modules whose lowest-weight spaces are isomorphic $A(V)$-modules are isomorphic.
  \end{quote}
  This is Zhu's uniqueness theorem, and it is proved for \emph{every} vertex operator algebra of the library, with one
  explicit datum \lean{Pr}. It is the \emph{uniqueness half} of Zhu's correspondence, and must be announced as such, not
  as ``Zhu's bijection''. The equivalence ``\lean{zhuBijective} $\iff$ uniqueness'' was checked in Lean against the
  library (the check is not part of this release).
\item \textbf{The datum \lean{Pr}.} \lean{Pr : BottomProjection (parabolicRepOfAMod W (regularAMod W))} is a single
  linear-algebra statement attached to $V$: in the module induced from the \emph{regular} $A(V)$-module along a
  parabolic subalgebra of the Frenkel--Zhu Lie algebra $\lthirteenVh$, the canonical bottom copy of $A(V)$ is
  embedded. The library proves that a \lean{BottomProjection} exists if and only if the bottom map is injective. It is a
  consequence of the Poincar\'e--Birkhoff--Witt theorem, hence true mathematics, and it is proved in the library
  for every $V$: \lean{voaPBWBottomProjection} builds it from the theorem \lean{VOA.pbw} (relative PBW,
  Chapter~\ref{ch:L02}, Section~\ref{l02:sec:pbwproof}). The wrapper \lean{zhuBijective_of_voaPBW} uses it and is
  unconditional; its \texttt{\#print axioms} is the standard three (at the snapshot, when
  \lean{VOA.pbw} was an axiom, it contained \lean{pbw}). So Zhu's bijection in its uniqueness form is now
  unconditional.
\item \textbf{The genuine bijection, for $C_2$-cofinite $V$ (new).} The two halves of Zhu's correspondence that the
  headline does not contain are now theorems: (1) the bottom of every irreducible module is a simple $A(V)$-module, for
  \emph{every} $V$ (\lean{allBottomRecovers}, \lean{isSimpleAMod_toAMod_of_irreducible}); (2) every simple $A(V)$-module is
  the bottom of an irreducible ordinary module, for $C_2$-cofinite $V$ (\lean{simpleAModsRealized_of_C2Cofinite}). Combined
  with the headline they give \lean{zhuCorrespondence_of_C2Cofinite} (Theorem~\ref{l13:thm:zhucorr}): for every $C_2$-cofinite
  $V$, $[M]\mapsto[M_h]$ is a bijection from $\Phi(V)$ onto the isomorphism classes of \emph{simple} $A(V)$-modules, where
  ``simple'' is the intrinsic notion, not an image. $C_2$ is used only for existence (2); Zhu's own existence theorem does
  not need it, because it produces admissible rather than ordinary modules.
\item \textbf{What is not proved.} Existence (2) without $C_2$, i.e.\ for $V$ whose $L(U)$ might have infinite-dimensional
  weight spaces; there the bijection is stated only onto its image (Section~\ref{l13:sec:scope}).
\item \textbf{A misnamed file.} The file \leanpath{VOA/\allowbreak Modules/\allowbreak ZhuVermaAnArchCritPath.lean}
  (``the critical path of the Zhu bijection'') is \emph{not} on the critical path of \lean{zhuBijective}, and its
  headline \lean{zhuBijection_critpath} is vacuous: its hypothesis bundle cannot be inhabited (checked in Lean; the check is
  not part of this release).
\item \textbf{The tower.} The library's circle-only $O_n(V)$ (\lean{On}, which omits Dong--Li--Mason's
  $(L_{-1}+L_0)a$ for $n\ge1$, Chapter~\ref{ch:L12}) is used on the critical path only at $n = 0$, where it equals
  $O(V)$. The construction of the induced module uses the corrected \lean{OnFull} and Dong--Jiang's full ideals
  \lean{OnmFull}; so that defect does not touch the headline.
\item \textbf{A redundant hypothesis, since removed.} \lean{pHL0} used to carry the datum \lean{Pr} without using it
  (checked in Lean; the check is not part of this release); the library has dropped it, and \lean{zhuBijective} now applies the hypothesis-free \lean{pHL0} directly.
\end{itemize}

%======================================================================
\section{Zhu's correspondence: the mathematics and the plan}\label{l13:sec:plan}
%======================================================================

\subsection{The statement in the literature}

Let $V$ be a vertex operator algebra. Recall from Chapter~\ref{ch:L12} (Section~\ref{l12:sec:zhu}) that Zhu's algebra is
$A(V) = V/O(V)$ with product $[u][v] = [u*v]$, where for homogeneous $u$
\begin{equation*}
 u*v = \Res_z Y(u,z)v\,\frac{(1+z)^{\wt u}}{z},\qquad u\circ v = \Res_z Y(u,z)v\,\frac{(1+z)^{\wt u}}{z^{2}},
 \qquad O(V) = \operatorname{span}\{u\circ v\}.
\end{equation*}
If $M = \bigoplus_{n\ge0}M_{h+n}$ is a $V$-module whose lowest weight $h$ is attained ($M_h\neq0$), then the zero modes
$o(u) = u_{(\wt u-1)}$ preserve $M_h$, and $[u]\mapsto o(u)|_{M_h}$ makes $M_h$ an $A(V)$-module
(Theorem~\ref{l12:thm:moduleMh}). Zhu's theorem says much more.

\begin{theorem}[Zhu 1996, Theorems~2.2.1 and~2.2.2]\label{l13:thm:zhu-lit}
Let $V$ be a vertex operator algebra.
\begin{enumerate}
\item If $M$ is an irreducible $\N$-gradable (admissible) $V$-module with lowest piece $M(0)$, then $M(0)$ is a simple
  $A(V)$-module.
\item For every simple $A(V)$-module $U$ there is an irreducible admissible $V$-module $L(U)$ with $L(U)(0)\cong U$.
\item The assignments $M\mapsto M(0)$ and $U\mapsto L(U)$ are mutually inverse bijections between isomorphism classes
  of irreducible admissible $V$-modules and isomorphism classes of simple $A(V)$-modules.
\end{enumerate}
\end{theorem}
Restricting to \emph{ordinary} modules (finite-dimensional $L_0$-eigenspaces), one gets a bijection between irreducible
ordinary modules and those simple $A(V)$-modules $U$ for which $L(U)$ has finite-dimensional weight spaces; when $V$
is $C_2$-cofinite these are all simple $A(V)$-modules (finite-dimensionality of $A(V)$ alone does not control the higher
weight spaces of $L(U)$).

Part~(3) contains two assertions of very different difficulty. \emph{Well-definedness and surjectivity onto the simple
modules} need the existence of $L(U)$, part~(2). \emph{Injectivity} --- an irreducible module is determined by its
bottom --- needs a universal object: a ``generalized Verma module'' $M(U)$ generated by $U$ that maps onto every module
whose bottom contains $U$. Both directions therefore require the same construction: given an $A(V)$-module $U$, build a
$V$-module $M(U)$ whose bottom is \emph{exactly} $U$.

\subsection{Why the construction is hard, and the tower}

To build $M(U)$ one must know what the excited levels $M(U)(n)$, $n\ge1$, are. The vectors of $M(U)(n)$ are formally
$a_{(p)}w$ with $w\in U$ and $\wt a - p - 1 = n$, subject to every relation that holds in every $V$-module. Frenkel and
Zhu found that these relations are encoded by quotients of $V$ itself: Dong, Li and Mason introduced the associative
algebras $A_n(V) = V/O_n(V)$ (``higher Zhu algebras'') acting on the first $n$ levels of any module, and Dong and Jiang
introduced $A_{n,m}(V) = V/O_{n,m}(V)$, an $A_n(V)$--$A_m(V)$-bimodule which maps level $m$ to level $n$. Their
construction is
\begin{equation}\label{l13:eq:MU}
 M(U) = \bigoplus_{n\ge0} A_{n,0}(V)\otimes_{A(V)}U,\qquad M(U)(n) = A_{n,0}(V)\otimes_{A(V)}U,
\end{equation}
with $a_{(p)}$ acting from level $n$ to level $\wt a-p-1+n$ by a residue product. The irreducible module is then
$L(U) = M(U)/J(U)$, $J(U)$ the largest graded submodule meeting the bottom trivially.

The delicate point is the bottom: one needs $M(U)(0) = A_{0,0}(V)\otimes_{A(V)}U\cong U$, i.e.\ $A_{0,0}(V) = A(V)$.
Dong and Jiang prove this \emph{representation-theoretically}: they realize the regular $A(V)$-module $A(V)$ as the
bottom of some module (by the induced-module construction of Dong--Li--Mason, which uses the
Poincar\'e--Birkhoff--Witt theorem) and read off $O_{0,0}(V)\subseteq O(V)$. This is precisely where the library's
datum \lean{Pr} enters.

\subsection{The plan of the formal proof}

The proof of \lean{zhuBijective} in the library is the following chain; each arrow is a section of this chapter.
\begin{enumerate}
\item \emph{The statement} (Section~\ref{l13:sec:statement}): the faithful forward map \lean{forwardMapIrr}; its
  bijectivity is equivalent to \emph{bottom rigidity} of irreducibles.
\item \emph{The Lie algebra $\lthirteenVh$ and induction} (Section~\ref{l13:sec:vhat}): the Frenkel--Zhu Lie algebra
  $\lthirteenVh = V\otimes\C[t,t^{-1}]/\operatorname{Im}(L_{-1}\otimes1+1\otimes\tfrac{d}{dt})$, its parabolic
  subalgebra $P_0$, the $P_0$-module structure on any $A(V)$-module, the induced module
  $\lthirteenUE(\lthirteenVh)\otimes_{\lthirteenUE(P_0)}U$, and the datum \lean{BottomProjection}. \lean{VOA.pbw} (an axiom at
  the snapshot, now a theorem) enters only here.
\item \emph{The tower} (Section~\ref{l13:sec:tower}): circle products, $O_n$, $A_n$, and the Dong--Jiang ideals
  $O_{n,m}$ and spaces $A_{n,m}$.
\item \emph{The induced module $M(U)$} (Section~\ref{l13:sec:MU}): the carrier~\eqref{l13:eq:MU}, its $V$-action (a weak
  module, unconditionally), and the ``gate'' $A_{0,0}(V) = A(V)$, proved from \lean{Pr} by realizing the regular module.
\item \emph{$L(U)$ and the injectivity argument} (Section~\ref{l13:sec:LU}): $J(U)$, $L(U)$, the universal map
  $M(\Omega(P))\to P$ for an irreducible $P$ with bottom $\Omega(P)$, the isomorphism $P\cong L(\Omega(P))$, and
  functoriality $U_1\cong U_2\Rightarrow L(U_1)\cong L(U_2)$. Together: irreducibles are determined by their bottoms.
\end{enumerate}

\begin{physics}
In a two-dimensional chiral conformal field theory with chiral algebra $V$, an irreducible module is a
\emph{superselection sector}, and its lowest-weight space $M_h$ is the space of \emph{primary states} (ground states of
the sector on the cylinder, all of energy $h-c/24$). The zero modes $o(u)$ of the fields are exactly the operators that
preserve energy, so they act on the primaries; $A(V)$ is the algebra they generate, with all the relations imposed by
the operator product expansion. Zhu's correspondence is the precise form of the physicist's slogan ``a sector is
determined by its primary multiplet'': every sector is generated from its primaries by the negative modes (the
descendants), and the only relations among descendants are those forced by the chiral algebra itself and by
\emph{null states}. The induced module $M(U)$ is the ``Verma module'' of all descendants of the primary multiplet $U$
with no null states removed; $J(U)$ is the space of null descendants (states orthogonal to everything); $L(U)$ is the
physical sector. The higher algebras $A_n(V)$ and $A_{n,m}(V)$ describe the fields as maps between \emph{excited}
levels: $A_{n,0}(V)\otimes U$ is the space of level-$n$ descendants of $U$.
\end{physics}

%======================================================================
\section{The statement that is proved}\label{l13:sec:statement}
%======================================================================

We begin at the top, with the theorem and the definitions that occur in its statement. The bundled types
\lean{AMod W} (an $A(V)$-module: a vector space $U$ with a linear $\rho:A(V)\to\End U$, unital and multiplicative),
\lean{AModEquiv W U U'} (a linear isomorphism intertwining $\rho$), \lean{LWModule W} (an ordinary module
\lean{isMod : OrdVModule W carrier} with a number \lean{h} such that all weights lie in $h+\N$ and $M_h\neq0$) and
\lean{LWModule.toAMod} (the bottom $M_h$ with $[u]\mapsto o(u)$) are Definition~\ref{l12:def:lwmodule} of
Chapter~\ref{ch:L12}; \lean{OrdVModule} is the library's ordinary module (Chapter~\ref{ch:01}): a weak module with an
internal direct-sum decomposition into $L_0$-eigenspaces $M_\alpha$, $\alpha\in\C$, each finite-dimensional, with real
parts bounded below.

\subsection{Irreducibility and isomorphism}

\begin{definition}[Irreducible lowest-weight module; isomorphism of modules]\label{l13:def:irr}
An ordinary module $M$ is \emph{irreducible} if $M\neq0$ and its only subspaces stable under every mode $u_{(n)}$ are $0$
and $M$. Two bundled lowest-weight modules are \emph{isomorphic} if there is a linear isomorphism of carriers
intertwining every mode: $e(u_{(n)}x) = u_{(n)}e(x)$ for all $u\in V$, $n\in\Z$.
\end{definition}
\noindent Files \leanpath{VOA/\allowbreak Modules/\allowbreak WeakModule.lean},
\leanpath{VOA/\allowbreak Modules/\allowbreak ZhuVermaAnArchRigidity.lean},
\leanpath{VOA/\allowbreak Modules/\allowbreak Rationality.lean} and
\leanpath{VOA/\allowbreak Modules/\allowbreak ZhuBijection.lean}.
\begin{leancode}
def IsIrreducible (N : OrdVModule W M) : Prop :=
  Nontrivial M ∧ ∀ P : WeakModule.Submodule' N.toWeakModule, P.carrier = ⊥ ∨ P.carrier = ⊤
def LWModule.IsIrreducibleLW {W : VertexOperatorAlgebra 𝒱} (P : LWModule W) : Prop :=
  OrdVModule.IsIrreducible P.isMod
structure ModuleEquiv (W : VertexOperatorAlgebra 𝒱) {M N : Type u} [AddCommGroup M]
    [Module ℂ M] [AddCommGroup N] [Module ℂ N] (A : OrdVModule W M) (B : OrdVModule W N) where
  /-- The underlying linear equivalence. -/
  toLinearEquiv : M ≃ₗ[ℂ] N
  /-- Intertwining of the actions: `e(u₍ₙ₎ x) = u₍ₙ₎ e(x)`. -/
  map_act : ∀ (u : 𝒱) (n : ℤ) (x : M),
    toLinearEquiv (A.toWeakModule.act u n x) = B.toWeakModule.act u n (toLinearEquiv x)
def LWModule.Iso (P Q : LWModule W) : Prop :=
  Nonempty (ModuleEquiv W P.isMod Q.isMod)
\end{leancode}

\begin{dictionary}
\lean{WeakModule.Submodule' N.toWeakModule} is a subspace of the carrier closed under every $u_{(n)}$ (no grading
required); \lean{P.carrier = ⊥ ∨ P.carrier = ⊤} says it is $0$ or everything; \lean{Nontrivial M} says $M\neq0$. So
\lean{OrdVModule.IsIrreducible} is the textbook irreducibility of the underlying weak module. Note that the submodules are
\emph{not} required to be graded; for an ordinary module this makes no difference, because every mode-stable subspace is
stable under $L_0 = \omega_{(1)}$ and hence is graded ($L_0$ acts semisimply with the weight spaces as eigenspaces;
the library proves this in \leanpath{ZhuVermaAnArchLUGraded.lean}, \lean{submodule_isGraded}).
\lean{ModuleEquiv} bundles a $\C$-linear equivalence with the intertwining of all modes; it automatically maps
$M_\alpha$ onto $N_\alpha$ because $L_0$ is a mode. \lean{LWModule.Iso} is mere existence (\lean{Nonempty}) of such an
equivalence, a \lean{Prop}, so that isomorphism classes can be formed as a quotient. All three are exactly the textbook
notions.
\end{dictionary}

\begin{physics}
An irreducible module is a sector that cannot be split into smaller sectors: every state can be reached from every other
by acting with modes of the chiral fields. An isomorphism of sectors is a relabelling of states that commutes with all
fields, i.e.\ the two sectors are physically the same.
\end{physics}

\subsection{The faithful domain and codomain}

\begin{definition}[$\Phi(V)$ and the image-defined $\lthirteenIrr A(V)$]\label{l13:def:irrclasses}
Let $\Phi(V)$ be the set of isomorphism classes of irreducible bundled lowest-weight modules. Call an $A(V)$-module $U$
\emph{a bottom of an irreducible} if there is an irreducible lowest-weight module $P$ with $P_h\cong U$ as
$A(V)$-modules, and let $\lthirteenIrr A(V)$ denote the set of isomorphism classes of such $U$. The \emph{faithful forward
map} is
\begin{equation*}
 \mathrm{fwd}:\Phi(V)\to\lthirteenIrr A(V),\qquad [P]\mapsto[P_h].
\end{equation*}
\end{definition}
\noindent File \leanpath{VOA/\allowbreak Modules/\allowbreak ZhuVermaAnArchFaithfulBijection.lean}.
\begin{leancode}
def IrrLWModule (W : VertexOperatorAlgebra 𝒱) : Type (u + 1) :=
  {P : LWModule W // P.IsIrreducibleLW}
def Iso (P Q : IrrLWModule W) : Prop := LWModule.Iso W P.1 Q.1
def IrrLWModuleIsoClass (W : VertexOperatorAlgebra 𝒱) : Type (u + 1) :=
  Quotient (IrrLWModule.setoid W)
def IsBottomOfIrr {W : VertexOperatorAlgebra 𝒱} (U : AMod W) : Prop :=
  ∃ P : LWModule W, P.IsIrreducibleLW ∧ Nonempty (AModEquiv W P.toAMod U)
def IrrAMod (W : VertexOperatorAlgebra 𝒱) : Type (u + 1) :=
  {U : AMod W // IsBottomOfIrr U}
def Iso (U V : IrrAMod W) : Prop := Nonempty (AModEquiv W U.1 V.1)
def IrrAModIsoClass (W : VertexOperatorAlgebra 𝒱) : Type (u + 1) :=
  Quotient (IrrAMod.setoid W)

noncomputable def IrrLWModule.toIrrAMod {W : VertexOperatorAlgebra 𝒱} (P : IrrLWModule W) :
    IrrAMod W :=
  ⟨P.1.toAMod, ⟨P.1, P.2, ⟨AModEquiv.refl W P.1.toAMod⟩⟩⟩
noncomputable def forwardMapIrr (W : VertexOperatorAlgebra 𝒱) :
    IrrLWModuleIsoClass W → IrrAModIsoClass W :=
  Quotient.lift
    (fun P => (Quotient.mk (IrrAMod.setoid W) (IrrLWModule.toIrrAMod P) : IrrAModIsoClass W))
    (fun _ _ ⟨e⟩ => Quotient.sound ⟨LWModule.isoToAModEquiv W e⟩)
\end{leancode}

\begin{dictionary}
\textbf{Decoding.} \lean{{P : LWModule W // P.IsIrreducibleLW}} is a \emph{subtype}: pairs $(P,\text{proof that }P$
is irreducible$)$; \lean{P.1} is the module and \lean{P.2} the proof. \lean{IrrLWModule.setoid W} is the equivalence
relation \lean{IrrLWModule.Iso} (isomorphism of the underlying modules) and \lean{Quotient} forms the set of classes; so
\lean{IrrLWModuleIsoClass W} is $\Phi(V)$. On the other side, \lean{IsBottomOfIrr U} is literally ``there exist an
irreducible $P$ and an $A(V)$-isomorphism $P_h\cong U$'', and \lean{IrrAModIsoClass W} is the set of $A(V)$-isomorphism
classes of such $U$. \lean{Quotient.lift f h} defines a function on classes from a function \lean{f} on representatives,
given a proof \lean{h} that isomorphic inputs give equal outputs; here the proof is \lean{LWModule.isoToAModEquiv}, the
easy fact that an isomorphism of modules restricts to an $A(V)$-isomorphism of bottoms (Chapter~\ref{ch:L12}). So
\lean{forwardMapIrr W} is $[P]\mapsto[P_h]$.

\textbf{Universe bookkeeping.} The carriers of modules live in the same universe \lean{Type u} as $V$; the types of
classes live one universe up. This is invisible mathematically: ``all modules'' means all modules whose carrier is a
set of the same size class as $V$, which loses nothing (every irreducible module is a quotient of an object built from
$V$ and a vector space of the same size).

\textbf{The crucial point: the codomain is the image.} In Zhu's theorem the codomain is the set of
\emph{simple} $A(V)$-modules. Here it is, by definition, the set of $A(V)$-modules that already arise as bottoms of
irreducibles. Every such module is simple, for every $V$ (Theorem~\ref{l13:thm:simplebottom-all}); the converse, that every
simple module is such a bottom, is proved for $C_2$-cofinite $V$, and for those $V$ the two sets are identified and the
correspondence becomes a genuine bijection (Theorem~\ref{l13:thm:zhucorr}). For general $V$ the codomain remains the image;
see Section~\ref{l13:sec:scope}.
\end{dictionary}

\begin{physics}
$\Phi(V)$ is the \emph{spectrum of sectors} of the chiral theory; $\lthirteenIrr A(V)$, as defined here, is the list of
primary multiplets \emph{that actually occur in some sector}. The forward map reads off the primary multiplet of each
sector.
\end{physics}

\subsection{The headline}

\begin{theorem}[Zhu's correspondence, uniqueness form]\label{l13:thm:headline}
Let $V$ be any vertex operator algebra of the library and assume the datum \lean{Pr} of
Definition~\ref{l13:def:bottomprojection} for the regular $A(V)$-module. Then $\mathrm{fwd}:\Phi(V)\to\lthirteenIrr A(V)$
is a bijection. Equivalently: if $P,Q$ are irreducible bundled lowest-weight $V$-modules and $P_h\cong Q_{h'}$ as
$A(V)$-modules, then $P\cong Q$. The datum \lean{Pr} is supplied, for every $V$, by the theorem \lean{VOA.pbw} (through \lean{voaPBWBottomProjection}).
\end{theorem}
\noindent Files \leanpath{VOA/\allowbreak Modules/\allowbreak ZhuVermaAnArchLUBijection.lean} and
\leanpath{VOA/\allowbreak Modules/\allowbreak ZhuVermaAnArchLUBijectionPBW.lean}.
\begin{leancode}
theorem zhuBijective (Pr : BottomProjection (parabolicRepOfAMod W (regularAMod W))) :
    Function.Bijective (forwardMapIrr W) :=
theorem zhuBijective_of_voaPBW (W : VertexOperatorAlgebra 𝒱) :
    Function.Bijective (forwardMapIrr W) :=
  zhuBijective
    (voaPBWBottomProjection (parabolicRepOfAMod W (regularAMod W)))
\end{leancode}
\noindent The equivalence with uniqueness is a library theorem, recorded in the file
\leanpath{VOA/\allowbreak Modules/\allowbreak ZhuVermaAnArchZhu534.lean}; the uniqueness statement is named
\lean{IrredBottomRigidity} and lives in the file \leanpath{VOA/\allowbreak Modules/\allowbreak ZhuVermaAnArchRigidity.lean}.
\begin{leancode}
def IrredBottomRigidity (W : VertexOperatorAlgebra 𝒱) : Prop :=
  ∀ P Q : LWModule W, P.IsIrreducibleLW → Q.IsIrreducibleLW →
    Nonempty (AModEquiv W P.toAMod Q.toAMod) → LWModule.Iso W P Q
theorem forwardMapIrr_surjective_uncond (W : VertexOperatorAlgebra 𝒱) :
    Function.Surjective (forwardMapIrr W) := by
theorem forwardMapIrr_bijective_iff_rigidity :
    Function.Bijective (forwardMapIrr W) ↔ IrredBottomRigidity W :=
  ⟨fun hb => irredBottomRigidity_of_injective hb.1, forwardMapIrr_bijective_of_rigidity⟩
\end{leancode}

\begin{dictionary}
\textbf{Hypotheses.} The only explicit hypothesis of \lean{zhuBijective} is \lean{Pr}; the VOA \lean{W} is implicit and
arbitrary (any \lean{VertexOperatorAlgebra 𝒱} on any complex vector space \lean{𝒱}). There is no simplicity, rationality,
$C_2$, or finiteness hypothesis on $V$. \lean{Pr} is a structure (data, not a proposition) whose meaning is explained in
Section~\ref{l13:sec:vhat}; its type mentions \lean{regularAMod W} ($A(V)$ as a module over itself) and
\lean{parabolicRepOfAMod} (Definition~\ref{l13:def:parabolicrep}). \lean{zhuBijective_of_voaPBW} has no hypothesis at all:
it applies \lean{zhuBijective} to the term \lean{voaPBWBottomProjection ...}, which is built from \lean{VOA.pbw}.

\textbf{Conclusion.} \lean{Function.Bijective f} is Mathlib's ``injective and surjective''.

\textbf{Exactly which bijection.} Between $\Phi(V)$ = isomorphism classes of irreducible \emph{ordinary} modules whose
weights lie in one coset $h+\N$ with $M_h\neq0$, and $\lthirteenIrr A(V)$ = isomorphism classes of $A(V)$-modules
occurring as such bottoms. Classes are taken up to isomorphism on both sides.

\textbf{Same, weaker, or stronger than Zhu?} Compared with Theorem~\ref{l13:thm:zhu-lit}(3) restricted to ordinary
modules:
\begin{itemize}
\item \emph{Surjectivity} is free: \lean{forwardMapIrr_surjective_uncond} proves it with no hypothesis, by unpacking the
  definition of the codomain. It contains none of Zhu's existence theorem~(2).
\item \emph{Injectivity} is exactly Zhu's uniqueness: by \lean{forwardMapIrr_bijective_iff_rigidity}, bijectivity
  of \lean{forwardMapIrr W} is \emph{equivalent} to \lean{IrredBottomRigidity W}, which reads: for irreducible $P,Q$, an
  $A(V)$-isomorphism $P_h\cong Q_{h'}$ implies $P\cong Q$. (The two lowest weights are then equal, because $[\omega]$ acts
  on the bottom by the lowest weight.)
\item \emph{Part~(1)} of Zhu's theorem (bottoms of irreducibles are simple) is not asserted by this theorem. It is a
  separate library theorem, valid for every $V$ (Theorem~\ref{l13:thm:simplebottom-all}), so the codomain consists of
  simple modules.
\end{itemize}
So the Lean theorem is \emph{weaker} than Zhu's full theorem~(3): it is precisely its uniqueness half. It is stated for a
wider class of VOAs than many textbook versions (no CFT-type assumption is needed since every library VOA is of CFT type).

\textbf{Could it be vacuous?} The hypothesis \lean{Pr} is satisfiable (it is a true consequence of PBW,
Section~\ref{l13:sec:vhat}), and for the conclusion there is no trick: \lean{IrredBottomRigidity} is a non-trivial
statement whenever $V$ has irreducible modules (e.g.\ the adjoint module of any simple VOA). A
check in Lean against the library (not part of this release) derives \lean{IrredBottomRigidity}-shaped statements \lean{zhu_rigidity}
and \lean{zhu_rigidity_voaPBW} directly from the two headlines, and conversely \lean{bijective_of_rigidity_statement}.
\end{dictionary}

\begin{warning}
\textbf{The name oversells (moderate).} The docstrings call \lean{zhuBijective} ``Zhu Theorem 5.3.4 (faithful form)''. It
is the \emph{uniqueness} half of Zhu's correspondence only: its codomain \lean{IrrAModIsoClass} is the image of the
forward map by definition, so surjectivity is a tautology (as the library's own \lean{forwardMapIrr_surjective_uncond}
and its docstring state). A reader who sees ``bijection $\Phi(V)\leftrightarrow\lthirteenIrr A(V)$'' must not conclude
that every simple $A(V)$-module is realized; that is a separate theorem, proved only under $C_2$-cofiniteness
(Section~\ref{l13:sec:scope}). For $C_2$-cofinite $V$ the genuine bijection onto all simple $A(V)$-modules is
\lean{zhuCorrespondence_of_C2Cofinite} (Theorem~\ref{l13:thm:zhucorr}); cite that, not \lean{zhuBijective}, when announcing
``Zhu's correspondence as a bijection''. Checked in Lean against the library (the check is not part of this release):
theorems \lean{zhu_rigidity}, \lean{bijective_of_rigidity_statement},
\lean{forwardMapIrr_surjective_free}.
\end{warning}

\textbf{Proof on paper (overview).} By \lean{forwardMapIrr_bijective_iff_rigidity} it suffices to prove uniqueness. Let
$P$ be irreducible with bottom $\Omega(P) = P_h$. In Section~\ref{l13:sec:LU} we construct, from the $A(V)$-module
$\Omega(P)$ alone, the module $L(\Omega(P))$ and prove $P\cong L(\Omega(P))$. We also prove that $L(U)$ depends only on
the isomorphism class of $U$. If $\Omega(P)\cong\Omega(Q)$ then
$P\cong L(\Omega(P))\cong L(\Omega(Q))\cong Q$. \qed

\begin{status}
\lean{zhuBijective}: proved, \emph{conditional on} the explicit datum \lean{Pr}; \texttt{\#print axioms VOA.zhuBijective}
= \lean{[propext, Classical.choice, Quot.sound]}. \lean{zhuBijective_of_voaPBW}: proved with no hypothesis; \textbf{unconditional};
\texttt{\#print axioms} = \lean{[propext, Classical.choice, Quot.sound]} in a later run, after \lean{VOA.pbw} became a theorem, not re-run for this edition (at the snapshot
the list also contained \lean{pbw}, then an axiom). \lean{forwardMapIrr_surjective_uncond}, \lean{forwardMapIrr_bijective_iff_rigidity}: unconditional,
standard three axioms. Except for the re-probe just stated, these results were obtained at the snapshot (not re-run for this edition) with
\texttt{lake env lean -DmaxSynthPendingDepth=3} on a scratch file (output in Section~\ref{l13:sec:checklist}).
\end{status}

\begin{physics}
The theorem says: \emph{a sector is determined by its primary multiplet}. Two sectors whose primaries carry isomorphic
representations of the zero-mode algebra are the same sector. Physically this is the statement that all descendant
states, their norms and their correlation functions are fixed by the chiral algebra once the primaries and the
action of zero modes on them are fixed --- the conformal bootstrap principle that ``the spectrum of primaries plus the
chiral algebra determines the representation''. What the theorem does \emph{not} say is which primary multiplets are
allowed; that is the classification problem (for minimal models: the Kac table), which is the content of the existence
half.
\end{physics}

\subsection{What lies outside the headline: the scope of the codomain}\label{l13:sec:scope}

We record where the two halves of Zhu's theorem that the headline does not contain stand, and how they combine with
the headline. None of these results is in this chapter's files; we quote their statements so that the reader can see
exactly what they assume. Since the snapshot, both halves have become theorems (the first for every $V$, the second for
$C_2$-cofinite $V$), and the library now states their combination as one bijection.

\begin{definition}[Simple $A(V)$-module; realizability]\label{l13:def:simpleamod}
An $A(V)$-module $U$ is \emph{simple} if $U\neq0$ and its only $A(V)$-submodules are $0$ and $U$. The predicate
\lean{SimpleAModsRealized W} says: every simple $A(V)$-module is the bottom of an irreducible graded module with
finite-dimensional weight spaces (an \lean{IrrRealization}).
\end{definition}
\noindent Files \leanpath{VOA/\allowbreak Modules/\allowbreak ZhuVermaAnArchLeaf3.lean},
\leanpath{VOA/\allowbreak Modules/\allowbreak ZhuVermaAnArchZhu534.lean} and (Chapter~\ref{ch:L15})
\leanpath{VOA/\allowbreak Modularity/\allowbreak GNProp8Closure.lean}.
\begin{leancode}
def IsSimpleAMod {W : VertexOperatorAlgebra 𝒱} (U : AMod W) : Prop :=
  Nontrivial U.carrier ∧ ∀ S : AModSubmodule U, S.carrier = ⊥ ∨ S.carrier = ⊤
def SimpleAModsRealized (W : VertexOperatorAlgebra 𝒱) : Prop :=
  ∀ U : AMod W, IsSimpleAMod U → Nonempty (IrrRealization W U)
theorem isBottomOfIrr_iff_isSimpleAMod (hrec : AllBottomRecovers W)
    (hreal : SimpleAModsRealized W) (U : AMod W) :
    IsBottomOfIrr U ↔ IsSimpleAMod U :=
theorem simpleAModsRealized_of_C2Cofinite {W : VertexOperatorAlgebra 𝒱}
    (hC2 : C2Cofinite W)
    (Pr : BottomProjection (parabolicRepOfAMod W (regularAMod W))) :
    SimpleAModsRealized W :=
\end{leancode}

\begin{dictionary}
\lean{AModSubmodule U} is a subspace stable under every $\rho(a)$. \lean{IrrRealization W U} packages an irreducible
\lean{GradedLWWeakModule} (Definition~\ref{l13:def:gradedlw}) with finite-dimensional weight spaces and an
$A(V)$-isomorphism of its bottom with $U$; by \lean{isBottomOfIrr_of_irrRealization} it puts $U$ into the codomain
\lean{IrrAMod W}. \lean{isBottomOfIrr_iff_isSimpleAMod} identifies the codomain with the simple $A(V)$-modules, under two
hypotheses: \lean{AllBottomRecovers W} (for the direction ``bottoms are simple'', Zhu's part~(1)) and
\lean{SimpleAModsRealized W} (Zhu's part~(2)). The second is proved in Chapter~\ref{ch:L15} from $C_2$-cofiniteness and
\lean{Pr}. The first is now a theorem for every $V$ (\lean{allBottomRecovers}, Theorem~\ref{l13:thm:simplebottom-all}
below); at the snapshot it was an open hypothesis, and Chapter~\ref{ch:L12}'s Theorem~\ref{l12:thm:simplebottom} proved the
irreducible case only under semisimplicity of $A(V)$ and \lean{Pr}. A finiteness hypothesis is unavoidable for
\lean{SimpleAModsRealized} as typed: the library proves (\lean{simpleAModsRealized_finiteDimensional}) that it forces every
simple $A(V)$-module to be finite-dimensional, which fails e.g.\ for $A(V^k(\mathfrak{sl}_2))\cong U(\mathfrak{sl}_2)$.
\end{dictionary}

\begin{status}
\lean{isBottomOfIrr_iff_isSimpleAMod}: conditional on \lean{AllBottomRecovers} and \lean{SimpleAModsRealized};
\texttt{\#print axioms} = standard three (re-run at an earlier commit, not re-run for this edition). \lean{simpleAModsRealized_of_C2Cofinite}: conditional on
\lean{C2Cofinite W} and \lean{Pr} (both explicit; there is no \lean{_voaPBW} wrapper); \texttt{\#print axioms} = standard three.
Since \lean{AllBottomRecovers W} is now a theorem for every $V$ (below), for a $C_2$-cofinite $V$ both hypotheses of
\lean{isBottomOfIrr_iff_isSimpleAMod} hold, and the codomain of \lean{forwardMapIrr} is exactly the simple $A(V)$-modules.
(The earlier text of this box called ``every element of the codomain is simple'' unproved; that is no longer so.)
\end{status}

\begin{physics}
Existence (``every allowed primary multiplet generates a sector'') is the half that physicists use to \emph{classify}
sectors: solve the zero-mode algebra's representation theory, then build the sector. In rational theories the zero-mode
algebra is finite-dimensional and semisimple and the list of sectors is finite; in non-rational theories (e.g.\ the
affine current algebra at generic level, whose zero modes generate $U(\mathfrak g)$) there are infinite-dimensional
primary multiplets that do not give ordinary sectors, which is why a finiteness hypothesis appears.
\end{physics}

\begin{theorem}[Bottoms of irreducible modules are simple; Zhu 1996, Theorem~2.2.1(1)]\label{l13:thm:simplebottom-all}
For every vertex operator algebra $V$ of the library and every irreducible bundled lowest-weight module $P$, the bottom
$P_h$ is a simple $A(V)$-module. More precisely, every lowest-weight module $P$ has the \emph{bottom-recovery property}: for
every $A(V)$-submodule $S\subseteq P_h$, the $V$-submodule of $P$ generated by $S$ meets $P_h$ in exactly $S$.
\end{theorem}
\noindent The predicate is defined in \leanpath{VOA/\allowbreak Modules/\allowbreak ZhuVermaAnArchLeaf3.lean} and
\leanpath{VOA/\allowbreak Modules/\allowbreak ZhuVermaAnArchAModScope.lean}:
\begin{leancode}
def BottomRecovers : Prop :=
  ∀ S : AModSubmodule (P.toAMod),
    vSubmoduleBottomCarrier S.genVSubmodule ≤ S.carrier

def AllBottomRecovers (W : VertexOperatorAlgebra 𝒱) : Prop :=
  ∀ P : LWModule W, BottomRecovers P
\end{leancode}
\noindent and the theorems are in \leanpath{VOA/\allowbreak Huang/\allowbreak LWBottomRecovery.lean} (Chapter~\ref{ch:L21}),
namespace \lean{VOA.Huang.LWBottomRecovery}, with section variables \lean{{W : VertexOperatorAlgebra 𝒱}} and
\lean{(P : LWModule W)}:
\begin{leancode}
theorem allBottomRecovers (W : VertexOperatorAlgebra 𝒱) : AllBottomRecovers W :=
  fun P => bottomRecovers P

theorem isSimpleAMod_toAMod_of_irreducible (hP : P.IsIrreducibleLW) : IsSimpleAMod P.toAMod :=
  isSimpleAMod_of_bottomRecovers P hP (bottomRecovers P)
\end{leancode}

\begin{dictionary}
\lean{S.genVSubmodule} is the weak $V$-submodule of $P$ generated by the subspace $S$ of the bottom (all vectors
$u^1_{(n_1)}\cdots u^k_{(n_k)}s$, $s\in S$); \lean{vSubmoduleBottomCarrier} is its intersection with the bottom $P_h$. So
\lean{BottomRecovers P} says: \emph{a mode word that carries a vector of $S$ back to the bottom lands in $S$}; equivalently,
the only way to return to the ground states is through the zero-mode algebra $A(V)$, which preserves $S$. The inclusion
$\supseteq$ is trivial, so this is an equality. \lean{allBottomRecovers} proves it for every lowest-weight module of every
$V$, with no hypothesis: no $C_2$, no semisimplicity, no datum \lean{Pr}. \lean{isSimpleAMod_toAMod_of_irreducible} deduces
simplicity of the bottom of an irreducible $P$: if $0\neq S\subsetneq P_h$, the submodule generated by $S$ is nonzero, hence all
of $P$ by irreducibility, and its bottom would be $S\neq P_h$, a contradiction. \textbf{Same as Zhu's part~(1)}, for ordinary
lowest-weight modules. It removes the extra hypotheses of Chapter~\ref{ch:L12}'s Theorem~\ref{l12:thm:simplebottom}. The
project had recorded \lean{AllBottomRecovers} as unproved; that gap is now closed.
\end{dictionary}

\begin{proof}[Proof on paper]
View $P$ as an admissible module with levels $P(k)=P_{h+k}$. A vector of the generated submodule lying in $P(0)$ is a sum of
mode words $u^1_{(n_1)}\cdots u^k_{(n_k)}s$ of total degree $0$. By the commutator and associativity formulas, any such word
can be rewritten as a combination of words in which all degree-lowering modes stand to the right of all degree-raising
ones; on the bottom the lowering modes act by $0$, and what remains are products of zero modes $o(u)$, i.e.\ the action of
$A(V)$ (this is the standard argument of Dong--Li--Mason that $\Omega(M)$-generated submodules are controlled by
$A(V)$, used by Zhu for part~(1)). Since $S$ is an $A(V)$-submodule, the result lies in $S$. The Lean proof applies the
general bottom recovery of \leanpath{VOA/\allowbreak Huang/\allowbreak BottomRecovery.lean}.
\end{proof}

\begin{physics}
The ground states of a sector cannot be ``re-entered'' by a detour through excited states except through the zero modes:
whatever a product of field modes does, if it starts and ends on ground states it acts as an element of the zero-mode
algebra. Hence the ground-state multiplet of an irreducible sector is an irreducible multiplet: ``one primary per sector,
up to zero-mode symmetry''.
\end{physics}

\begin{status}
\lean{VOA.Huang.LWBottomRecovery.allBottomRecovers}, \lean{VOA.Huang.LWBottomRecovery.isSimpleAMod_toAMod_of_irreducible}:
\texttt{[propext, Classical.choice, Quot.sound]}, run at an earlier commit (not re-run for this edition) with \texttt{lake env lean -DmaxSynthPendingDepth=3}.
\textbf{Unconditional.}
\end{status}

\subsection{Zhu's correspondence as one bijection, for \texorpdfstring{$C_2$}{C2}-cofinite \texorpdfstring{$V$}{V}}\label{l13:sec:zhucorr}

With both halves available, the library states Zhu's correspondence in the form one expects: a bijection between $\Phi(V)$
and isomorphism classes of simple $A(V)$-modules, where the codomain is defined intrinsically (``simple''), not as the image of
the map.

\begin{definition}[Simple $A(V)$-modules up to isomorphism; the bottom as a simple module]\label{l13:def:simpleamodclass}
Let $\mathrm{Simple}A(V)$ be the collection of simple $A(V)$-modules (in the sense of Definition~\ref{l13:def:simpleamod}),
and $\mathrm{Simple}A(V)/{\cong}$ its set of isomorphism classes. For an irreducible lowest-weight module $P$, its bottom
$P_h$ is an element of $\mathrm{Simple}A(V)$ by Theorem~\ref{l13:thm:simplebottom-all}. The forgetful map
$\mathrm{forget}:\lthirteenIrr\mathrm{LW}(V)/{\cong}\to\Phi(V)$ sends the class of a bundled lowest-weight module to the class of
the underlying irreducible ordinary module; it is a bijection.
\end{definition}
\noindent File \leanpath{VOA/\allowbreak Modularity/\allowbreak ZhuModularityCanonical.lean}.
\begin{leancode}
def SimpleAMod (W : VertexOperatorAlgebra 𝒱) : Type (u + 1) :=
  {U : AMod W // IsSimpleAMod U}

def SimpleAMod.setoid (W : VertexOperatorAlgebra 𝒱) : Setoid (SimpleAMod W) where
  r U V := Nonempty (AModEquiv W U.1 V.1)

def SimpleAModIsoClass (W : VertexOperatorAlgebra 𝒱) : Type (u + 1) :=
  Quotient (SimpleAMod.setoid W)

noncomputable def IrrLWModule.toSimpleAMod {W : VertexOperatorAlgebra 𝒱} (P : IrrLWModule W) :
    SimpleAMod W :=
  ⟨P.1.toAMod,
    (IrrLWModule.toIrrAMod P).2.isSimpleAMod (VOA.Huang.LWBottomRecovery.allBottomRecovers W)⟩
\end{leancode}
\noindent File \leanpath{VOA/\allowbreak Modularity/\allowbreak LowestWeightExistence.lean} (Chapter~\ref{ch:L15}).
\begin{leancode}
def forgetClass (W : VertexOperatorAlgebra 𝒱) : IrrLWModuleIsoClass W → Φ W :=
  Quotient.lift (fun P => (Quotient.mk (IrrModule.setoid W) (IrrLWModule.forget P) : Φ W))
    (fun _ _ h => Quotient.sound h)

noncomputable def irrLWModuleIsoClass_equiv_Φ (W : VertexOperatorAlgebra 𝒱) :
    IrrLWModuleIsoClass W ≃ Φ W where
  toFun := forgetClass W
  invFun := lowestWeightClass W
\end{leancode}

\begin{dictionary}
\lean{{U : AMod W // IsSimpleAMod U}} is the subtype of bundled $A(V)$-modules that are simple (nonzero, with only the
trivial $A(V)$-submodules): an \emph{intrinsic} predicate on $U$, unlike \lean{IsBottomOfIrr}, which asks for an irreducible
$V$-module realising $U$. The setoid is $A(V)$-isomorphism (\lean{AModEquiv}, Chapter~\ref{ch:L12}); the displayed
\lean{where} clause shows the relation, and its three proof fields (reflexivity, symmetry, transitivity, borrowed from
\lean{AMod.setoid}) are omitted. \lean{IrrLWModule.toSimpleAMod P} is the bottom \lean{P.1.toAMod} together with the proof
that it is simple, obtained from \lean{IsBottomOfIrr.isSimpleAMod} and the theorem \lean{allBottomRecovers}; no hypothesis on
$V$. \lean{forgetClass} forgets the lowest weight $h$ (an irreducible bundled lowest-weight module \emph{is} an irreducible
ordinary module, \lean{IrrLWModule.forget}); \lean{irrLWModuleIsoClass_equiv_Φ} is the equivalence
$\lthirteenIrr\mathrm{LW}(V)/{\cong}\;\simeq\;\Phi(V)$ whose forward map is \lean{forgetClass} and whose inverse chooses the
conformal weight (existence of a lowest weight for irreducible modules; we show its first two fields, the proofs of the two
inverse laws are omitted). This identifies the domain \lean{IrrLWModuleIsoClass} of the headline with Chapter~\ref{ch:L12}'s
$\Phi(V)$.
\end{dictionary}

\begin{physics}
$\mathrm{Simple}A(V)/{\cong}$ is the list of \emph{allowed} irreducible primary multiplets, computed from the zero-mode algebra
alone, before asking whether each of them generates a sector.
\end{physics}

\begin{status}
Definitions. \lean{VOA.IrrLWModule.toSimpleAMod}: \texttt{[propext, Classical.choice, Quot.sound]} (run at an earlier commit, not re-run for this edition).
\end{status}

\begin{theorem}[Zhu's correspondence as one bijection; Zhu 1996, Theorem~2.2.2 / Gannon, Theorem~5.3.4, for $C_2$-cofinite $V$]\label{l13:thm:zhucorr}
Let $V$ be $C_2$-cofinite. There is a bijection
\begin{equation*}
  e:\Phi(V)\xrightarrow{\ \sim\ }\mathrm{Simple}A(V)/{\cong}
\end{equation*}
such that for every irreducible lowest-weight module $P$ with bottom $P_h$, $e([P])=[P_h]$.
\end{theorem}
\noindent File \leanpath{VOA/\allowbreak Modularity/\allowbreak ZhuModularityCanonical.lean}.
\begin{leancode}
theorem zhuCorrespondence_of_C2Cofinite (W : VertexOperatorAlgebra 𝒱) (hC2 : C2Cofinite W) :
    ∃ e : Φ W ≃ SimpleAModIsoClass W, ∀ P : IrrLWModule W,
      e (forgetClass W (Quotient.mk (IrrLWModule.setoid W) P)) =
        Quotient.mk (SimpleAMod.setoid W) (IrrLWModule.toSimpleAMod P) := by
\end{leancode}

\begin{dictionary}
\textbf{Decoding.} \lean{Φ W ≃ SimpleAModIsoClass W} is a Lean \emph{equivalence}: a function with a two-sided inverse,
i.e.\ a bijection. The second conjunct pins down \emph{which} bijection: for every irreducible lowest-weight module \lean{P}, the
class of \lean{P} in $\Phi(V)$ (obtained by forming its class \lean{Quotient.mk (IrrLWModule.setoid W) P} and forgetting the
lowest weight with \lean{forgetClass}) is sent to the class of its bottom. Since every element of $\Phi(V)$ is of this form
(\lean{irrLWModuleIsoClass_equiv_Φ}), $e$ is completely determined: it is $[M]\mapsto[M_h]$.

\textbf{Hypotheses.} Only \lean{hC2 : C2Cofinite W}. The datum \lean{Pr} is supplied inside the proof by
\lean{voaPBWBottomProjection} (the theorem \lean{VOA.pbw}).

\textbf{Same, weaker, or stronger than Zhu?} Compared with Theorem~\ref{l13:thm:zhu-lit}(3) restricted to ordinary modules:
\begin{itemize}
\item the codomain is the set of \emph{all} simple $A(V)$-modules up to isomorphism, defined intrinsically, so this is no longer the
  image-defined statement of the headline: surjectivity is genuine content (existence of modules);
\item injectivity is the headline's uniqueness, valid without $C_2$;
\item well-definedness into simple modules is Theorem~\ref{l13:thm:simplebottom-all}, valid without $C_2$;
\item $C_2$-cofiniteness is used \emph{only} for surjectivity, through \lean{simpleAModsRealized_of_C2Cofinite}. Zhu's existence
  theorem needs no such hypothesis, but it yields an \emph{admissible} module $L(U)$; to land in $\Phi(V)$, whose elements are
  \emph{ordinary} modules, one needs the weight spaces of $L(U)$ to be finite-dimensional, which $C_2$ guarantees. For
  $C_2$-cofinite $V$ every simple $A(V)$-module is finite-dimensional (since $A(V)$ is), as Zhu's ordinary-module version
  requires.
\end{itemize}
So, for $C_2$-cofinite $V$, this is exactly Zhu's correspondence between irreducible ordinary modules and simple
$A(V)$-modules. For $V$ that are not $C_2$-cofinite the library has only the uniqueness half (the headline).

\textbf{Could it be vacuous?} No: both sides are nonempty whenever $V$ has an irreducible module (e.g.\ the adjoint module of a
simple $V$), and the second conjunct ties $e$ to the actual modules, so it is not an abstract cardinality statement.
\end{dictionary}

\begin{proof}[Proof on paper]
By Theorem~\ref{l13:thm:simplebottom-all}, \lean{AllBottomRecovers W} holds; by \lean{simpleAModsRealized_of_C2Cofinite} (with
\lean{Pr} from \lean{VOA.pbw}), \lean{SimpleAModsRealized W} holds. So \lean{isBottomOfIrr_iff_isSimpleAMod} identifies the
image-defined codomain \lean{IrrAModIsoClass W} of the headline with \lean{SimpleAModIsoClass W}. Compose: $\Phi(V)\simeq$
\lean{IrrLWModuleIsoClass W} (inverse of \lean{irrLWModuleIsoClass_equiv_Φ}), then the headline's bijection
\lean{forwardMapIrr} (\lean{zhuBijective_of_voaPBW}), then this identification. Each step sends $[P]$ to $[P_h]$.
\end{proof}

\begin{physics}
For a $C_2$-cofinite chiral algebra, \emph{sectors are in one-to-one correspondence with irreducible representations of the
zero-mode algebra}: to classify the sectors, classify the simple modules of the finite-dimensional algebra $A(V)$. This is how
the sectors of minimal models, WZW models and lattice theories are classified in practice.
\end{physics}

\begin{status}
\lean{VOA.zhuCorrespondence_of_C2Cofinite}: \texttt{[propext, Classical.choice, Quot.sound]}, run at an earlier commit (not re-run for this edition) with
\texttt{lake env lean -DmaxSynthPendingDepth=3} on a scratch file containing the text of
\leanpath{ZhuModularityCanonical.lean} (which had no olean at that commit; its imports had). \textbf{\texttt{\#print
axioms}-clean, conditional on} \lean{C2Cofinite W} only; uses the theorem \lean{VOA.pbw}. This is the form to cite as ``Zhu's
correspondence (bijection)''.
\end{status}

%======================================================================
\section{The Frenkel--Zhu Lie algebra, induction, and the datum \texttt{Pr}}\label{l13:sec:vhat}
%======================================================================

This section explains the single input of the headline. We need a Lie algebra $\lthirteenVh$ attached to $V$, a
subalgebra $P_0$ of it, a way to make any $A(V)$-module a $P_0$-module, and induction from $P_0$ to $\lthirteenVh$.

\subsection{The Lie algebra \texorpdfstring{$\widehat V$}{V-hat}}

\begin{definition}[Frenkel--Zhu, Borcherds; Lie algebra of modes]\label{l13:def:vhat}
Let $V\otimes\C[t,t^{-1}]$ have basis-free elements $a(m) = a\otimes t^m$ ($a\in V$, $m\in\Z$), and let $D$ be the subspace
spanned by $(L_{-1}a)(m) + m\,a(m-1)$. Put $\lthirteenVh = V\otimes\C[t,t^{-1}]/D$ with bracket
\begin{equation}\label{l13:eq:bracket}
 [a(m),b(k)] = \sum_{i\ge0}\binom mi\,(a_{(i)}b)(m+k-i).
\end{equation}
The sum is finite because $a_{(i)}b = 0$ for $i\gg0$. Grade $\lthirteenVh$ by $\deg a(m) = \wt a - m - 1$ for homogeneous
$a$.
\end{definition}
\noindent Files \leanpath{VOA/\allowbreak Modules/\allowbreak ZhuVermaAnArchVhat.lean} and
\leanpath{VOA/\allowbreak Modules/\allowbreak ZhuVermaAnArchVhatJac3.lean}.
\begin{leancode}
abbrev VLoop (𝒱 : Type u) [AddCommGroup 𝒱] [Module ℂ 𝒱] : Type u := ℤ →₀ 𝒱
theorem vloopBracket_single (W : VertexOperatorAlgebra 𝒱) (m k : ℤ) (a b : 𝒱) :
    vloopBracket W (single m a) (single k b)
      = ∑ᶠ i : ℕ, (gBinom m i) • single (m + k - i) (W.smul a i b) := by
noncomputable def loopDerivSubmodule (W : VertexOperatorAlgebra 𝒱) : Submodule ℂ (VLoop 𝒱) :=
  Submodule.span ℂ {x | ∃ (a : 𝒱) (m : ℤ),
    x = single m (W.L (-1) a) + (m : ℂ) • single (m - 1) a}
abbrev VHatCarrier (W : VertexOperatorAlgebra 𝒱) : Type u := VLoop 𝒱 ⧸ loopDerivSubmodule W
noncomputable def gen (W : VertexOperatorAlgebra 𝒱) (a : 𝒱) (m : ℤ) : VHatCarrier W :=
  mk W (single m a)
noncomputable def deg (W : VertexOperatorAlgebra 𝒱) (d : ℤ) : Submodule ℂ (VHatCarrier W) :=
  Submodule.span ℂ {x | ∃ (w m : ℤ) (a : 𝒱), a ∈ W.grading w ∧ w - m - 1 = d ∧ x = gen W a m}
noncomputable def lieRing (W : VertexOperatorAlgebra 𝒱) : LieRing (VHatCarrier W) :=
  lieRingOfData W (bracketDescends W)
    (vhatAntisymm W (bracketDescends W)) (vhatLeibniz_holds W)
\end{leancode}

\begin{dictionary}
\textbf{Decoding.} \lean{ℤ →₀ 𝒱} is Mathlib's type of finitely supported functions $\Z\to V$, i.e.\
$V\otimes\C[t,t^{-1}]$; \lean{single m a} is $a\otimes t^m = a(m)$. \lean{vloopBracket W} is the bilinear extension of
\eqref{l13:eq:bracket}; \lean{vloopBracket_single} states its value on generators, with \lean{gBinom m i}
$=\binom mi$ for $m\in\Z$, $i\in\N$ and \lean{W.smul a i b} $= a_{(i)}b$; \lean{∑ᶠ} is a finitely supported sum, and its
support is finite by the truncation axiom (the library proves this, \lean{brLevel_support_finite'}). \lean{W.L (-1)} is
$L_{-1}$. \lean{loopDerivSubmodule W} is $D$, \lean{VHatCarrier W} is $\lthirteenVh$, \lean{gen W a m} is the class of
$a(m)$, and \lean{deg W d} is the span of the homogeneous generators of degree $d$.

\textbf{The Lie structure is a theorem.} \lean{lieRing W} is assembled from three proved facts: the bracket descends to the
quotient by $D$ (\lean{bracketDescends}: $D$ is an ideal, because $L_{-1}$ is a derivation of every product and
$(L_{-1}a)_{(i)} = -i\,a_{(i-1)}$), antisymmetry (\lean{vhatAntisymm}, from skew-symmetry
$Y(a,z)b = e^{zL_{-1}}Y(b,-z)a$, valid only modulo $D$), and the Jacobi identity (\lean{vhatLeibniz_holds}, from the
Borcherds identity). There is no central extension: the conformal data already live in $V$. This is the
Lie algebra of Frenkel--Zhu, Borcherds and Dong--Li--Mason. It is attached locally as an instance (\lean{instLieRingVHat})
in the files that need it.

\textbf{Same as the literature?} Yes. In the literature the quotient is by the image of
$L_{-1}\otimes1+1\otimes\frac{d}{dt}$; on $a\otimes t^m$ this operator gives $(L_{-1}a)\otimes t^m + m\,a\otimes t^{m-1}$,
exactly the generators of $D$.
\end{dictionary}

\textbf{Proof on paper (that $\lthirteenVh$ is a Lie algebra).} The commutator formula of a vertex algebra says that on any
module $[a_{(m)},b_{(k)}] = \sum_i\binom mi(a_{(i)}b)_{(m+k-i)}$, and the translation axiom says
$(L_{-1}a)_{(m)} = -m\,a_{(m-1)}$. So \eqref{l13:eq:bracket} is the bracket of mode operators, and $D$ lies in the kernel of
$a(m)\mapsto a_{(m)}$ on every module. To prove the Lie axioms \emph{abstractly} (without a faithful module) one checks
them on generators: (i) $D$ is an ideal because $L_{-1}(a_{(i)}b) = (L_{-1}a)_{(i)}b + a_{(i)}L_{-1}b$ and
$(L_{-1}a)_{(i)}b = -i\,a_{(i-1)}b$, which turns $[D,\cdot]$ into combinations of generators of $D$; (ii) antisymmetry
follows from skew-symmetry $b_{(k)}a = \sum_j(-1)^{k+j+1}\frac{1}{j!}L_{-1}^j(a_{(k+j)}b)$, the $L_{-1}^j$ terms being
absorbed by $D$; (iii) the Jacobi identity on three generators is the Borcherds identity with $\ell = 0$, after a
Vandermonde-type rearrangement of binomial coefficients. \qed

\begin{status}
\lean{VHat.lieRing}, \lean{VHat.bracketDescends}, \lean{VHat.vhatLeibniz_holds}: proved, unconditional;
\texttt{\#print axioms} = \lean{[propext, Classical.choice, Quot.sound]}.
\end{status}

\begin{physics}
$\lthirteenVh$ is the Lie algebra of \emph{all Laurent modes of all chiral fields}, with the commutators dictated by the
singular part of the OPE --- the ``mode algebra'' of the chiral algebra. For the free boson it contains the Heisenberg
algebra $[J_m,J_k] = m\,\delta_{m+k,0}$ (from $J(z)J(w)\sim1/(z-w)^2$), for an affine current algebra it contains the
affine Kac--Moody algebra, and for the stress tensor the Virasoro algebra. Every sector of the chiral theory is a
representation of $\lthirteenVh$; the converse is false, because the OPE also contains \emph{regular} terms (normal-ordered
products) that a mere Lie-algebra representation need not respect. That defect is repaired below by a quotient.
\end{physics}

\subsection{The parabolic subalgebra and the \texorpdfstring{$A(V)$}{A(V)}-modules}

\begin{definition}[$P_n$; parabolic representation]\label{l13:def:parabolicrep}
For $n\in\N$ let $P_n\subseteq\lthirteenVh$ be the span of the homogeneous generators of degree $0$ or $\le-n-1$; its
\emph{deep part} is the span of those of degree $\le-n-1$. A \emph{parabolic representation} of level $n$ on a vector
space $U$ is a Lie-algebra homomorphism $P_n\to\End U$ which kills the deep part and sends $\lthirteenvac(-1)$ to the
identity.
\end{definition}
\noindent Files \leanpath{VOA/\allowbreak Modules/\allowbreak ZhuVermaAnArchDLM3Deg.lean} and
\leanpath{VOA/\allowbreak Modules/\allowbreak ZhuVermaAnArchDLM3Induced.lean}.
\begin{leancode}
noncomputable def degSpan (W : VertexOperatorAlgebra 𝒱) (S : Set ℤ) :
    Submodule ℂ (VHatCarrier W) :=
  Submodule.span ℂ {x | ∃ (w m : ℤ) (a : 𝒱), a ∈ W.grading w ∧ w - m - 1 ∈ S ∧ x = gen W a m}
def deepSet (n : ℕ) : Set ℤ := {d : ℤ | d ≤ -(n : ℤ) - 1}
def parabolicSet (n : ℕ) : Set ℤ := {d : ℤ | d = 0 ∨ d ≤ -(n : ℤ) - 1}
noncomputable def deepPart (W : VertexOperatorAlgebra 𝒱) (n : ℕ) : Submodule ℂ (VHatCarrier W) :=
  degSpan W (deepSet n)
noncomputable def parabolic (W : VertexOperatorAlgebra 𝒱) (n : ℕ) : Submodule ℂ (VHatCarrier W) :=
  degSpan W (parabolicSet n)
noncomputable def Pn (W : VertexOperatorAlgebra 𝒱) (n : ℕ) : LieSubalgebra ℂ (VHatCarrier W) :=
  { parabolic W n with
    lie_mem' := fun hx hy => degSpan_lie_mem W (parabolicSet_add n) hx hy }
structure ParabolicRep (W : VertexOperatorAlgebra 𝒱) (n : ℕ) (U : Type u)
    [AddCommGroup U] [Module ℂ U] where
  /-- The `P_n`-action on `U`. -/
  act : Pn W n →ₗ⁅ℂ⁆ End ℂ U
  /-- The deep part `⨁_{p>n} V̂(−p)` acts trivially (DLM3 (4.1)). -/
  deep_zero : ∀ g : Pn W n, (g : VHatCarrier W) ∈ deepPart W n → act g = 0
  /-- The vacuum loop mode `𝟙(−1)` acts as the identity (the unit of `A_n(𝒱)`). -/
  vacuum_one : act ⟨gen W W.vacuum (-1), vacuumGen_mem_parabolic W n⟩ = 1
\end{leancode}

\begin{dictionary}
\lean{LieSubalgebra ℂ (VHatCarrier W)} is a subspace closed under the bracket; closure is \lean{degSpan_lie_mem} with the
observation that the degree set $\{0\}\cup(-\infty,-n-1]$ is closed under addition. \lean{Pn W n →ₗ⁅ℂ⁆ End ℂ U} is a
Lie-algebra homomorphism into the commutator Lie algebra of $\End U$. The two fields are the two normalizations of
Dong--Li--Mason's (4.1). Note the sign convention: \lean{deg} is $\wt a - m - 1$, so ``degree $\le-n-1$'' means the mode
$a_{(m)}$ \emph{lowers} weight by at least $n+1$; on level $n$ of a module these modes act by zero.
\end{dictionary}

\begin{physics}
$P_0$ is spanned by the zero modes (degree $0$) and all modes that lower the energy. On the ground states of a sector the
lowering modes vanish and the zero modes act through $A(V)$: this is exactly a parabolic representation of level $0$. The
choice of $P_0$ is the chiral-algebra analogue of choosing a Borel-plus-Cartan subalgebra when constructing a Verma module:
``annihilation operators kill the vacuum''.
\end{physics}

\begin{definition}[The regular $A(V)$-module; the parabolic representation of an $A(V)$-module]\label{l13:def:regular}
The \emph{regular module} is $A(V)$ acting on itself by left multiplication. Every $A(V)$-module $U$ becomes a
level-$0$ parabolic representation by $a(m)\mapsto\rho([\pi_{m+1}a])$, where $\pi_w$ is the projection onto $V_w$: a
degree-$0$ generator $a(\wt a-1)$ acts as $[a]$, every other homogeneous generator by $0$.
\end{definition}
\noindent File \leanpath{VOA/\allowbreak Modules/\allowbreak ZhuVermaAnArchDLM3Zero.lean}.
\begin{leancode}
noncomputable def regularAMod (W : VertexOperatorAlgebra 𝒱) : AMod W where
  carrier := A W
  ρ := Amul (zhuRingData_holds W)
  ...
noncomputable def zhuLoop (W : VertexOperatorAlgebra 𝒱) : VLoop 𝒱 →ₗ[ℂ] A W :=
  Finsupp.lsum ℂ (fun m : ℤ => (_root_.VOA.mk W).comp (gradeComp W (m + 1)))
noncomputable def zhuDeg0 (W : VertexOperatorAlgebra 𝒱) : VHatCarrier W →ₗ[ℂ] A W :=
  Submodule.liftQ (loopDerivSubmodule W) (zhuLoop W) (by
  ...
noncomputable def parabolicRepOfAMod (W : VertexOperatorAlgebra 𝒱) (U : AMod W) :
    ParabolicRep W 0 U.carrier where
  act :=
    { toFun := fun g => U.ρ (zhuDeg0 W (g : VHatCarrier W))
  ...
\end{leancode}

\begin{dictionary}
\lean{Amul (zhuRingData_holds W)} is left multiplication in the ring $A(V)$ (Chapter~\ref{ch:L12},
Theorem~\ref{l12:thm:zhuring}); \lean{regularAMod W} is the bundled $A(V)$-module $A(V)$. It is infinite-dimensional in
general (for the Heisenberg VOA $A(V)\cong\C[x]$); this matters because an \lean{OrdVModule} cannot have it as bottom,
which is why the library works with weak modules in the construction. \lean{zhuLoop W} is $a(m)\mapsto[\pi_{m+1}a]$ on $V\otimes\C[t,t^{-1}]$, where \lean{gradeComp W w} is the projection $\pi_w$ onto $V_w$ along the grading, and \lean{zhuDeg0 W} is the induced linear map
$\lthirteenVh\to A(V)$, $a(m)\mapsto[\pi_{m+1}a]$; it is well defined on the quotient by $D$ because
$(L_{-1}+L_0)w\in O(V)$ (\lean{zhuLgen_mem_O}). \lean{parabolicRepOfAMod W U} is the parabolic representation
$g\mapsto\rho(\text{zhuDeg0}(g))$. Its three obligations are discharged by proofs: the Lie-homomorphism property
(\lean{zhuDeg0_bracket_parabolic}, which on the degree-$0$ part is Zhu's commutator formula
$u*v-v*u\equiv\Res_zY(u,z)v(1+z)^{\wt u-1}\pmod{O(V)}$, \lean{zhu_comm_mod_O}), the vanishing on the deep part, and
$\rho([\lthirteenvac]) = 1$. This is Dong--Li--Mason's Proposition~3.1 at $n=0$: $a(\wt a-1)\mapsto[a]$ is a surjective
Lie homomorphism from the degree-$0$ part of $\lthirteenVh$ to $A(V)$ with its commutator bracket.
\end{dictionary}

\begin{status}
\lean{regularAMod}, \lean{zhuDeg0}, \lean{parabolicRepOfAMod}: definitions whose obligations are proved, unconditional;
\texttt{\#print axioms} = standard three.
\end{status}

\begin{physics}
The zero mode $o(a)$ of a field of weight $\wt a$ is its mode $a(\wt a-1)$; ``$a(\wt a-1)\mapsto[a]$'' is the
statement that on ground states the zero modes multiply according to Zhu's product. The regular module is the
``universal multiplet'': any primary multiplet is a quotient of copies of it.
\end{physics}

\subsection{The induced module and the bottom map}

\begin{definition}[Induced module, bottom map, levels]\label{l13:def:indmod}
For a parabolic representation $P$ of level $n$ on $U$, let $\lthirteenUE(\lthirteenVh)$ be the universal enveloping algebra
and
\begin{equation*}
 M_n(U) = \lthirteenUE(\lthirteenVh)\otimes_{\lthirteenUE(P_n)}U
 = \bigl(\lthirteenUE(\lthirteenVh)\otimes_\C U\bigr)\big/\operatorname{span}\{y\,g\otimes u - y\otimes g\cdot u :
 y\in\lthirteenUE(\lthirteenVh),\ g\in P_n,\ u\in U\}.
\end{equation*}
The \emph{bottom map} is $U\to M_n(U)$, $u\mapsto[1\otimes u]$. The \emph{level} $k$ part $M_n(U)(k)$ is spanned by
$[y\otimes u]$ with $y$ of degree $k-n$ (degree extended multiplicatively from $\lthirteenVh$ to $\lthirteenUE(\lthirteenVh)$).
\end{definition}
\noindent Files \leanpath{VOA/\allowbreak Modules/\allowbreak ZhuVermaAnArchMUVhat.lean},
\leanpath{VOA/\allowbreak Modules/\allowbreak ZhuVermaAnArchDLM3Induced.lean},
\leanpath{VOA/\allowbreak Modules/\allowbreak ZhuVermaAnArchDLM3Admissible.lean} and
\leanpath{VOA/\allowbreak Modules/\allowbreak ZhuVermaAnArchDLM3Level.lean}.
\begin{leancode}
abbrev UEnv (W : VertexOperatorAlgebra 𝒱) : Type u :=
  UniversalEnvelopingAlgebra ℂ (VHatCarrier W)
abbrev IndCover (W : VertexOperatorAlgebra 𝒱) (U : Type u) [AddCommGroup U] [Module ℂ U] :
    Type u := UEnv W ⊗[ℂ] U
def balancingSet (P : ParabolicRep W n U) : Set (IndCover W U) :=
  {t | ∃ (y : UEnv W) (g : Pn W n) (u : U),
    t = (y * uι W (g : VHatCarrier W)) ⊗ₜ[ℂ] u - y ⊗ₜ[ℂ] (P.act g u)}
noncomputable def balancing (P : ParabolicRep W n U) : Submodule ℂ (IndCover W U) :=
  Submodule.span ℂ (balancingSet P)
abbrev IndMod (P : ParabolicRep W n U) : Type u := IndCover W U ⧸ balancing P
noncomputable def indBottom (P : ParabolicRep W n U) : U →ₗ[ℂ] IndMod P :=
  (indMk P).comp (TensorProduct.mk ℂ (UEnv W) U 1)
noncomputable def levCover (P : ParabolicRep W n U) (k : ℤ) : Submodule ℂ (IndMod P) :=
  Submodule.span ℂ {z | ∃ y ∈ uenvDeg W (k - (n : ℤ)), ∃ u : U, z = indMk P (y ⊗ₜ[ℂ] u)}
\end{leancode}

\begin{dictionary}
\lean{UniversalEnvelopingAlgebra ℂ (VHatCarrier W)} is Mathlib's $\lthirteenUE(\lthirteenVh)$, and \lean{uι W} its
canonical Lie map $\lthirteenVh\to\lthirteenUE(\lthirteenVh)$. \lean{⊗[ℂ]} and \lean{⊗ₜ[ℂ]} are Mathlib's tensor product of
vector spaces and its pure tensors. The balancing relations use only generators $g\in P_n$, not all of
$\lthirteenUE(P_n)$; since $\lthirteenUE(P_n)$ is generated by $P_n$, the span is the same, so \lean{IndMod P} is
$\lthirteenUE(\lthirteenVh)\otimes_{\lthirteenUE(P_n)}U$. \lean{indMk P} is the quotient map, \lean{indBottom P} is
$u\mapsto[1\otimes u]$, and \lean{uenvDeg W d} is the span of products of homogeneous generators of total degree $d$.
The $\lthirteenVh$-action on \lean{IndMod P} is left multiplication (\lean{indRho}); it is a representation of the Lie algebra
for free.

\textbf{Two facts proved without any PBW input.} The levels span (\lean{levCover_spanning}); $M_0(U)(d) = 0$ for $d<0$
and $M_0(U)(0)$ is the image of the bottom map (\lean{levCover_eq_bot_of_neg}, \lean{levCover_zero_eq_range}); and the
levels are linearly independent:
\end{dictionary}

\begin{leancode}
theorem levCover_iSupIndep (P : ParabolicRep W n U) : iSupIndep (levCover P) := by
\end{leancode}
\noindent (file \leanpath{VOA/\allowbreak Modules/\allowbreak ZhuVermaAnArchDLM3PBWIndependence.lean}).
\textbf{Proof on paper.} Let $\epsilon = \omega(1)\in\lthirteenVh$ (the class of $\omega\otimes t$, i.e.\ the mode $L_0$; the
library's \lean{degreeElement}). By \eqref{l13:eq:bracket}, $[\epsilon,a(m)] = (\wt a-m-1)\,a(m)$ for homogeneous $a$: the
adjoint action of $\epsilon$ is multiplication by the degree. Extend $\operatorname{ad}\epsilon$ to a derivation $\partial$ of
$\lthirteenUE(\lthirteenVh)$; then $\partial y = d\,y$ for $y$ of degree $d$. Define $E$ on $\lthirteenUE(\lthirteenVh)\otimes U$ by
$E(y\otimes u) = \partial y\otimes u$. It preserves the balancing relations: $\partial(yg)\otimes u - \partial y\otimes gu =
(\partial y\,g\otimes u - \partial y\otimes gu) + y[\epsilon,g]\otimes u$, and $[\epsilon,g]$ is a multiple of $g$ that is $0$ if
$\deg g = 0$ and deep if $g$ is deep, so it acts on $U$ by zero and $y[\epsilon,g]\otimes u$ is itself a balancing relation. Hence
$E$ descends to $M_n(U)$, and level $k$ lies in the eigenspace of eigenvalue $k-n$. Eigenspaces for distinct eigenvalues are
independent. \qed

\begin{physics}
$M_0(U)$ is the space of all formal descendants $a^1_{(m_1)}\cdots a^k_{(m_k)}|u\rangle$ of a primary multiplet, modulo
only the commutation relations of the mode algebra and ``lowering modes annihilate primaries''. The level is the energy
above the primaries. Independence of levels is the statement that energy is a good quantum number: it is measured by an
operator.
\end{physics}

\subsection{The datum \texttt{BottomProjection} and its equivalents}

\begin{definition}[Bottom projection; DLM's PBW step]\label{l13:def:bottomprojection}
A \emph{bottom projection} for a level-$n$ parabolic representation $P$ on $U$ is a linear map $\pi:M_n(U)\to U$ with
$\pi([1\otimes u]) = u$ for all $u$, and $\pi = 0$ on every level $M_n(U)(k)$, $k\neq n$. The \emph{PBW freeness} of $P$ is
the conjunction of: the bottom map is injective, and the levels are independent.
\end{definition}
\noindent Files \leanpath{VOA/\allowbreak Modules/\allowbreak ZhuVermaAnArchDLM3Survival.lean} and
\leanpath{VOA/\allowbreak Modules/\allowbreak ZhuVermaAnArchDLM3PBWSplit.lean}.
\begin{leancode}
structure BottomProjection (P : ParabolicRep W n U) where
  /-- The projection onto the bottom `U` of `M_n(U)`. -/
  proj : IndMod P →ₗ[ℂ] U
  /-- It retracts the bottom map `u ↦ [1 ⊗ u]`. -/
  proj_bottom : ∀ u : U, proj (indBottom P u) = u
  /-- It annihilates every level other than `n` (DLM3: "`U*` annihilates `⨁_{i≠n} M(U)(i)`"). -/
  proj_levCover : ∀ ⦃k : ℤ⦄, k ≠ (n : ℤ) → ∀ ⦃z : IndMod P⦄, z ∈ levCover P k → proj z = 0
structure PBWFreeness (P : ParabolicRep W n U) : Prop where
  /-- **PBW, part 1** — the bottom `U` embeds: `u ↦ [1 ⊗ u]` is injective. -/
  bottom_injective : Function.Injective (indBottom P)
  /-- **Historical PBW part 2, now a theorem** — the levels `M_n(U)(k) = U(V̂)_{k−n}U` are
  independent. `levCover_iSupIndep` supplies this field unconditionally; together with
  `levCover_spanning` it gives the internal `ℤ`-grading of `M_n(U)`. -/
  levels_iSupIndep : iSupIndep (levCover P)
\end{leancode}
\noindent The relations between them (files \leanpath{VOA/\allowbreak Modules/\allowbreak ZhuVermaAnArchDLM3PBWIndependence.lean}
and \leanpath{VOA/\allowbreak Modules/\allowbreak ZhuVermaAnArchDLM3BottomProjectionBoundary.lean}):
\begin{leancode}
theorem pbwFreeness_iff_bottom_injective (P : ParabolicRep W n U) :
    PBWFreeness P ↔ Function.Injective (indBottom P) :=
noncomputable def bottomProjectionOfInjective (Q : ParabolicRep W 0 U)
    (hinj : Function.Injective (indBottom Q)) : BottomProjection Q :=
theorem nonempty_bottomProjection_iff_bottom_injective (Q : ParabolicRep W 0 U) :
    Nonempty (BottomProjection Q) ↔ Function.Injective (indBottom Q) := by
\end{leancode}

\begin{dictionary}
\textbf{\lean{BottomProjection}} is a structure carrying \emph{data} (the map \lean{proj}) and two equations. \lean{⦃k⦄} are
strict-implicit binders (inferred); \lean{levCover P k} is $M_n(U)(k)$. Its use in the proof is Dong--Li--Mason's: to
extend a linear functional $u'$ on $U$ to $M_n(U)$ by zero off level $n$, which defines the radical $J$ below.

\textbf{What it really asserts.} By \lean{nonempty_bottomProjection_iff_bottom_injective}, at level $0$ a bottom projection
exists \emph{iff} the bottom map $u\mapsto[1\otimes u]$ is injective. (Given injectivity, the projection is constructed
from the independence of levels, which is a theorem.) Likewise \lean{PBWFreeness P} is equivalent to injectivity.
So the datum \lean{Pr} in the headline is precisely:
\begin{quote}
\emph{in $\lthirteenUE(\lthirteenVh)\otimes_{\lthirteenUE(P_0)}A(V)$, the canonical copy $1\otimes A(V)$ of the regular
$A(V)$-module is not collapsed.}
\end{quote}
Note \lean{Pr} is required \emph{only for the regular module}, not for every $U$; the proof transports what it needs to
arbitrary $U$ through the Dong--Jiang construction (Section~\ref{l13:sec:MU}).

\textbf{Is it a genuine open input?} No longer. \lean{voaPBWBottomProjection} produces
\lean{BottomProjection (parabolicRepOfAMod W (regularAMod W))} for every $V$ from the theorem \lean{VOA.pbw} (formerly a project
axiom, now a theorem). At the snapshot it was open in Lean (supplied only by the then-axiom). The library also proves a
bottom projection for a \emph{different} $A(V)$-module, the weight-zero space of the adjoint module
(\lean{adjointBottomProjection}, below), which shows the structure is satisfiable but does not discharge \lean{Pr}.
Mathematically it is \emph{true for every} $V$: it is a special case of the Poincar\'e--Birkhoff--Witt theorem (next
theorem), and the library supplies it from \lean{VOA.pbw}.

\textbf{How could it be wrong?} Two traps were checked. (a) A bottom projection could be vacuous if the levels were not
independent (then ``$\pi = 0$ on other levels'' might force $\pi = 0$); independence is a theorem. (b) Requiring it for
\emph{all} $U$ would be harmless too (PBW is true for all $U$), but the headline asks for less.
\end{dictionary}

\begin{status}
\lean{BottomProjection}, \lean{PBWFreeness}: definitions. \lean{levCover_iSupIndep}, \lean{pbwFreeness_iff_bottom_injective},
\lean{bottomProjectionOfInjective}, \lean{nonempty_bottomProjection_iff_bottom_injective}: proved, unconditional,
standard three axioms.
\end{status}

\begin{physics}
The datum says that in the space of formal descendants of the universal primary multiplet, the primaries themselves are
not accidentally identified with zero by the commutation relations. Physicists take this for granted (``the Verma module
has a basis of ordered monomials''); mathematically it is the PBW theorem.
\end{physics}

\subsection{How \texorpdfstring{\lean{VOA.pbw}}{VOA.pbw} supplies the datum}

\begin{theorem}[Bottom injectivity from PBW]\label{l13:thm:voapbw}
For every Lie algebra $\mathfrak g$ over $\C$, every Lie subalgebra $\mathfrak p$ and every $\mathfrak p$-module $U$, the map
$U\to\lthirteenUE(\mathfrak g)\otimes_{\lthirteenUE(\mathfrak p)}U$, $u\mapsto1\otimes u$, is injective. Consequently every
parabolic representation of $\lthirteenVh$ has injective bottom map, is PBW-free, and (at level $0$) has a bottom projection.
\end{theorem}
\noindent Files \leanpath{VOA/\allowbreak Foundations/\allowbreak AffineLie.lean} (the theorem \lean{VOA.pbw}; Chapter~\ref{ch:L02},
Section~\ref{l02:sec:axiom}) and \leanpath{VOA/\allowbreak Modules/\allowbreak ZhuVermaAnArchDLM3PBWAxiom.lean}.
\begin{leancode}
structure PBWPrinciple : Prop where
  /-- PBW nonvanishing for every affine vacuum module. -/
  affineVacuum : ∀ {𝔤 : Type u} [LieRing 𝔤] [LieAlgebra ℂ 𝔤]
      (κ : LinearMap.BilinForm ℂ 𝔤) (hκ : κ.IsInvariant) (k : ℂ),
    PBW κ hκ k
  /-- The canonical bottom embeds in every module induced from a Lie subalgebra. -/
  inducedBottom : ∀ {g U : Type u} [LieRing g] [LieAlgebra ℂ g]
      [AddCommGroup U] [Module ℂ U] (p : LieSubalgebra ℂ g)
      (rho : p →ₗ⁅ℂ⁆ Module.End ℂ U),
    Function.Injective (PBWInduction.bottom p rho)
theorem pbw : PBWPrinciple :=
  ⟨fun κ hκ k => AffineLie.vacuumVector_ne_zero_of_induction κ hκ k,
    fun p rho => PBWTheorem.bottom_injective p rho⟩
theorem voaPBWBottomInjective (P : ParabolicRep W n U) :
    Function.Injective (indBottom P) :=
  VOA.PBWPrinciple.inducedBottom
    (VOA.pbw : VOA.PBWPrinciple.{u}) (Pn W n) P.act
theorem voaPBWFreeness (P : ParabolicRep W n U) : PBWFreeness P :=
  pbwFreeness_of_bottom_injective P (voaPBWBottomInjective P)
noncomputable def voaPBWBottomProjection (Q : ParabolicRep W 0 U) : BottomProjection Q :=
  bottomProjectionOfInjective Q (voaPBWBottomInjective Q)
\end{leancode}

\begin{dictionary}
\textbf{The PBW principle.} \lean{PBWPrinciple} has two fields (Chapter~\ref{ch:L02} identifies both with classical PBW
consequences and presents their Lean proof). Only \lean{inducedBottom} is used here: for any Lie algebra \lean{g}, Lie subalgebra \lean{p} and Lie
representation \lean{rho : p →ₗ⁅ℂ⁆ Module.End ℂ U}, the bottom map into
\lean{PBWInduction.Induced p rho} $=\lthirteenUE(\mathfrak g)\otimes_\C U/\operatorname{span}\{y\,x\otimes u - y\otimes\rho(x)u\}$
is injective. The library proves that, for \lean{g = VHatCarrier W}, \lean{p = Pn W n}, \lean{rho = P.act}, this
generic quotient and bottom map are \emph{definitionally} the ones of Definition~\ref{l13:def:indmod}
(\lean{pbwInduction_balancing_eq}, \lean{pbwInduction_bottom_eq}, both \lean{rfl}). So \lean{voaPBWBottomInjective} is a
literal specialization; no translation is hidden.

\textbf{Universe.} \lean{VOA.PBWPrinciple.{u}} instantiates the theorem at the universe of $V$; the statement quantifies over Lie
algebras in every universe, which is harmless.
\end{dictionary}

\textbf{Proof on paper (of the mathematics behind \lean{VOA.pbw}).} By PBW, if $\{x_i\}_{i\in I}$ is an ordered basis of
$\mathfrak g$ extending a basis of $\mathfrak p$ (with the elements of the complement placed first), the ordered monomials form a
basis of $\lthirteenUE(\mathfrak g)$, and $\lthirteenUE(\mathfrak g)$ is a free right $\lthirteenUE(\mathfrak p)$-module on the
ordered monomials in the complement. Hence $\lthirteenUE(\mathfrak g)\otimes_{\lthirteenUE(\mathfrak p)}U\cong
\bigoplus_{\text{monomials }m}m\otimes U$, and the summand of the empty monomial is $1\otimes U\cong U$. For
$\mathfrak g = \lthirteenVh$, which is infinite-dimensional, one needs PBW for arbitrary Lie algebras over a field, which is
classical (the symmetrization map is an isomorphism of coalgebras in characteristic $0$; or Bergman's diamond lemma). Mathlib
does not contain PBW; the library proves the needed relative form itself (\lean{PBWTheorem.bottom_injective},
Chapter~\ref{ch:L02}, Theorem~\ref{l02:thm:bottominj}), by Humphreys' relative construction rather than via the full PBW
basis. \qed

\begin{status}
\lean{voaPBWBottomInjective}, \lean{voaPBWFreeness}, \lean{voaPBWBottomProjection}: \texttt{\#print axioms} =
\lean{[propext, Classical.choice, Quot.sound]} (re-probed at a later commit, after \lean{VOA.pbw} became a theorem, not re-run for this edition; at the snapshot the list also contained
\lean{pbw}). Unconditional; they use the theorem \lean{VOA.pbw} (field \lean{inducedBottom}). Every other declaration of
this chapter is free of \lean{pbw}.
\end{status}

\begin{physics}
PBW is the algebraic statement that a Fock-space-like basis exists: ordering creation operators to the left and
annihilation operators to the right, every state of the induced module is uniquely a combination of ordered products of
creation operators acting on the primaries. It is used tacitly whenever a physicist counts descendants (e.g.\ the character
$q^h/\prod(1-q^n)$ of a Virasoro Verma module).
\end{physics}

\begin{theorem}[A satisfiability check: the adjoint bottom]\label{l13:thm:adjoint}
For the $A(V)$-module $V_0$ (the weight-zero space of the adjoint module, $=\C\lthirteenvac$), the bottom map is injective,
without PBW, and hence a bottom projection exists.
\end{theorem}
\noindent File \leanpath{VOA/\allowbreak Modules/\allowbreak ZhuVermaAnArchDLM3AdjointBottom.lean}.
\begin{leancode}
noncomputable abbrev adjointBottomAMod (W : VertexOperatorAlgebra 𝒱) : AMod W :=
  lowestAMod (adjointModule_isLowestWeight W)
noncomputable def adjointBottomProjection (W : VertexOperatorAlgebra 𝒱) :
    BottomProjection (parabolicRepOfAMod W (adjointBottomAMod W)) :=
  bottomProjectionOfInjective _ (adjoint_indBottom_injective W)
\end{leancode}

\begin{dictionary}
The proof uses the general principle \lean{indBottom_injective_of_equivariant_embedding}: if $U$ embeds $P_0$-equivariantly
into \emph{some} $\lthirteenVh$-representation $M$, then the universal property of induction gives a left inverse of the
bottom map. For $U = V_0\subseteq V$ one takes $M = V$ with the action $a(m)\mapsto a_{(m)}$. This instance is \emph{not}
the regular module, as the file states; it is evidence that \lean{BottomProjection} is satisfiable, nothing more.
\end{dictionary}

\textbf{Proof on paper.} Given an equivariant injective $\varphi:U\to M$, the universal property of $M_0(U)$ gives a
$\lthirteenVh$-map $\Phi:M_0(U)\to M$ with $\Phi([1\otimes u]) = \varphi(u)$; if $[1\otimes u] = 0$ then $\varphi(u) = 0$
and $u = 0$. \qed

\begin{status}
\lean{adjointBottomProjection}: proved, unconditional, standard three axioms, no \lean{pbw}.
\end{status}

\begin{physics}
The vacuum is a primary of every chiral algebra, and the vacuum sector exists; therefore the formal descendants of the vacuum
cannot collapse the vacuum itself. The same argument would prove \lean{Pr} if one already had a sector whose primaries
contain a faithful copy of the regular module --- but constructing such a sector is exactly what Zhu's theorem is for, so
the argument would be circular.
\end{physics}

\begin{warning}
\textbf{Stale docstrings about the trust boundary (documentation only).} Several module docstrings on the critical path still say that the
project axiom does \emph{not} cover the PBW step and ``must not be cited here'', or that the independence of levels is an
unproved PBW input: \leanpath{ZhuVermaAnArchDLM3Survival.lean} (``the project's sanctioned \lean{pbw} axiom is about affine
Lie algebras, not $\lthirteenUE(\lthirteenVh)$ over $\lthirteenUE(P_0)$; it must not be cited here''), and
\leanpath{ZhuVermaAnArchDLM3Zero.lean} (``independence of the levels is DLM3's PBW invocation \ldots which the sanctioned
\lean{pbw} axiom does \textbf{not} cover''). Both are superseded: \lean{PBWPrinciple.inducedBottom} covers exactly this step
(\lean{voaPBWBottomInjective}), and level independence is the theorem \lean{levCover_iSupIndep}. The mathematics of the
files is unaffected. (Both docstrings, and the word ``axiom'' in them, are further outdated since \lean{VOA.pbw} became a
theorem.)
\end{warning}

%======================================================================
\section{The tower: circle products, \texorpdfstring{$O_n$}{O\_n}, \texorpdfstring{$A_n$}{A\_n} and the Dong--Jiang spaces \texorpdfstring{$A_{n,m}$}{A\_\{n,m\}}}\label{l13:sec:tower}
%======================================================================

All products in this section are residues of the form
\begin{equation}\label{l13:eq:res}
 \Res_z Y(a,z)b\,\frac{(1+z)^{\wt a+s}}{z^{p}} = \sum_{i\ge0}\binom{\wt a+s}{i}\,a_{(i-p)}b\qquad(a\text{ homogeneous}),
\end{equation}
for integers $s$ and $p$, extended bilinearly to inhomogeneous $a$ through the grading $V = \bigoplus_kV_k$. The sum is finite
when $\wt a+s\ge0$ (only $i\le\wt a+s$ contribute). The library has one generic implementation, \lean{resHom W s p a b}, and
expresses every product of the tower through it.

\subsection{The generic residue and the circle products}

\begin{definition}[Generic residue; $n$-th circle product; $O_n$ and $A_n$ of the library]\label{l13:def:circ}
\lean{resHom W s p a b} is \eqref{l13:eq:res}. The $n$-th circle product is $a\circ_nb = \Res_zY(a,z)b\,(1+z)^{\wt a+n}/z^{2n+2}$,
i.e.\ \eqref{l13:eq:res} with $s = n$, $p = 2n+2$. The library's $O_n^{\mathrm{lib}}(V)$ is the span of all $a\circ_nb$ and
$A_n^{\mathrm{lib}}(V) = V/O_n^{\mathrm{lib}}(V)$ (Definition~\ref{l12:def:an} of Chapter~\ref{ch:L12}).
\end{definition}
\noindent Files \leanpath{VOA/\allowbreak Modules/\allowbreak ZhuVermaAnBimodule.lean} and
\leanpath{VOA/\allowbreak Modules/\allowbreak ZhuVermaAn.lean}.
\begin{leancode}
noncomputable def resFamily (W : VertexOperatorAlgebra 𝒱) (s p : ℤ) (b : 𝒱) (k : ℤ) :
    (W.grading k) →ₗ[ℂ] 𝒱 where
  toFun ak := ∑ i ∈ Finset.range ((k + s).toNat + 1),
    cbz (k + s) i • W.smul (ak : 𝒱) ((i : ℤ) - p) b
noncomputable def resHom (W : VertexOperatorAlgebra 𝒱) (s p : ℤ) (a b : 𝒱) : 𝒱 :=
  letI := DirectSum.IsInternal.chooseDecomposition W.grading W.isInternal
  (DFinsupp.lsum ℂ (resFamily W s p b)) (DirectSum.decompose W.grading a)
noncomputable def circFamily (W : VertexOperatorAlgebra 𝒱) (n : ℕ) (b : 𝒱) (k : ℤ) :
    (W.grading k) →ₗ[ℂ] 𝒱 where
  toFun ak := ∑ i ∈ Finset.range ((k + n).toNat + 1),
    cbz (k + (n : ℤ)) i • W.smul (ak : 𝒱) ((i : ℤ) - 2 * (n : ℤ) - 2) b
noncomputable def circ (W : VertexOperatorAlgebra 𝒱) (n : ℕ) (a b : 𝒱) : 𝒱 :=
  letI := DirectSum.IsInternal.chooseDecomposition W.grading W.isInternal
  (DFinsupp.lsum ℂ (circFamily W n b)) (DirectSum.decompose W.grading a)
def On (W : VertexOperatorAlgebra 𝒱) (n : ℕ) : Submodule ℂ 𝒱 :=
  Submodule.span ℂ {x | ∃ a b : 𝒱, x = circ W n a b}
abbrev An (W : VertexOperatorAlgebra 𝒱) (n : ℕ) : Type u := 𝒱 ⧸ On W n
\end{leancode}

\begin{dictionary}
\lean{cbz m i} is $\binom mi$ for $m,i\in\Z$ (zero for $i<0$), \lean{W.smul a j b} is $a_{(j)}b$. \lean{resFamily W s p b k}
is the formula on the weight-$k$ slice $V_k$ (the sum runs over $i = 0,\dots,\max(k+s,0)$, beyond which the binomial
vanishes when $k+s\ge0$; when $k+s<0$ the range is just $i=0$, \lean{(k + s).toNat = 0}). \lean{DirectSum.decompose W.grading a}
splits $a$ into homogeneous components and \lean{DFinsupp.lsum} sums the slice maps over them: this is ``extend bilinearly over
the grading''. The truncation of the sum at $\max(k+s,0)$ when $k+s<0$ is a convention that never matters below, because
every use has $s\ge0$ and $k\ge0$ (library VOAs have $V_k = 0$ for $k<0$). \lean{circ} is \lean{resHom W n (2n+2)} written
out separately (\lean{circ_eq_resHom}).
\end{dictionary}

\begin{physics}
A residue $\Res_zY(a,z)b\,(1+z)^{\wt a+s}/z^p$ is a weighted combination of OPE coefficients $a_{(j)}b$, i.e.\ of the fields
appearing in the operator product $a(z)b(0)$. Such combinations arise when one computes how a mode of $a$ acts on a state
produced by $b$ \emph{at a fixed energy}: on the cylinder $z = e^{w}-1$, the factor $(1+z)^{\wt a}$ is the conformal
transformation from the plane, and $z^{-p}$ selects how many levels the mode moves. The circle products are the combinations that
are guaranteed to act by zero on the lowest $n+1$ levels of every sector.
\end{physics}

\begin{theorem}[$O_0 = O$]\label{l13:thm:On0}
$O_0^{\mathrm{lib}}(V) = O(V)$; hence $A_0^{\mathrm{lib}}(V) = A(V)$.
\end{theorem}
\noindent File \leanpath{VOA/\allowbreak Modules/\allowbreak ZhuVermaAnDescent15.lean}.
\begin{leancode}
theorem On_zero_eq_O (W : VertexOperatorAlgebra 𝒱) : On W 0 = O W := by
\end{leancode}
\textbf{Proof on paper.} Both are spans; $a\circ_0b$ is exactly Zhu's generator $\Res_zY(a,z)b(1+z)^{\wt a}/z^2$ of $O(V)$
(Chapter~\ref{ch:L12}), so the generating sets agree. \qed

\begin{status}
Proved, unconditional; standard three axioms.
\end{status}

\begin{physics}
$O(V)$ is the space of fields whose zero modes vanish on every primary; at level $n$, $O_n(V)$ is the space of fields
whose zero modes vanish on the first $n$ excited levels of every sector.
\end{physics}

\subsection{The corrected \texorpdfstring{$O_n(V)$}{O\_n(V)} and the product \texorpdfstring{$*_n$}{*\_n}}

\begin{definition}[Dong--Li--Mason's $O_n(V)$ and $*_n$]\label{l13:def:onfull}
$O_n(V) = \operatorname{span}\{a\circ_nb\} + \operatorname{span}\{(L_{-1}+L_0)a\}$ and $A_n(V) = V/O_n(V)$. For homogeneous $a$,
\begin{equation*}
 a*_nb = \sum_{j=0}^n(-1)^j\binom{n+j}{n}\Res_zY(a,z)b\,\frac{(1+z)^{\wt a+n}}{z^{n+j+1}}.
\end{equation*}
\end{definition}
\noindent Files \leanpath{VOA/\allowbreak Modules/\allowbreak ZhuVermaAnArch.lean},
\leanpath{VOA/\allowbreak Modules/\allowbreak ZhuVermaAnArchAlgebra.lean},
\leanpath{VOA/\allowbreak Modules/\allowbreak ZhuVermaAnBimodule.lean} and
\leanpath{VOA/\allowbreak Modules/\allowbreak ZhuVermaAnArchAnIdeal.lean}.
\begin{leancode}
def OnL (W : VertexOperatorAlgebra 𝒱) : Submodule ℂ 𝒱 :=
  Submodule.span ℂ {x | ∃ a : 𝒱, x = W.L (-1) a + W.L 0 a}
def OnFull (W : VertexOperatorAlgebra 𝒱) (n : ℕ) : Submodule ℂ 𝒱 :=
  On W n ⊔ OnL W
abbrev AnFull (W : VertexOperatorAlgebra 𝒱) (n : ℕ) : Type u := 𝒱 ⧸ OnFull W n
noncomputable def leftProd (W : VertexOperatorAlgebra 𝒱) (n : ℕ) (a b : 𝒱) : 𝒱 :=
  ∑ j ∈ Finset.range (n + 1),
    ((-1 : ℂ) ^ j * cbz ((j : ℤ) + n) (n : ℤ)) •
      resHom W (n : ℤ) ((n : ℤ) + (j : ℤ) + 1) a b
structure AnIdeal (W : VertexOperatorAlgebra 𝒱) (n : ℕ) : Prop where
  /-- **Left ideal** `Oₙ *ₙ 𝒱 ⊆ Oₙ`: `g` in the first (acting) slot, over the corrected `Oₙ`. -/
  left : ∀ g ∈ OnFull W n, ∀ b : 𝒱, leftProd W n g b ∈ OnFull W n
  /-- **Right ideal** `𝒱 *ₙ Oₙ ⊆ Oₙ`: `g` in the second slot ranges over the corrected `Oₙ`. -/
  right : ∀ a : 𝒱, ∀ g ∈ OnFull W n, leftProd W n a g ∈ OnFull W n
theorem anIdeal_holds (W : VertexOperatorAlgebra 𝒱) (n : ℕ) : AnIdeal W n :=
\end{leancode}

\begin{dictionary}
\lean{⊔} is the sum of subspaces. \lean{OnFull W n} is exactly Dong--Li--Mason's (2.1); \lean{leftProd W n a b} is their
$a*_nb$ with the external coefficients $(-1)^j\binom{j+n}{n}$ (\lean{cbz ((j : ℤ) + n) (n : ℤ)}) and pole order $n+j+1$; at
$n = 0$ it is Zhu's product (\lean{leftProd_zero : leftProd W 0 a b = zhuMul W a b}). \lean{AnIdeal W n} says $O_n(V)$ is a
two-sided ideal for $*_n$, and \lean{anIdeal_holds} proves it for every $n$ (Dong--Li--Mason, Lemma~2.2). Associativity of
$*_n$ modulo $O_n(V)$ for $n\ge1$ is \emph{not} proved (the library reduces it to a named target \lean{AssocHom W n} and proves
only $n = 0$); it is not needed for the headline.
\end{dictionary}

\begin{warning}
\textbf{Status of the $O_n$ defect recorded in Chapter~\ref{ch:L12} (entry 49 of the project's defect list).} The circle-only
\lean{On W n} omits Dong--Li--Mason's generators $(L_{-1}+L_0)a$ for $n\ge1$. This chapter adds two facts. (i) The library also
defines the correct ideal, \lean{OnFull} $=$ \lean{On} $\sqcup$ \lean{OnL}, and it is \lean{OnFull} (through the Dong--Jiang
ideals below) that the construction of $M(U)$ uses; \lean{On} enters the headline's proof only at $n = 0$, where
\lean{OnFull_zero_eq_On_zero} and \lean{On_zero_eq_O} show all versions agree with $O(V)$. So the defect does not affect
\lean{zhuBijective}. (ii) The docstring of \lean{OnFull} asserts that for $n\ge1$ it is ``strictly larger than the circle-only
\lean{On W n}''; no such strictness is proved in the library, and we do not know whether it holds for all $V$ (documentation only, unverified
claim).
\end{warning}

\begin{status}
\lean{anIdeal_holds}: proved, unconditional, standard three axioms. Associativity of $A_n(V)$, $n\ge1$: open (named
target), off the critical path.
\end{status}

\begin{physics}
$A_n(V)$ is the algebra of ``generalized zero modes'' $o_{n,n}(a)$ acting on level $n$ of every sector; $*_n$ is their
composition law. That $O_n(V)$ is an ideal is the consistency of this composition with the null relations.
\end{physics}

\subsection{The Dong--Jiang bimodules \texorpdfstring{$A_{n,m}(V)$}{A\_\{n,m\}(V)}}

The induced module needs the maps from level $m$ to level $n$, not only the algebras acting on one level. Dong and Jiang
define, for homogeneous $u$,
\begin{align}
 u*^n_{m,p}v &= \sum_{i=0}^{p}(-1)^i\binom{m+n-p+i}{i}\Res_zY(u,z)v\,\frac{(1+z)^{\wt u+m}}{z^{m+n-p+i+1}},\label{l13:eq:djprod}\\
 u\circ^n_mv &= \Res_zY(u,z)v\,\frac{(1+z)^{\wt u+m}}{z^{n+m+2}},\qquad
 \text{grading generator } L_{-1}u+(L_0+m-n)u.\notag
\end{align}
The special case $u*^n_mv := u*^n_{m,m}v$ is the right $A_m(V)$-action and $u\,\bar*^n_m\,v := u*^n_{m,n}v$ the left $A_n(V)$-action;
at $m = n$ both are $*_n$.

\begin{definition}[Dong--Jiang ideals and spaces]\label{l13:def:onm}
$O'_{n,m}(V)$ is the span of all $u\circ^n_mv$ and all $L_{-1}u+(L_0+m-n)u$. $O''_{n,m}(V)$ is the span of the
\emph{associator defects}
\begin{equation*}
 u*^n_{m,p_3}\Bigl((a*^{p_3}_{p_1,p_2}b)*^{p_3}_{m,p_1}c - a*^{p_3}_{m,p_2}(b*^{p_2}_{m,p_1}c)\Bigr),\qquad a,b,c,u\in V,\ p_i\in\N.
\end{equation*}
$O'''_{n,m}(V)$ is the span of $(a*^n_{p,p}g)*^n_{m,p}v$ with $g\in O'_{p,p}(V)$, $p\in\N$. The full ideal is
$O_{n,m}(V) = O'_{n,m}+O''_{n,m}+O'''_{n,m}$ and $A_{n,m}(V) = V/O_{n,m}(V)$.
\end{definition}
\noindent Files \leanpath{VOA/\allowbreak Modules/\allowbreak ZhuVermaAnArchAnmFull.lean} and
\leanpath{VOA/\allowbreak Modules/\allowbreak ZhuVermaAnArchAnm.lean}.
\begin{leancode}
noncomputable def anmGen (W : VertexOperatorAlgebra 𝒱) (n m p : ℕ) (u v : 𝒱) : 𝒱 :=
  ∑ i ∈ Finset.range (p + 1),
    ((-1 : ℂ) ^ i * cbz ((m : ℤ) + (n : ℤ) - (p : ℤ) + i) (i : ℤ))
      • resHom W (m : ℤ) ((m : ℤ) + (n : ℤ) - (p : ℤ) + (i : ℤ) + 1) u v
noncomputable def anmRight (W : VertexOperatorAlgebra 𝒱) (n m : ℕ) (u v : 𝒱) : 𝒱 :=
  ∑ i ∈ Finset.range (m + 1),
    ((-1 : ℂ) ^ i * cbz ((n : ℤ) + i) (i : ℤ)) • resHom W (m : ℤ) ((n : ℤ) + (i : ℤ) + 1) u v
noncomputable def anmCirc (W : VertexOperatorAlgebra 𝒱) (n m : ℕ) (u v : 𝒱) : 𝒱 :=
  resHom W (m : ℤ) ((n : ℤ) + (m : ℤ) + 2) u v
noncomputable def anmLgen (W : VertexOperatorAlgebra 𝒱) (n m : ℕ) (u : 𝒱) : 𝒱 :=
  W.L (-1) u + W.L 0 u + ((m : ℂ) - (n : ℂ)) • u
def OnmL (W : VertexOperatorAlgebra 𝒱) (n m : ℕ) : Submodule ℂ 𝒱 :=
  Submodule.span ℂ {x | ∃ u : 𝒱, x = anmLgen W n m u}
def Onm (W : VertexOperatorAlgebra 𝒱) (n m : ℕ) : Submodule ℂ 𝒱 :=
  Submodule.span ℂ {x | ∃ u v : 𝒱, x = anmCirc W n m u v} ⊔ OnmL W n m
noncomputable def assocDefect (W : VertexOperatorAlgebra 𝒱) (n m p₁ p₂ p₃ : ℕ) (a b c u : 𝒱) : 𝒱 :=
  anmGen W n m p₃ u
    (anmGen W p₃ m p₁ (anmGen W p₃ p₁ p₂ a b) c
      - anmGen W p₃ m p₂ a (anmGen W p₂ m p₁ b c))
def OnmAssoc (W : VertexOperatorAlgebra 𝒱) (n m : ℕ) : Submodule ℂ 𝒱 :=
  Submodule.span ℂ
    {x | ∃ (p₁ p₂ p₃ : ℕ) (a b c u : 𝒱), x = assocDefect W n m p₁ p₂ p₃ a b c u}
noncomputable def levelMix (W : VertexOperatorAlgebra 𝒱) (n m p : ℕ) (a g v : 𝒱) : 𝒱 :=
  anmGen W n m p (anmGen W n p p a g) v
def OnmLevel (W : VertexOperatorAlgebra 𝒱) (n m : ℕ) : Submodule ℂ 𝒱 :=
  Submodule.span ℂ
    {x | ∃ (p : ℕ) (a v : 𝒱) (g : 𝒱), g ∈ Onm W p p ∧ x = levelMix W n m p a g v}
def OnmFull (W : VertexOperatorAlgebra 𝒱) (n m : ℕ) : Submodule ℂ 𝒱 :=
  Onm W n m ⊔ OnmAssoc W n m ⊔ OnmLevel W n m
abbrev AnmFull (W : VertexOperatorAlgebra 𝒱) (n m : ℕ) : Type u := 𝒱 ⧸ OnmFull W n m
\end{leancode}

\begin{dictionary}
\textbf{Line-by-line.} \lean{anmGen W n m p u v} is \eqref{l13:eq:djprod}: the range \lean{Finset.range (p + 1)} is
$i = 0,\dots,p$; the coefficient is $(-1)^i\binom{m+n-p+i}{i}$; the residue has shift $s = m$ and pole $m+n-p+i+1$.
\lean{anmRight W n m} is the special case $p = m$ (the library proves \lean{anmGen_eq_anmRight}). \lean{anmCirc} and \lean{anmLgen}
are $u\circ^n_mv$ and the shifted grading generator; note the shift $+(m-n)$, which at $m = n$ gives back $(L_{-1}+L_0)u$.
\lean{Onm} is $O'_{n,m}$, \lean{OnmAssoc} is $O''_{n,m}$, \lean{OnmLevel} is $O'''_{n,m}$ (with $g$ ranging over
$O'_{p,p} = $ \lean{Onm W p p}), and \lean{OnmFull} is their sum. Note that the indices $p_i$ range over all of $\N$, as in
Dong--Jiang: $O''_{n,m}$ and $O'''_{n,m}$ involve \emph{all} levels of the tower at once.

\textbf{Fidelity to Dong--Jiang.} We compared each Lean definition with the formulas of Dong--Jiang (\emph{Math.\ Z.} 259, §2)
as transcribed, with line references to the arXiv text, in the docstring of \leanpath{ZhuVermaAnArchAnmFull.lean}; they agree
term by term. (We did not re-derive Dong--Jiang's formulas themselves.) The library
also proves the expected diagonal identity \lean{Onm_diag_eq_OnFull : Onm W n n = OnFull W n} ($O'_{n,n} = O_n$) and the
$(0,0)$ anchor \lean{Onm_zero_zero_eq_O : Onm W 0 0 = O W}.

\textbf{Why the full ideal.} Dong and Jiang note that they cannot prove $O_{n,m} = O'_{n,m}$ (their Remark~2.10); $V/O'_{n,m}$
is not known to be a bimodule. The library follows them and works with \lean{OnmFull}; this is the conservative choice (a
bigger ideal, a smaller quotient, more relations available).
\end{dictionary}

\begin{physics}
The products $*^n_{m,p}$ are the composition laws of transition operators between energy levels; the ideals $O_{n,m}(V)$ are the
relations they satisfy in every sector. Writing them down explicitly is the algebraic counterpart of computing all three-point
functions $\langle\text{level }n|u(z)|\text{level }m\rangle$ allowed by the chiral algebra.
\end{physics}

\begin{definition}[The diagonal collapse, Dong--Jiang Proposition~4.6]\label{l13:def:djann}
\lean{DJAnnEqAn W m} is the statement $O_{m,m}(V)\subseteq O_m(V)$, equivalently (since $O_m\subseteq O_{m,m}$ is proved)
$A_{m,m}(V) = A_m(V)$.
\end{definition}
\noindent File \leanpath{VOA/\allowbreak Modules/\allowbreak ZhuVermaAnArchDJCarrier.lean}.
\begin{leancode}
def DJAnnEqAn (W : VertexOperatorAlgebra 𝒱) (m : ℕ) : Prop :=
  OnmFull W m m ≤ OnFull W m
\end{leancode}

\begin{dictionary}
The reverse inclusion \lean{onFull_le_OnmFull_diag} is a theorem, so \lean{DJAnnEqAn W m} says exactly that the associator and
level-mixing summands at $(m,m)$ add nothing to $O_m(V)$. Only $m = 0$ is used: \lean{DJAnnEqAn W 0} is $O_{0,0}(V)\subseteq O(V)$
(using \lean{OnFull_zero_eq_On_zero} and \lean{On_zero_eq_O}). It is derived from \lean{Pr} in Theorem~\ref{l13:thm:gate}.
\end{dictionary}

\begin{status}
All definitions above; no hypotheses. \texttt{\#print axioms VOA.OnmFull} = standard three.
\end{status}

\begin{physics}
$A_{n,m}(V)$ is the space of ``transition operators'': for a field $u$, the mode that maps level $m$ of a sector to level $n$
is $o_{n,m}(u) = u_{(\wt u-1+m-n)}$, and $O_{n,m}(V)$ is the space of fields all of whose level-$m$-to-level-$n$ modes vanish in
every sector. The three summands are, respectively: the analogue of Zhu's relation; the requirement that composing transition
operators through intermediate levels is associative; and the requirement that the relations of intermediate levels are
respected. $A_{n,0}(V)\otimes_{A(V)}U$ is then ``the span of all level-$n$ states $o_{n,0}(u)|w\rangle$ obtainable from the
primary multiplet $U$''.
\end{physics}

%======================================================================
\section{The induced module \texorpdfstring{$M(U)$}{M(U)} and the gate \texorpdfstring{$A_{0,0}(V) = A(V)$}{A\_\{0,0\}(V)=A(V)}}\label{l13:sec:MU}
%======================================================================

\subsection{The carrier}

\begin{definition}[Level-$m$ modules; the Dong--Jiang carrier]\label{l13:def:djcarrier}
An $A_m(V)$-module (\lean{AmMod W m}) is a vector space $U$ with a linear $\rho:A_m(V)\to\End U$ such that
$\rho[\lthirteenvac] = 1$ and $\rho[u*_mv] = \rho[u]\rho[v]$. For such $U$ and $n\in\N$ put
\begin{align*}
 M(U)(n) = A_{n,m}(V)\otimes_{A_m(V)}U &:= \bigl(A_{n,m}(V)\otimes_\C U\bigr)\big/
 \operatorname{span}\{[x*^n_ma]\otimes w - [x]\otimes\rho[a]w\},\\
 M(U) &= \bigoplus_{n\ge0}M(U)(n).
\end{align*}
Every $A(V)$-module is an $A_0(V)$-module via $A_0(V) = A(V)$ (\lean{AMod.toAmModZero}).
\end{definition}
\noindent Files \leanpath{VOA/\allowbreak Modules/\allowbreak ZhuVermaAnArchDJCarrier.lean},
\leanpath{VOA/\allowbreak Modules/\allowbreak ZhuVermaAnArchDJCarrierGraded.lean} and
\leanpath{VOA/\allowbreak Modules/\allowbreak ZhuVermaAnArchLU.lean}.
\begin{leancode}
structure AmMod (W : VertexOperatorAlgebra 𝒱) (m : ℕ) where
  /-- The underlying carrier. -/
  carrier : Type u
  [addCommGroup : AddCommGroup carrier]
  [module : Module ℂ carrier]
  /-- The `ℂ`-linear action of `A_m(𝒱)`. -/
  ρ : AnFull W m →ₗ[ℂ] Module.End ℂ carrier
  /-- The unit `[𝟙]` acts as the identity. -/
  ρ_one : ρ (mkAnFull W m W.vacuum) = 1
  /-- Multiplicativity on every representable Frenkel–Zhu product `*_m`. -/
  ρ_mul : ∀ u v : 𝒱,
    ρ (mkAnFull W m (leftProd W m u v)) = ρ (mkAnFull W m u) * ρ (mkAnFull W m v)
noncomputable def AMod.toAmModZero (U : AMod W) : AmMod W 0 where
  carrier := U.carrier
  ρ := (U.ρ).comp (anFullZeroEquiv W).toLinearMap
  ρ_one := U.ρ_one
  ρ_mul u v := by
    show U.ρ (VOA.mk W (leftProd W 0 u v))
      = U.ρ (VOA.mk W u) * U.ρ (VOA.mk W v)
    rw [leftProd_zero]
    exact U.ρ_mul u v
abbrev DJCover (W : VertexOperatorAlgebra 𝒱) (U : AmMod W m) (n : ℕ) : Type u :=
  AnmFull W n m ⊗[ℂ] U.carrier
def djBalancingSet (W : VertexOperatorAlgebra 𝒱) (U : AmMod W m) (n : ℕ) :
    Set (DJCover W U n) :=
  {z | ∃ (x a : 𝒱) (w : U.carrier),
    z = mkAnmFull W n m (anmRight W n m x a) ⊗ₜ[ℂ] w
        - mkAnmFull W n m x ⊗ₜ[ℂ] U.ρ (mkAnFull W m a) w}
noncomputable def djBalancing (W : VertexOperatorAlgebra 𝒱) (U : AmMod W m) (n : ℕ) :
    Submodule ℂ (DJCover W U n) :=
  Submodule.span ℂ (djBalancingSet W U n)
abbrev DJPiece (W : VertexOperatorAlgebra 𝒱) (U : AmMod W m) (n : ℕ) : Type u :=
  DJCover W U n ⧸ djBalancing W U n
abbrev DJCarrier (W : VertexOperatorAlgebra 𝒱) (U : AmMod W m) : Type u :=
  ⨁ n : ℕ, DJPiece W U n
noncomputable def DJCarrier.deg (W : VertexOperatorAlgebra 𝒱) (U : AmMod W m) (n : ℕ) :
    Submodule ℂ (DJCarrier W U) := LinearMap.range (DJCarrier.incl W U n)
\end{leancode}

\begin{dictionary}
\textbf{The balanced tensor product.} \lean{AmMod} mirrors \lean{AMod}, with \lean{AnFull W m} $= A_m(V)$ and
\lean{leftProd W m} $=*_m$. Since associativity of $A_m(V)$ is not available for $m\ge1$, the library cannot form Mathlib's
\lean{TensorProduct (AnFull W m)}; instead \lean{DJPiece W U n} is the $\C$-tensor product divided by the span of the
\emph{representative-level} balancing relations $[x*^n_ma]\otimes w - [x]\otimes\rho[a]w$. This is the universal space in which
Dong--Jiang's balancing holds; when $A_m(V)$ is a ring acting on the right of $A_{n,m}(V)$ it coincides with
$A_{n,m}(V)\otimes_{A_m(V)}U$ (the relations generate the same subspace). On the critical path only $m = 0$ occurs, where
$A_0(V) = A(V)$ \emph{is} a ring and the right action of $A(V)$ on $A_{n,0}(V)$ descends unconditionally
(\lean{djRightSlot_zero}); so there the construction is literally $A_{n,0}(V)\otimes_{A(V)}U$.

\textbf{The direct sum.} \lean{⨁ n : ℕ, DJPiece W U n} is Mathlib's external direct sum (finitely supported families);
\lean{DJCarrier.deg W U n} is the image of the $n$-th summand. It is an internal direct-sum decomposition
(\lean{DJCarrier.isInternal_deg}).

\textbf{\lean{toAmModZero}.} Transports the $A(V)$-action along \lean{anFullZeroEquiv : AnFull W 0 ≃ A W}; its
multiplicativity is \lean{leftProd_zero} ($*_0 = *$) and Zhu's $\rho_{\mathrm{mul}}$. Nothing is lost.
\end{dictionary}

\begin{physics}
$M(U)(n)$ is the space of level-$n$ descendants of the primary multiplet $U$, written as ``(transition operator from level $0$
to level $n$) applied to (a primary)''. The balancing relation says that applying a zero mode to the primary first, or
absorbing it into the transition operator, gives the same state.
\end{physics}

\subsection{The action of \texorpdfstring{$V$}{V}}

\begin{definition}[Dong--Jiang's modes]\label{l13:def:djmode}
For homogeneous $u$ and $p\in\Z$, the mode $u_{(p)}$ maps $M(U)(n)$ to $M(U)(N)$, $N = \wt u - p - 1 + n$, by
\begin{equation*}
 u_{(p)}\bigl([v]\otimes w\bigr) = [u*^N_{m,n}v]\otimes w\quad(N\ge0),\qquad u_{(p)} = 0\text{ on }M(U)(n)\text{ if }N<0,
\end{equation*}
and is extended linearly over the grading of $u$. For $m = 0$ this is \lean{djModeZero W U p u}.
\end{definition}
\noindent File \leanpath{VOA/\allowbreak Modules/\allowbreak ZhuVermaAnArchDJAction.lean} (and
\leanpath{VOA/\allowbreak Modules/\allowbreak ZhuVermaAnArchMUFullMode.lean} for the target level).
\begin{leancode}
def DJLeftSlot (W : VertexOperatorAlgebra 𝒱) (N m n : ℕ) : Prop :=
  ∀ u : 𝒱, ∀ g ∈ OnmFull W n m, anmGen W N m n u g ∈ OnmFull W N m
def DJActionDescent (W : VertexOperatorAlgebra 𝒱) (m : ℕ) : Prop :=
  ∀ N n : ℕ, DJLeftSlot W N m n
theorem djActionDescent_zero (W : VertexOperatorAlgebra 𝒱) : DJActionDescent W 0 :=
  fun N n => djLeftSlot_zero W N n
noncomputable def djCoverAct (W : VertexOperatorAlgebra 𝒱) (U : AmMod W m) (N n : ℕ)
    (h : DJLeftSlot W N m n) (u : 𝒱) : DJCover W U n →ₗ[ℂ] DJCover W U N :=
  TensorProduct.map (AnmFull.leftGen W N m n h u) LinearMap.id
noncomputable def djPieceAct (W : VertexOperatorAlgebra 𝒱) (U : AmMod W m) (N n : ℕ)
    (h : DJLeftSlot W N m n) (u : 𝒱) : DJPiece W U n →ₗ[ℂ] DJPiece W U N :=
  (djBalancing W U n).liftQ ((mkDJPiece W U N).comp (djCoverAct W U N n h u))
    (djBalancing_le_ker_coverAct U N n h u)
def muTargetLevel (k p : ℤ) (n : ℕ) : ℤ := k - p - 1 + (n : ℤ)
noncomputable def djLevelIncl (W : VertexOperatorAlgebra 𝒱) (U : AmMod W m)
    (hD : DJActionDescent W m) (k p : ℤ) (n : ℕ) (u : 𝒱) :
    DJPiece W U n →ₗ[ℂ] DJCarrier W U :=
  if 0 ≤ muTargetLevel k p n then
    (DJCarrier.incl W U (muTargetLevel k p n).toNat).comp
      (djPieceAct W U (muTargetLevel k p n).toNat n (hD _ n) u)
  else 0
noncomputable def djMode (W : VertexOperatorAlgebra 𝒱) (U : AmMod W m)
    (hD : DJActionDescent W m) (p : ℤ) (a : 𝒱) : Module.End ℂ (DJCarrier W U) :=
  letI := DirectSum.IsInternal.chooseDecomposition W.grading W.isInternal
  (DFinsupp.lsum ℂ (djModeFamily W U hD p)) (DirectSum.decompose W.grading a)
noncomputable def djModeZero (W : VertexOperatorAlgebra 𝒱) (U : AmMod W 0) (p : ℤ) (a : 𝒱) :
    Module.End ℂ (DJCarrier W U) :=
  djMode W U (djActionDescent_zero W) p a
\end{leancode}

\begin{dictionary}
\lean{DJLeftSlot W N m n} is Dong--Jiang's Lemma~2.8 (left inclusion): $V*^N_{m,n}O_{n,m}(V)\subseteq O_{N,m}(V)$, which makes
$[v]\mapsto[u*^N_{m,n}v]$ well defined $A_{n,m}(V)\to A_{N,m}(V)$ (\lean{AnmFull.leftGen}). At base level $m = 0$ it is
\emph{proved} for all $N,n$ (\lean{djActionDescent_zero}); in general it is reduced to a named residual. \lean{djCoverAct}
applies this map on the first tensor factor; \lean{djPieceAct} descends it to the balanced piece (the balancing relations are
respected, \lean{djBalancing_le_ker_coverAct}). \lean{muTargetLevel k p n} is $k-p-1+n$; \lean{djLevelIncl} is the case split of
Definition~\ref{l13:def:djmode} for $u\in V_k$; \lean{djModeFamily}/\lean{djMode} assemble the components over levels (by the
universal property of \lean{⨁}) and over the grading of $u$. \lean{djModeZero} is the hypothesis-free $m = 0$ instance. So
\lean{djModeZero W U p u} is the operator $u_{(p)}$ on $M(U)$; in particular \lean{djModeZero W U 1 W.conformalVector} is
$L_0 = \omega_{(1)}$.
\end{dictionary}

\begin{physics}
The mode $u_{(p)}$ of a field of weight $\wt u$ changes the energy by $\wt u-p-1$; acting on a level-$n$ descendant it produces a
level-$(\wt u-p-1+n)$ descendant, computed by composing transition operators. If the target level would be negative the state is
annihilated --- the statement that there are no states below the primaries.
\end{physics}

\begin{theorem}[$M(U)$ is a weak module; Dong--Jiang, Theorem~4.10]\label{l13:thm:djweak}
For every $A(V)$-module $U$ (more precisely, every $A_0(V)$-module), the modes of Definition~\ref{l13:def:djmode} make $M(U)$
a weak $V$-module. The decomposition $M(U) = \bigoplus_nM(U)(n)$ is a grading: $u_{(p)}M(U)(n)\subseteq M(U)(\wt u-p-1+n)$.
\end{theorem}
\noindent File \leanpath{VOA/\allowbreak Modules/\allowbreak ZhuVermaAnArchDJCarrierWeak.lean}.
\begin{leancode}
noncomputable def djCarrierWeak (W : VertexOperatorAlgebra 𝒱) (U : AmMod W 0) :
    WeakModule W (DJCarrier W U) :=
  weakModuleOfRep (djRhoZero W U) (djActModeTruncated W U) (djVacuumCreation W U)
    (isModuleJacobi_djRhoZero W U)
\end{leancode}

\begin{dictionary}
\lean{WeakModule W M} is the library's weak module (Chapter~\ref{ch:01}): a vertex operator $Y_M$ satisfying truncation,
vacuum creation $\lthirteenvac_{(m)} = \delta_{m,-1}$, and the Jacobi (Borcherds) identity. \lean{weakModuleOfRep ρ htr hvac hjac}
builds one from a representation $\rho$ of $\lthirteenVh$ (\lean{VHatRep W M} $=$ \lean{VHatCarrier W →ₗ⁅ℂ⁆ Module.End ℂ M},
with $u_{(p)} = \rho(u(p))$) satisfying truncation (\lean{ActModeTruncated}), vacuum creation and the module Jacobi identity
(\lean{IsModuleJacobi}). Here $\rho = $ \lean{djRhoZero W U}, whose modes are \lean{djModeZero}. All three obligations are
\emph{theorems} for every $U$; there is no hypothesis. The library records the action as \lean{djCarrierWeak_act}:
the weak-module mode $u_{(p)}$ is \lean{djModeZero W U p u}.
\end{dictionary}

\textbf{Proof on paper.} (i) $\rho$ is a Lie-algebra representation of $\lthirteenVh$: the translation relation
$(L_{-1}u)_{(p)} = -p\,u_{(p-1)}$ holds because the grading generator $L_{-1}u+(L_0+m-n)u$ lies in $O_{N,m}$, and the
commutator formula $[u_{(p)},v_{(q)}] = \sum_i\binom pi(u_{(i)}v)_{(p+q-i)}$ is Dong--Jiang's Lemma~4.7, a binomial identity
between the products $*^N_{m,n}$ modulo $O_{N,m}$. (ii) Truncation: $u_{(p)}$ lowers the level for $p$ large and is zero once
the target level is negative. (iii) Vacuum: $\lthirteenvac*^N_{m,n}v = v$ at $N = n$ and the other modes vanish. (iv) The Jacobi
identity of a module follows from the commutator formula plus the \emph{iterate} formula
$(u_{(\ell)}v)_{(N)} = \sum_i(-1)^i\binom\ell i\bigl(u_{(\ell-i)}v_{(N+i)} - (-1)^\ell v_{(\ell+N-i)}u_{(i)}\bigr)$; for $\ell\ge0$ it
is a consequence of the commutator formula, and for $\ell<0$ it is Dong--Jiang's Lemma~4.9, which holds on $M(U)$ because the
associator defects (the summand $O''$) have been divided out. \qed

\begin{status}
\lean{djCarrierWeak}: definition with all obligations proved, unconditional; \texttt{\#print axioms} = standard three.
\end{status}

\begin{physics}
The descendants of a primary multiplet, with the fields acting by the rules forced by the OPE, form a consistent
representation of the full chiral algebra (including normal-ordered products) --- \emph{provided} the associativity of the
OPE has been imposed, which is what quotienting by $O''$ does. This is the chiral-algebra version of ``Verma modules are
modules''.
\end{physics}

\subsection{Graded lowest-weight weak modules}

\begin{definition}[Graded lowest-weight weak module]\label{l13:def:gradedlw}
A \emph{graded lowest-weight weak module} of lowest weight $h\in\C$ is a weak module $M$ with a decomposition
$M = \bigoplus_{\alpha\in\C}M_\alpha$ such that $L_0 = \alpha$ on $M_\alpha$, the real parts of the occurring $\alpha$ are bounded
below, every occurring $\alpha$ lies in $h+\N$, and $M_h\neq0$. It is \emph{irreducible} if $M\neq0$ and its only
mode-stable subspaces are $0$ and $M$.
\end{definition}
\noindent Files \leanpath{VOA/\allowbreak Modules/\allowbreak ZhuVermaAnArchDJCarrierWeak.lean},
\leanpath{VOA/\allowbreak Modules/\allowbreak ZhuVermaAnArchLUGraded.lean} and
\leanpath{VOA/\allowbreak Modules/\allowbreak ZhuVermaAnArchZhu534.lean}.
\begin{leancode}
structure GradedLWWeakModule (W : VertexOperatorAlgebra 𝒱) (M : Type v) [AddCommGroup M]
    [Module ℂ M] (h : ℂ) extends WeakModule W M where
  /-- The complex weight grading `M = ⊕_α M_α`. -/
  weight : ℂ → Submodule ℂ M
  /-- The grading is a genuine internal direct sum. -/
  isInternal : DirectSum.IsInternal weight
  /-- **5.3.1c.** `M_α` is the `L₀`-eigenspace of eigenvalue `α`. -/
  L0_weight : ∀ (α : ℂ) (x : M), x ∈ weight α → mode (YM W.conformalVector) 1 x = α • x
  /-- Bounded below in the real part of the weight. -/
  boundedBelow : ∃ B : ℝ, ∀ α : ℂ, weight α ≠ ⊥ → B ≤ α.re
  /-- `h` is a lowest weight: every nonzero weight space sits at `h + ℕ`. -/
  isLow : ∀ α : ℂ, weight α ≠ ⊥ → ∃ n : ℕ, α = h + n
  /-- The lowest weight is *realized*: the bottom layer is nonzero. -/
  realized : weight h ≠ ⊥
def IsIrreducible (N : GradedLWWeakModule W M h) : Prop :=
  Nontrivial M ∧ ∀ Q : WeakModule.Submodule' N.toWeakModule, Q.carrier = ⊥ ∨ Q.carrier = ⊤
def FiniteWeights (N : GradedLWWeakModule W M hw) : Prop :=
  ∀ α : ℂ, FiniteDimensional ℂ (N.weight α)
def toOrdVModule (N : GradedLWWeakModule W M hw) (hfin : N.FiniteWeights) : OrdVModule W M where
  toWeakModule := N.toWeakModule
  weight := N.weight
  isInternal := N.isInternal
  L0_weight := N.L0_weight
  finrank := hfin
  boundedBelow := N.boundedBelow
def toLWModule (N : GradedLWWeakModule W M hw) (hfin : N.FiniteWeights) : LWModule W where
  carrier := M
  isMod := N.toOrdVModule hfin
  h := hw
  isLow := N.isLow
  realized := N.realized
\end{leancode}

\begin{dictionary}
\lean{mode (YM W.conformalVector) 1 x} is $\omega_{(1)}x = L_0x$ (recall \lean{mode f n} is the coefficient $f_{-n-1}$, so
\lean{mode (YM u) n} is $u_{(n)}$). \lean{GradedLWWeakModule} is \emph{exactly} \lean{OrdVModule} plus lowest weight
\lean{h} with all weights in $h+\N$ and $M_h\neq0$, \emph{minus} the requirement that weight spaces be finite-dimensional. The
reason for the omission is substantive: the gate below needs $M(U)$ for the regular module $U = A(V)$, whose bottom is
infinite-dimensional in general. \lean{FiniteWeights} restores finiteness, and \lean{toOrdVModule}/\lean{toLWModule} convert
back; these conversions change nothing else (every field is copied), so a \lean{GradedLWWeakModule} with \lean{FiniteWeights} is
literally an \lean{LWModule}. \lean{IsIrreducible} is the same notion as \lean{OrdVModule.IsIrreducible}, and
\lean{toLWModule_isIrreducibleLW} transports it (its proof is the identity).
\end{dictionary}

\begin{physics}
A graded lowest-weight weak module is a sector whose energy levels may be infinitely degenerate --- harmless for building
sectors, but not a ``physical'' sector with a well-defined character. The finiteness of $L_0$-eigenspaces is what makes the
partition function $\tr q^{L_0-c/24}$ meaningful.
\end{physics}

\subsection{The gate: \texorpdfstring{$O_{0,0}(V) = O(V)$}{O\_\{0,0\}(V)=O(V)} from \texttt{Pr}}

\begin{theorem}[Survival of the bottom; the gate]\label{l13:thm:gate}
Let \lean{Pr} be a bottom projection for the parabolic representation of the regular $A(V)$-module. Then:
\begin{enumerate}
\item (Dong--Li--Mason, Theorem~4.1 at $n=0$.) For every $A(V)$-module $U$ with a bottom projection, $U$ embeds into the
  weak module $\lthirteenMbar_0(U) = M_0(U)/(\text{defects})$ as its bottom.
\item (Dong--Jiang, Proposition~4.6 at $(0,0)$.) $O_{0,0}(V) = O(V)$; hence $A_{0,0}(V) = A(V)$, and for every $A(V)$-module $U$
  the bottom level of $M(U)$ is $M(U)(0)\cong U$.
\end{enumerate}
\end{theorem}
\noindent Files \leanpath{VOA/\allowbreak Modules/\allowbreak ZhuVermaAnArchDLM3Survive.lean} and
\leanpath{VOA/\allowbreak Modules/\allowbreak ZhuVermaAnArchDLM3Induced.lean}.
\begin{leancode}
abbrev IndBar (P : ParabolicRep W n U) : Type u :=
  IndMod P ⧸ VHatRep.defectSubmodule (indRho P)
theorem indBarBottom_injective_of_bottomProjection (W : VertexOperatorAlgebra 𝒱) (U : AMod W)
    (Pr : BottomProjection (parabolicRepOfAMod W U)) :
    Function.Injective (indBarBottom (parabolicRepOfAMod W U)) :=
theorem gate_of_pbwDatum (Pr : BottomProjection (parabolicRepOfAMod W (regularAMod W))) :
    OnmFull W 0 0 = O W :=
theorem djAnnEqAn_zero_of_pbwDatum
    (Pr : BottomProjection (parabolicRepOfAMod W (regularAMod W))) : DJAnnEqAn W 0 :=
noncomputable def djCarrierDegZeroEquivOfPbwDatum
    (Pr : BottomProjection (parabolicRepOfAMod W (regularAMod W))) (V : AmMod W 0) :
    DJCarrier.deg W V 0 ≃ₗ[ℂ] V.carrier :=
\end{leancode}

\begin{dictionary}
\textbf{\lean{IndBar}.} \lean{VHatRep.defectSubmodule (indRho P)} is the smallest $\lthirteenVh$-stable subspace of $M_0(U)$
containing every \emph{iterate defect}
$(u_{(\ell)}v)_{(N)}x - \sum_i(-1)^i\binom\ell i\bigl(u_{(\ell-i)}v_{(N+i)}x - (-1)^\ell v_{(\ell+N-i)}u_{(i)}x\bigr)$, i.e.\ the
failure of the module associativity law. The quotient is a weak $V$-module (\lean{indBarWeakModule}, unconditional).
\lean{indBarBottom} is $u\mapsto$ class of $[1\otimes u]$. The library's quotient imposes the defects for \emph{all} vectors, which
is a priori more than Dong--Li--Mason's (they impose them on the generating space); this only makes the survival statement
(1) \emph{stronger}.

\textbf{Item (1)} is \lean{indBarBottom_injective_of_bottomProjection}: its hypothesis is a bottom projection \emph{for that
$U$}; we use it for $U$ the regular module, which is what \lean{Pr} is.

\textbf{Item (2).} \lean{gate_of_pbwDatum} is literally $O_{0,0}(V) = O(V)$, \lean{djAnnEqAn_zero_of_pbwDatum} is
Definition~\ref{l13:def:djann} at $m = 0$, and \lean{djCarrierDegZeroEquivOfPbwDatum Pr V} is a linear isomorphism
$M(V)(0)\cong V$ for every $A_0(V)$-module $V$ (here the bound variable \lean{V} is the module, not the VOA). It is $A(V)$-equivariant
(\lean{djCarrierWeak_bottom_equiv_rho}). All three take \lean{Pr} as the only hypothesis.
\end{dictionary}

\textbf{Proof on paper.} \emph{Step~1 (survival).} Let $U$ have a bottom projection $\pi:M_0(U)\to U$. Following Dong--Li--Mason,
define the radical $J = \{x\in M_0(U): \pi(y\cdot x) = 0\text{ for all }y\in\lthirteenUE(\lthirteenVh)\}$. It is
$\lthirteenVh$-stable by construction, and $J\cap(1\otimes U) = 0$: if $[1\otimes u]\in J$ then $u = \pi([1\otimes u]) = 0$. The heart
of the matter is that every iterate defect lies in $J$ (\lean{defectPairsToZero_of_bottomProjection}); then the defect submodule is
contained in $J$, so it meets the bottom trivially, which is (1). To show that a defect $\delta$ lies in $J$ one must show
$\pi(y\delta) = 0$ for all $y$. Since $\pi$ kills all levels except $0$, only the level-$0$ component matters. Dong--Li--Mason's
Proposition~4.3 shows that for $x = [1\otimes u]$ the level-zero part of the defect is computed by Zhu's relation in $A(V)$ and
vanishes; their Propositions~4.6--4.7 propagate this from the bottom to all of $M_0(U)$, using the commutator formula to move
modes past the defect (the library proves the needed ``mode-shift'' identity for every truncated representation,
\lean{defectModeShift_of_truncated}). \emph{Step~2 (annihilation, Dong--Jiang Lemma~4.1 and Proposition~4.5).} In any weak module with
a non-negative level structure, the zero mode $o(x)$ of $x\in O_{0,0}(V)$ kills the bottom level: for $x\in O'_{0,0}$ this is Zhu's
argument; for the associator and level-mixing summands it is Dong--Jiang's composition law
$o_{n,m}(a*^n_{m,p}b) = o_{n,p}(a)\,o_{p,m}(b)$ on level $m$ (proved in the library as \lean{djCompositionLev_holds}).
\emph{Step~3 (the regular module).} Apply Step~1 to the regular module $U = A(V)$: $A(V)$ embeds as the bottom of the weak module
$\lthirteenMbar_0(A(V))$, with $o(x)$ acting on the bottom as left multiplication by $[x]$. For $x\in O_{0,0}(V)$, Step~2 gives
$0 = o(x)[\lthirteenvac] = [x*\lthirteenvac] = [x]$, i.e.\ $x\in O(V)$. The reverse inclusion is $O(V) = O'_{0,0}(V)\subseteq O_{0,0}(V)$.
\emph{Step~4.} With $A_{0,0}(V) = A(V)$, the balanced tensor $A_{0,0}(V)\otimes_{A(V)}U = A(V)\otimes_{A(V)}U\cong U$. \qed

\begin{status}
\lean{indBarBottom_injective_of_bottomProjection}, \lean{gate_of_pbwDatum}, \lean{djAnnEqAn_zero_of_pbwDatum},
\lean{djCarrierDegZeroEquivOfPbwDatum}: proved, \emph{conditional on} the explicit \lean{Pr}; \texttt{\#print axioms} = standard
three (the explicit hypothesis is not reported by \texttt{\#print axioms}). Instantiated with \lean{voaPBWBottomProjection} (built from the
theorem \lean{VOA.pbw}) they hold for every $V$; the \lean{Pr} hypothesis is therefore always dischargeable.
\end{status}

\begin{physics}
The gate is the statement ``the primaries of the universal Verma-type module are exactly the primary multiplet we started
with''. Its proof is a beautiful piece of reasoning a physicist would recognize: to show that a relation among fields kills all
primaries, exhibit one sector in which the primaries form the \emph{regular representation} of the zero-mode algebra --- the
most faithful possible --- and test the relation there. The PBW datum guarantees that this universal sector does not collapse.
\end{physics}

\begin{definition}[$M(U)$ as a graded lowest-weight weak module]\label{l13:def:djgradedlw}
Let $U$ be a nonzero $A(V)$-module and $h\in\C$ such that $L_0$ acts on $M(U)(n)$ by $h+n$ for all $n$. Then $M(U)$, with
$M(U)_{h+n} = M(U)(n)$ and $M(U)_\alpha = 0$ for $\alpha\notin h+\N$, is a graded lowest-weight weak module of lowest weight $h$
(the bottom is nonzero by the gate).
\end{definition}
\noindent File \leanpath{VOA/\allowbreak Modules/\allowbreak ZhuVermaAnArchDJCarrierWeak.lean}.
\begin{leancode}
noncomputable def djCWeight (W : VertexOperatorAlgebra 𝒱) (U : AmMod W 0) (h : ℂ) (α : ℂ) :
    Submodule ℂ (DJCarrier W U) :=
  ⨆ (n : ℕ) (_ : h + n = α), DJCarrier.deg W U n
noncomputable def djGradedLW_of_pbwDatum
    (Pr : BottomProjection (parabolicRepOfAMod W (regularAMod W))) (U : AmMod W 0)
    (hU : Nontrivial U.carrier) (h : ℂ)
    (hL0 : ∀ (n : ℕ) (x : DJCarrier W U), x ∈ DJCarrier.deg W U n →
      djModeZero W U 1 W.conformalVector x = (h + n) • x) :
    GradedLWWeakModule W (DJCarrier W U) h :=
  djGradedLW W U h hL0 (djCarrier_deg_zero_ne_bot Pr U hU)
\end{leancode}

\begin{dictionary}
\lean{⨆ (n : ℕ) (_ : h + n = α), DJCarrier.deg W U n} is the sum of the levels $n$ with $h+n = \alpha$: the level-$n$ piece is
placed at weight $h+n$. The hypotheses are \lean{Pr} (to know the bottom is $U$, hence nonzero), \lean{hU : Nontrivial U.carrier}
($U\neq0$), and \lean{hL0}: $L_0 = h+n$ on level $n$. The last is a genuine condition on $U$: it holds iff $[\omega]$ acts on $U$
by the scalar $h$ (for a nonzero module); for the bottom $\Omega(P)$ of a lowest-weight module $P$ it is a theorem
(\lean{pHL0}, Section~\ref{l13:sec:LU}).
\end{dictionary}

\begin{status}
Definition; obligations proved given \lean{Pr}, \lean{hU}, \lean{hL0}; standard three axioms.
\end{status}

\begin{physics}
For $M(U)$ to be graded by energy, all primaries must have the same conformal weight $h$, i.e.\ $L_0$ must be a scalar on the
multiplet. Then level $n$ has energy $h+n$.
\end{physics}

%======================================================================
\section{The irreducible quotient \texorpdfstring{$L(U)$}{L(U)} and the uniqueness argument}\label{l13:sec:LU}
%======================================================================

\subsection{The maximal bottom-trivial submodule}

\begin{definition}[$J$ and $L = M/J$]\label{l13:def:maxJ}
Let $M$ be a graded lowest-weight weak module of lowest weight $h$. A submodule $Q$ (mode-stable subspace) is \emph{graded} if
$Q = \bigoplus_\alpha(Q\cap M_\alpha)$, and \emph{bottom-trivial} if $Q\cap M_h = 0$. Let $J(M)$ be the sum of all graded
bottom-trivial submodules and $L(M) = M/J(M)$. For $M = M(U)$ write $J(U)$ and $L(U)$.
\end{definition}
\noindent Files \leanpath{VOA/\allowbreak Modules/\allowbreak ZhuVermaAnArchLUGraded.lean} and
\leanpath{VOA/\allowbreak Modules/\allowbreak ZhuVermaAnArchLU.lean}.
\begin{leancode}
structure BotTrivialGraded (N : GradedLWWeakModule W M h) where
  /-- The underlying weak submodule. -/
  Q : WeakModule.Submodule' N.toWeakModule
  /-- It is graded. -/
  Q_graded : IsGradedSub N Q
  /-- It meets the bottom `M_h` trivially. -/
  Q_disjoint : Q.carrier ⊓ N.weight h = ⊥
noncomputable def maxJ (N : GradedLWWeakModule W M h) :
    WeakModule.Submodule' N.toWeakModule where
  carrier := ⨆ T : BotTrivialGraded N, T.Q.carrier
theorem maxJ_bottom_disjoint (N : GradedLWWeakModule W M h) :
    (maxJ N).carrier ⊓ N.weight h = ⊥ := by
noncomputable def LU : GradedLWWeakModule W (M ⧸ (maxJ N).carrier) h :=
  quotient N (maxJ N) (maxJ_isGraded N) (maxJ_bottom_disjoint N)

variable (Pr : BottomProjection (parabolicRepOfAMod W (regularAMod W))) (U : AmMod W 0)
  (hU : Nontrivial U.carrier) (h : ℂ)
  (hL0 : ∀ (n : ℕ) (x : DJCarrier W U), x ∈ DJCarrier.deg W U n →
    djModeZero W U 1 W.conformalVector x = (h + n) • x)
noncomputable def djJ : WeakModule.Submodule' (djCarrierWeak W U) :=
  maxJ (djGradedLW_of_pbwDatum Pr U hU h hL0)
noncomputable def djLU :
    GradedLWWeakModule W (DJCarrier W U ⧸ (djJ Pr U hU h hL0).carrier) h :=
  LU (djGradedLW_of_pbwDatum Pr U hU h hL0)
\end{leancode}

\begin{dictionary}
\lean{IsGradedSub N Q} is the gradedness of $Q$; \lean{⨆} of submodules is their sum. \lean{maxJ N} is $J(M)$ (the library checks
mode-stability of the sum). \lean{maxJ_bottom_disjoint} is the key point: \emph{$J(M)$ itself is bottom-trivial}. \lean{LU N} is
the quotient $M/J(M)$ with the induced grading, again a graded lowest-weight weak module of lowest weight $h$ (its bottom is
nonzero because $J$ meets $M_h$ trivially). \lean{djJ}, \lean{djLU} are the instances for $M = M(U)$ of
Definition~\ref{l13:def:djgradedlw}; the \lean{variable} line shows the hypotheses they carry: \lean{Pr}, \lean{hU}
($U\neq0$), a weight \lean{h}, and \lean{hL0} ($L_0 = h+n$ on level $n$).

\textbf{Possible trap.} In the literature $J(U)$ is usually defined as the \emph{largest submodule meeting $U$ trivially}, without
the word ``graded''. For an ordinary module the two agree, because every submodule is graded ($L_0$-eigenspace decomposition). For
a graded weak module the library proves the same thing, with no finiteness: every mode-stable subspace is graded
(\lean{submodule_isGraded}, by $L_0$-semisimplicity, using Lagrange interpolation polynomials in $L_0$). So nothing is lost.
\end{dictionary}

\textbf{Proof on paper (that $J$ is bottom-trivial).} Let $x\in J\cap M_h$. Then $x = q_1+\dots+q_r$ with $q_i\in Q_i$, $Q_i$
graded and bottom-trivial. Take the weight-$h$ component of both sides: since $x\in M_h$, the left side is $x$; since $Q_i$ is
graded, the weight-$h$ component of $q_i$ lies in $Q_i\cap M_h = 0$. So $x = 0$. \qed

\begin{status}
\lean{maxJ}, \lean{LU}: definitions; \lean{maxJ_bottom_disjoint}: proved, unconditional (no finiteness). \lean{djJ},
\lean{djLU}: conditional on the displayed \lean{Pr}, \lean{hU}, \lean{hL0}; standard three axioms.
\end{status}

\begin{physics}
$J(U)$ is the space of \emph{null descendants}: states that can never be brought back to the primaries by any field
(``orthogonal to everything'', in the presence of an inner product). Dividing them out produces the irreducible physical
sector $L(U)$; e.g.\ for the Virasoro algebra at $c = 1/2$, $h = 1/16$, the null vector at level $2$ generates part of $J$.
\end{physics}

\begin{theorem}[$L(U)$ is irreducible; its bottom is $U$; Dong--Jiang, Theorem~4.13]\label{l13:thm:LUirr}
If $U$ is a simple $A(V)$-module satisfying the hypotheses of Definition~\ref{l13:def:djgradedlw}, then $L(U)$ is irreducible, and
its bottom is isomorphic to $U$ as an $A(V)$-module.
\end{theorem}
\noindent Files \leanpath{VOA/\allowbreak Modules/\allowbreak ZhuVermaAnArchLUGraded.lean},
\leanpath{VOA/\allowbreak Modules/\allowbreak ZhuVermaAnArchLU.lean} and
\leanpath{VOA/\allowbreak Modules/\allowbreak ZhuVermaAnArchZhu534.lean}.
\begin{leancode}
def GeneratedByBottom : Prop :=
  ∀ Q : WeakModule.Submodule' N.toWeakModule, N.weight h ≤ Q.carrier → Q.carrier = ⊤
def BottomIrreducible : Prop :=
  ∀ S : Submodule ℂ M, S ≤ N.weight h → S ≠ ⊥ →
    (∀ (a : 𝒱) (p : ℤ) (x : M), x ∈ S → N.toWeakModule.act a p x ∈ N.weight h →
      N.toWeakModule.act a p x ∈ S) →
    S = N.weight h
theorem LU_isIrreducible (hgen : GeneratedByBottom N) (hirr : BottomIrreducible N) :
    IsIrreducible (LU N) := by
theorem djLU_isIrreducible (hsimple : IsSimpleAmMod U) :
    IsIrreducible (djLU Pr U hU h hL0) :=
  LU_isIrreducible (djGeneratedByBottom Pr U hU h hL0) (djBottomIrreducible Pr U hU h hL0 hsimple)
noncomputable def djLU_bottom_U : (djLU Pr U hU h hL0).weight h ≃ₗ[ℂ] U.carrier :=
  ((djLU_bottom_equiv Pr U hU h hL0).symm).trans
    ((LinearEquiv.ofEq _ _ (djCWeight_zero U h)).trans (djCarrierWeak_bottom_equiv Pr U))
theorem djLU_isBottomOfIrr (hsimple : IsSimpleAMod U) : IsBottomOfIrr U :=
  isBottomOfIrr_of_irrRealization (djIrrRealization Pr U hU hval hL0 hfin hsimple)
\end{leancode}

\begin{dictionary}
\lean{GeneratedByBottom N}: a submodule containing $M_h$ is everything. \lean{BottomIrreducible N}: a nonzero subspace $S\subseteq
M_h$ stable under every mode that keeps it in $M_h$ is all of $M_h$. \lean{LU_isIrreducible} derives irreducibility of $M/J$ from
these two; for $M(U)$ both are discharged (\lean{djGeneratedByBottom}: Dong--Jiang's ``$M(U)$ is generated by $U$'';
\lean{djBottomIrreducible}: from simplicity of $U$, \lean{IsSimpleAmMod U}). \lean{djLU_bottom_U} is the linear isomorphism
$L(U)_h\cong U$ (it is $A(V)$-equivariant, \lean{djLU_bottomAModEquiv}). \lean{djLU_isBottomOfIrr} (with the section variables
\lean{Pr}, \lean{hU}, \lean{hval}, \lean{hL0} and \lean{hfin : FiniteWeights}) says: under these hypotheses, $U$ lies in the codomain
\lean{IrrAMod W} of the headline. It is the existence half of Zhu's theorem for one $U$, \emph{conditional on finiteness of the
weight spaces of $L(U)$}; it is not used by the headline.

\textbf{Warning on a literature shortcut.} The library's docstring notes that Dong--Li--Mason's own Lemma~4.8 deduces irreducibility
from complete reducibility, valid only for rational $V$; the argument here uses neither rationality nor finiteness.
\end{dictionary}

\textbf{Proof on paper.} A nonzero submodule $\bar Q$ of $L(U)$ pulls back to a submodule $Q\supsetneq J$ of $M(U)$. It is graded
(every submodule is), so if $Q\cap M(U)_h = 0$ then $Q\subseteq J$ by maximality, a contradiction; hence $Q\cap M(U)_h\neq0$. This
intersection is stable under the zero modes, i.e.\ an $A(V)$-submodule of $M(U)_h\cong U$; by simplicity it is all of $M(U)_h$, and
since $M(U)$ is generated by its bottom, $Q = M(U)$, so $\bar Q = L(U)$. \qed

\begin{status}
\lean{LU_isIrreducible}: proved given its two named inputs. \lean{djLU_isIrreducible}, \lean{djLU_bottom_U}: proved, conditional on
\lean{Pr}, \lean{hU}, \lean{hL0} (and \lean{hsimple}); \lean{djLU_isBottomOfIrr}: additionally on \lean{hfin}. Standard three
axioms.
\end{status}

\begin{physics}
Removing the null descendants of an irreducible primary multiplet leaves an irreducible sector whose primaries are exactly that
multiplet: no new primaries are created and none are lost. This is how one constructs, e.g., the minimal-model sectors from
Verma modules.
\end{physics}

\subsection{An irreducible module is \texorpdfstring{$L$}{L} of its bottom}

\begin{definition}[The bottom of $P$ and the universal map]\label{l13:def:phibar}
For a bundled lowest-weight module $P$ let $\Omega(P) = P_h$ with its $A(V)$-action (an $A_0(V)$-module). The \emph{universal
map} $\bar\varphi:M(\Omega(P))\to P$ is the unique linear map that is the inclusion on the bottom and intertwines all modes:
$\bar\varphi([v]\otimes w) = o_{n,0}(v)w$ on level $n$, where $o_{n,0}(v)$ is the mode of $v$ that raises the level by $n$.
\end{definition}
\noindent File \leanpath{VOA/\allowbreak Modules/\allowbreak ZhuVermaAnArchLUUniversalIrr.lean}.
\begin{leancode}
noncomputable def pBotAmMod (P : LWModule W) : AmMod W 0 :=
  (lowestAMod P.isLow).toAmModZero
noncomputable def pPhiBar (P : LWModule W) :
    DJCarrier W (pBotAmMod P) →ₗ[ℂ] P.carrier :=
  oLevUnivMap P.isMod.toWeakModule (pAdmLev P) (pBotAmMod P) (pBotIncl P)
    (pBotIncl_lev P) (pBotIncl_act P)
theorem pPhiBar_intertwine (P : LWModule W) (p : ℤ) (a : 𝒱) (x : DJCarrier W (pBotAmMod P)) :
    pPhiBar P (djModeZero W (pBotAmMod P) p a x)
      = P.isMod.toWeakModule.act a p (pPhiBar P x) :=
theorem pPhiBar_surjective (P : LWModule W) (hP : P.IsIrreducibleLW) :
    Function.Surjective (pPhiBar P) := by
theorem pPhiBarKer_eq_djJ (P : LWModule W)
    (Pr : BottomProjection (parabolicRepOfAMod W (regularAMod W)))
    (hU : Nontrivial (pBotAmMod P).carrier)
    (hL0 : ∀ (n : ℕ) (x : DJCarrier W (pBotAmMod P)), x ∈ DJCarrier.deg W (pBotAmMod P) n →
      djModeZero W (pBotAmMod P) 1 W.conformalVector x = (P.h + n) • x)
    (hP : P.IsIrreducibleLW) :
    (pPhiBarKer P).carrier = (djJ Pr (pBotAmMod P) hU P.h hL0).carrier :=
\end{leancode}

\begin{dictionary}
\lean{lowestAMod P.isLow} is $P_h$ as an \lean{AMod} (Chapter~\ref{ch:L12}); \lean{pBotAmMod P} is it as an $A_0(V)$-module.
\lean{pPhiBar P} is $\bar\varphi$, built by \lean{oLevUnivMap} from the bottom inclusion \lean{pBotIncl P}, the fact that it lands in
level $0$, and its $A(V)$-equivariance; it is defined for \emph{every} $P$, without hypotheses. \lean{pPhiBar_intertwine}: it is
a module map. \lean{pPhiBar_surjective}: onto, if $P$ is irreducible (its image is a nonzero submodule). \lean{pPhiBarKer P} is its
kernel as a submodule, and \lean{pPhiBarKer_eq_djJ} identifies it with $J(\Omega(P))$, under \lean{Pr}, \lean{hU}
($\Omega(P)\neq0$), \lean{hL0} and irreducibility of $P$.
\end{dictionary}

\textbf{Proof on paper.} \emph{Existence of $\bar\varphi$.} On level $n$ define $\bar\varphi([v]\otimes w) = o_{n,0}(v)w$. This
respects $O_{n,0}(V)$ because, by Dong--Jiang's annihilation (Step~2 of Theorem~\ref{l13:thm:gate}, applied in the module $P$ with
its levels $P_{h+k}$), $o_{n,0}(x)$ kills $P_h$ for $x\in O_{n,0}$; it respects the balancing because
$o_{n,0}(x*^n_0a)w = o_{n,0}(x)o(a)w$ (composition law). It intertwines the modes by the same composition law. \emph{Kernel.} The
kernel $K$ is a submodule; on the bottom $\bar\varphi$ is the inclusion $M(\Omega(P))(0)\cong\Omega(P)\hookrightarrow P$ (this uses
the gate, hence \lean{Pr}), so $K$ is bottom-trivial and $K\subseteq J$. Conversely $\bar\varphi(J)$ is a submodule of the irreducible
$P$ meeting $P_h$ trivially (because $J$ is graded and bottom-trivial and $\bar\varphi$ preserves weights), so it is $0$, i.e.\
$J\subseteq K$. (The library phrases this as: $K$ is bottom-trivial and maximal, and $J$ is the unique such submodule,
\lean{djJ_unique}.) \qed

\begin{status}
\lean{pBotAmMod}, \lean{pPhiBar}: definitions, no hypotheses. \lean{pPhiBar_intertwine}: unconditional.
\lean{pPhiBar_surjective}: conditional on irreducibility of $P$. \lean{pPhiBarKer_eq_djJ}: conditional on \lean{Pr}, \lean{hU},
\lean{hL0}, irreducibility. All: standard three axioms.
\end{status}

\begin{physics}
Every sector is a quotient of the universal ``Verma-type'' module of its primaries: map each formal descendant to the actual
descendant state. In an irreducible sector, the kernel of this map is exactly the space of null descendants.
\end{physics}

\begin{theorem}[$L_0$ on $M(\Omega(P))$; finiteness]\label{l13:thm:hl0}
For every bundled lowest-weight module $P$ and every $n$, $L_0$ acts on $M(\Omega(P))(n)$ as $h+n$. If moreover $P$ is
irreducible, then $L(\Omega(P))$ has finite-dimensional weight spaces.
\end{theorem}
\noindent Files \leanpath{VOA/\allowbreak Modules/\allowbreak ZhuVermaAnArchHL0.lean} and
\leanpath{VOA/\allowbreak Modules/\allowbreak ZhuVermaAnArchLUFiniteWeights.lean}.
\begin{leancode}
theorem pHL0 (P : LWModule W) :
    ∀ (n : ℕ) (x : DJCarrier W (pBotAmMod P)), x ∈ DJCarrier.deg W (pBotAmMod P) n →
      djModeZero W (pBotAmMod P) 1 W.conformalVector x = (P.h + n) • x := by
theorem pFiniteWeights (P : LWModule W)
    (Pr : BottomProjection (parabolicRepOfAMod W (regularAMod W)))
    (hU : Nontrivial (pBotAmMod P).carrier)
    (hL0 : ∀ (n : ℕ) (x : DJCarrier W (pBotAmMod P)), x ∈ DJCarrier.deg W (pBotAmMod P) n →
      djModeZero W (pBotAmMod P) 1 W.conformalVector x = (P.h + n) • x)
    (hP : P.IsIrreducibleLW) :
    (djLU Pr (pBotAmMod P) hU P.h hL0).FiniteWeights := by
\end{leancode}

\begin{dictionary}
\lean{pHL0} gives the datum \lean{hL0} of Definition~\ref{l13:def:djgradedlw} for $U = \Omega(P)$, $h = P_h$'s weight. It takes
\lean{Pr} as an argument but does not use it (see the warning). \lean{pFiniteWeights} gives \lean{FiniteWeights} for
$L(\Omega(P))$, so that $L(\Omega(P))$ can be converted into an \lean{LWModule} by \lean{toLWModule}.
\end{dictionary}

\textbf{Proof on paper.} \emph{$L_0$.} By Dong--Jiang's commutator Lemma~2.3 applied to $\omega\in V_2$,
$\omega\,\bar*^n_0\,v - v*^n_0\omega - \Res_z(1+z)Y(\omega,z)v\in O'_{n,0}(V)$, and $\Res_z(1+z)Y(\omega,z)v = L_{-1}v+L_0v$; subtracting
the grading generator $L_{-1}v+(L_0-n)v\in O'_{n,0}$ gives $[\omega\,\bar*^n_0\,v] = [v*^n_0\omega]+n[v]$ in $A_{n,0}(V)$. On
$[v]\otimes w$, the left side is $L_0([v]\otimes w)$; by balancing $[v*^n_0\omega]\otimes w = [v]\otimes[\omega]w = h\,[v]\otimes w$.
So $L_0 = h+n$. \emph{Finiteness.} The first-isomorphism theorem gives a linear isomorphism $L(\Omega(P))\cong P$ (next theorem)
compatible with $L_0$, hence weight space by weight space; $P$'s weight spaces are finite-dimensional. \qed

\begin{warning}
\textbf{Redundant hypothesis (moderate) --- fixed.} The current \lean{pHL0}, quoted above, has no \lean{Pr}. When this chapter
was first written, \lean{pHL0} took \lean{Pr : BottomProjection (parabolicRepOfAMod W (regularAMod W))} but
its proof does not use it; the library's docstring says so (``retains an explicit, unused \ldots{} parameter''), and
\lean{genHL0} in \leanpath{VOA/\allowbreak Modularity/\allowbreak HL0Scoped.lean} (Chapter~\ref{ch:L15}) states the core without it.
A check in Lean against the library (not part of this release), theorem \lean{pHL0_without_Pr}, proves the conclusion of \lean{pHL0} with no \lean{Pr}
(from \lean{genHL0} and \lean{LWModule.toAMod_omegaScalar}); axioms: standard three. Harmless for the headline, which has \lean{Pr}
anyway.
\end{warning}

\begin{status}
\lean{pHL0}: proved, no hypotheses beyond \lean{P}. \lean{pFiniteWeights}: proved, conditional on \lean{Pr}, \lean{hU}, \lean{hL0},
irreducibility. Standard three axioms.
\end{status}

\begin{physics}
All level-$n$ descendants of a primary of weight $h$ have weight $h+n$: the zero mode of the stress tensor counts the level.
Finiteness: in a physical sector each energy level is finitely degenerate.
\end{physics}

\begin{theorem}[$P\cong L(\Omega(P))$; Dong--Jiang, Theorem~4.13]\label{l13:thm:realizes}
Every irreducible bundled lowest-weight module $P$ is isomorphic to $L(\Omega(P))$.
\end{theorem}
\noindent File \leanpath{VOA/\allowbreak Modules/\allowbreak ZhuVermaAnArchLURealizesIrr.lean}.
\begin{leancode}
theorem pRealizesLWIso (P : LWModule W)
    (Pr : BottomProjection (parabolicRepOfAMod W (regularAMod W)))
    (hU : Nontrivial (pBotAmMod P).carrier)
    (hL0 : ∀ (n : ℕ) (x : DJCarrier W (pBotAmMod P)), x ∈ DJCarrier.deg W (pBotAmMod P) n →
      djModeZero W (pBotAmMod P) 1 W.conformalVector x = (P.h + n) • x)
    (hP : P.IsIrreducibleLW)
    (hfin : (djLU Pr (pBotAmMod P) hU P.h hL0).FiniteWeights) :
    LWModule.Iso W ((djLU Pr (pBotAmMod P) hU P.h hL0).toLWModule hfin) P :=
\end{leancode}

\begin{dictionary}
The conclusion is \lean{LWModule.Iso} between $L(\Omega(P))$ (converted by \lean{toLWModule hfin}) and $P$. The hypotheses:
\lean{Pr}; \lean{hU : Nontrivial (pBotAmMod P).carrier} (true for every $P$, since \lean{P.realized}; the library proves it as
\lean{pBotNontrivial}); \lean{hL0} (supplied by \lean{pHL0}); \lean{hP} irreducibility; \lean{hfin} (supplied by
\lean{pFiniteWeights}). So the only non-trivial inputs are \lean{Pr} and irreducibility.
\end{dictionary}

\textbf{Proof on paper.} $\bar\varphi:M(\Omega(P))\to P$ is surjective with kernel $J(\Omega(P))$; the first isomorphism theorem gives
$L(\Omega(P)) = M/J\cong P$, and the isomorphism intertwines the modes because $\bar\varphi$ does. \qed

\begin{status}
Proved, conditional on \lean{Pr} (other hypotheses dischargeable as stated); standard three axioms.
\end{status}

\begin{physics}
Every irreducible sector is \emph{the} irreducible quotient of the universal module of its primaries. There is no room for two
different sectors with the same primaries.
\end{physics}

\subsection{Functoriality and the assembly}

\begin{theorem}[$L(\cdot)$ is invariant under isomorphism of bottoms]\label{l13:thm:functorial}
If $U_1\cong U_2$ as $A(V)$-modules (both nonzero, with $L_0$ acting on level $n$ by $h_i+n$, $h_1 = h_2$), then
$L(U_1)\cong L(U_2)$.
\end{theorem}
\noindent Files \leanpath{VOA/\allowbreak Modules/\allowbreak ZhuVermaAnArchLUFunctorial.lean} and
\leanpath{VOA/\allowbreak Modules/\allowbreak ZhuVermaAnArchLUBijection.lean}.
\begin{leancode}
noncomputable def djFunctorEquiv (U₁ U₂ : AmMod W 0) (e : AmModHom U₁ U₂) :
    DJCarrier W U₁ ≃ₗ[ℂ] DJCarrier W U₂ :=
  DirectSum.congrLinearEquiv (fun n => djFunctorPiece U₁ U₂ e n)
theorem djLUFunctorLWIso' (Pr : BottomProjection (parabolicRepOfAMod W (regularAMod W)))
    {U₁ U₂ : AmMod W 0} (e : AmModHom U₁ U₂) (hU₁ : Nontrivial U₁.carrier)
    (hU₂ : Nontrivial U₂.carrier) (h₁ h₂ : ℂ) (heq : h₁ = h₂)
    (hL0₁ : ∀ (n : ℕ) (x : DJCarrier W U₁), x ∈ DJCarrier.deg W U₁ n →
      djModeZero W U₁ 1 W.conformalVector x = (h₁ + n) • x)
    (hL0₂ : ∀ (n : ℕ) (x : DJCarrier W U₂), x ∈ DJCarrier.deg W U₂ n →
      djModeZero W U₂ 1 W.conformalVector x = (h₂ + n) • x)
    (hfin₁ : (djLU Pr U₁ hU₁ h₁ hL0₁).FiniteWeights)
    (hfin₂ : (djLU Pr U₂ hU₂ h₂ hL0₂).FiniteWeights) :
    LWModule.Iso W ((djLU Pr U₁ hU₁ h₁ hL0₁).toLWModule hfin₁)
      ((djLU Pr U₂ hU₂ h₂ hL0₂).toLWModule hfin₂) := by
theorem pH_eq {P Q : LWModule W} (ψ : AModEquiv W P.toAMod Q.toAMod) : P.h = Q.h := by
\end{leancode}

\begin{dictionary}
\lean{AmModHom U₁ U₂} is (despite its name) an $A_0(V)$-module \emph{isomorphism} datum (a linear equivalence intertwining $\rho$);
\lean{amModHomOfBottomEquiv} produces one from an \lean{AModEquiv} of bottoms. \lean{djFunctorEquiv} is $M(U_1)\cong M(U_2)$,
$\bigoplus_n(\mathrm{id}\otimes e)$; the library proves it intertwines the modes (\lean{djFunctorEquiv_natural}) and maps $J(U_1)$
onto $J(U_2)$ (\lean{djFunctorEquiv_maps_djJ}), so it descends to $L(U_1)\cong L(U_2)$. \lean{pH_eq}: isomorphic bottoms force equal
lowest weights, because $[\omega]$ acts on the bottom by the lowest weight and the bottom is nonzero.
\end{dictionary}

\textbf{Proof on paper.} $\mathrm{id}\otimes e$ respects the balancing relations because $e$ intertwines the $A(V)$-actions, and commutes
with the modes because they act on the first factor. An isomorphism of modules carries the maximal bottom-trivial graded submodule to
the maximal one. \qed

\begin{status}
Proved, conditional on \lean{Pr} and the displayed data; standard three axioms.
\end{status}

\begin{physics}
Relabelling the primaries by an isomorphism of zero-mode representations relabels the whole sector: the sector is a functor of its
primary multiplet.
\end{physics}

\begin{theorem}[Assembly]\label{l13:thm:assembly}
A family $K$ of lowest-weight modules with $K(P)\cong P$ for irreducible $P$, and $K(P)\cong K(Q)$ whenever $P,Q$ are irreducible with
isomorphic bottoms, implies bottom rigidity and hence bijectivity of \lean{forwardMapIrr}. Taking $K(P) = L(\Omega(P))$ proves the
headline.
\end{theorem}
\noindent Files \leanpath{VOA/\allowbreak Modules/\allowbreak ZhuVermaAnArchLURealizesIrr.lean} and
\leanpath{VOA/\allowbreak Modules/\allowbreak ZhuVermaAnArchLUBijection.lean}.
\begin{leancode}
structure DJRealizationFamily (W : VertexOperatorAlgebra 𝒱) where
  /-- The realizing lowest-weight module (e.g. `L(Ω(P)) = (djLU …).toLWModule hfin`). -/
  K : LWModule W → LWModule W
  /-- Every irreducible `P` is isomorphic to `K P` (`realizesIrr`, `pRealizesLWIso`). -/
  realizes : ∀ P : LWModule W, P.IsIrreducibleLW → LWModule.Iso W (K P) P
  /-- `K` depends on `P` only through its bottom iso-class (functoriality of `L(·)`). -/
  isoInvariant : ∀ P Q : LWModule W, P.IsIrreducibleLW → Q.IsIrreducibleLW →
    Nonempty (AModEquiv W P.toAMod Q.toAMod) → LWModule.Iso W (K P) (K Q)
theorem DJRealizationFamily.rigidity (D : DJRealizationFamily W) : IrredBottomRigidity W := by
theorem DJRealizationFamily.bijective (D : DJRealizationFamily W) :
    Function.Bijective (forwardMapIrr W) :=
  forwardMapIrr_bijective_of_rigidity D.rigidity
noncomputable def djK (Pr : BottomProjection (parabolicRepOfAMod W (regularAMod W)))
    (hHL0 : ∀ (P : LWModule W) (n : ℕ) (x : DJCarrier W (pBotAmMod P)),
      x ∈ DJCarrier.deg W (pBotAmMod P) n →
      djModeZero W (pBotAmMod P) 1 W.conformalVector x = (P.h + n) • x)
    (P : LWModule W) : LWModule W := by
  classical
  exact if hP : P.IsIrreducibleLW then
    (djLU Pr (pBotAmMod P) (pBotNontrivial P) P.h (hHL0 P)).toLWModule
      (pFiniteWeights P Pr (pBotNontrivial P) (hHL0 P) hP)
  else P
theorem zhuBijective_of_hHL0
    (Pr : BottomProjection (parabolicRepOfAMod W (regularAMod W)))
    (hHL0 : ∀ (P : LWModule W) (n : ℕ) (x : DJCarrier W (pBotAmMod P)),
      x ∈ DJCarrier.deg W (pBotAmMod P) n →
      djModeZero W (pBotAmMod P) 1 W.conformalVector x = (P.h + n) • x) :
    Function.Bijective (forwardMapIrr W) :=
  (djRealizationFamily Pr hHL0).bijective
\end{leancode}

\begin{dictionary}
\lean{DJRealizationFamily W} is a bundle (data) of a function \lean{K : LWModule W → LWModule W} and two properties; it is
\emph{not} a hypothesis of the headline: \lean{djRealizationFamily Pr hHL0} constructs it, with
\lean{K = djK Pr hHL0}. \lean{djK} sends an irreducible $P$ to $L(\Omega(P))$ and any other $P$ to itself (a ``junk value'', never
inspected because both properties only concern irreducible $P$; \lean{dite}/\lean{if hP : ... then ... else} is Lean's
case split on a proposition, available by classical logic). \lean{zhuBijective_of_hHL0} is the assembly, with the $L_0$ datum
\lean{hHL0} as a hypothesis, and \lean{zhuBijective} discharges \lean{hHL0} by \lean{pHL0}.
\end{dictionary}

\textbf{Proof on paper (of the headline, complete).} Let $P,Q$ be irreducible with $\psi:P_h\cong Q_{h'}$ as $A(V)$-modules. Then
$h = h'$ (\lean{pH_eq}). By Theorem~\ref{l13:thm:realizes}, $P\cong L(\Omega(P))$ and $Q\cong L(\Omega(Q))$; by
Theorem~\ref{l13:thm:functorial} applied to $\psi$, $L(\Omega(P))\cong L(\Omega(Q))$. Hence $P\cong Q$: bottom rigidity. By
\lean{forwardMapIrr_bijective_iff_rigidity}, \lean{forwardMapIrr} is bijective. The ingredients used: the Lie algebra
$\lthirteenVh$ and the Dong--Li--Mason survival theorem (for the regular module, via \lean{Pr}) to get the gate
$A_{0,0}(V) = A(V)$; the Dong--Jiang module $M(U)$; the maximal submodule $J$; the universal map $\bar\varphi$. \qed

\begin{status}
\lean{DJRealizationFamily.rigidity}, \lean{DJRealizationFamily.bijective}: unconditional (given the bundle).
\lean{zhuBijective_of_hHL0}: conditional on \lean{Pr} and \lean{hHL0}. Standard three axioms.
\end{status}

\begin{physics}
The chain $P\cong L(\Omega(P))\cong L(\Omega(Q))\cong Q$ is the mathematical form of the bootstrap statement ``a sector is fixed by its
primaries''. The proof never needs a classification of sectors, rationality, or modular invariance: only the associativity of the
OPE (encoded in the tower) and the PBW structure of the mode algebra.
\end{physics}

%======================================================================
\section{The 298 files by role}\label{l13:sec:files}
%======================================================================

\subsection{How the critical path was determined}

We did not rely on file names or docstrings. A scratch Lean file importing
\leanpath{VOA/\allowbreak Modules/\allowbreak ZhuVermaAnArchLUBijectionPBW.lean} walked, from \lean{zhuBijective_of_voaPBW}, the
constants occurring in the types \emph{and proof terms} of all declarations reachable from it, and recorded the module of each.
Of the $298$ files \leanpath{ZhuVerma*.lean}, exactly $101$ (about $38\,000$ lines) contain a constant used by the headline; the
other $197$ (about $52\,000$ lines) are imported by the umbrella but contribute nothing to it. (Being in the import closure is
weaker: $236$ of the files are imported, transitively, by the headline's file.)

\subsection{The critical path, grouped by the sections of this chapter}

\begin{itemize}
\item \emph{Statement} (Section~\ref{l13:sec:statement}): \leanpath{FaithfulBijection}, \leanpath{Rigidity},
  \leanpath{Zhu534} (surjectivity, the iff with rigidity, \lean{toLWModule}), \leanpath{AModScope} (\lean{OmegaScalar}),
  \leanpath{LUBijection}, \leanpath{LUBijectionPBW}.
\item \emph{The Lie algebra $\lthirteenVh$ and induction} (Section~\ref{l13:sec:vhat}): \leanpath{Vhat}, \leanpath{VhatDescent},
  \leanpath{VhatSkew}, \leanpath{VhatJacobi}, \leanpath{VhatJac2}, \leanpath{VhatJac3}, \leanpath{MUVhat}, \leanpath{VHatModeData},
  \leanpath{VHatRepMU}, \leanpath{MUjacobi}, \leanpath{MUjacobiClose}, \leanpath{MUModuleDataBuild}, \leanpath{MUmod},
  \leanpath{DLM3Deg}, \leanpath{DLM3Induced}, \leanpath{DLM3Admissible}, \leanpath{DLM3Level}, \leanpath{DLM3Zero},
  \leanpath{DLM3Survival}, \leanpath{DLM3PBWSplit}, \leanpath{DLM3PBWIndependence}, \leanpath{DLM3PBWAxiom}.
\item \emph{The tower} (Section~\ref{l13:sec:tower}): \leanpath{An} (=\leanpath{ZhuVermaAn}), \leanpath{AnBimodule},
  \leanpath{AnDescent}, \leanpath{AnDescent4}, \leanpath{AnDescent15}, \leanpath{Arch}, \leanpath{Algebra}, \leanpath{AnAssoc},
  \leanpath{Anm}, \leanpath{AnmFull}, \leanpath{Filtration}, \leanpath{AltVandermonde}, \leanpath{ResSkew}, \leanpath{SkewModOnm},
  \leanpath{SkewGradingGen}, \leanpath{Units}, \leanpath{Rgen}, \leanpath{CollapseAssoc}, \leanpath{LeftGradingResidualFull},
  \leanpath{MUcirc}, \leanpath{MUcross}, \leanpath{MUFullMode}, \leanpath{MUFullRho}, \leanpath{MUFullBracketClose},
  \leanpath{DJ22}--\leanpath{DJ28*} (Dong--Jiang Lemmas~2.2--2.8: the closure of the ideals under the products).
\item \emph{$M(U)$ and the gate} (Section~\ref{l13:sec:MU}): \leanpath{DJCarrier}, \leanpath{DJCarrierGraded}, \leanpath{DJAction},
  \leanpath{DJ48}, \leanpath{DJ48BdryHolds} (Lemma~4.7, commutator), \leanpath{DJ49Star}, \leanpath{DJ51},
  \leanpath{DJCarrierIterate} (Lemma~4.9, iterate), \leanpath{DJCarrierWeak}, \leanpath{DJGeneration}, \leanpath{IterateNonneg},
  \leanpath{GradedCloseDeep}, \leanpath{GradedModeVanish}, \leanpath{BottomSurvival}, \leanpath{BottomAn},
  \leanpath{BottomDefectClose}, \leanpath{FreenessDescent}, \leanpath{DefectDerivation}, \leanpath{DLM3Pairing},
  \leanpath{DLM3Prop43}, \leanpath{DLM3Propagate}, \leanpath{DLM3Survive}, \leanpath{Gate}, \leanpath{GateB1},
  \leanpath{GateB1Ann}, \leanpath{GateAnnN0}, \leanpath{ZeroModeLaw}, \leanpath{UnivMap}, \leanpath{VModQuot}.
\item \emph{$L(U)$ and uniqueness} (Section~\ref{l13:sec:LU}): \leanpath{LUGraded}, \leanpath{LU}, \leanpath{LUUniversal},
  \leanpath{LUUniversalIrr}, \leanpath{LURealizesIrr}, \leanpath{HL0}, \leanpath{LUFiniteWeights}, \leanpath{LUFunctorial},
  \leanpath{LUIsoInvariant}.
\end{itemize}
(Prefix \leanpath{ZhuVermaAnArch} omitted.)

\subsection{Off the critical path}

The $197$ remaining files are earlier or alternative routes, kept in the library. By role:
\begin{itemize}
\item \emph{Earlier constructions of the $\lthirteenVh$-action and of $M(U)$} (roughly $70$ files): the $\otimes_\C$ ``MU track''
  (\leanpath{MU*} except the few on the path, \leanpath{Sd*}, \leanpath{Hole*}, \leanpath{CarrierBridge*},
  \leanpath{ConstructionAUniform}, \leanpath{GenuineRho}, \leanpath{AGenuineRho}, \leanpath{Rho*}, \leanpath{LevelShiftRho},
  \leanpath{ActionNaturalLevelShift}, \leanpath{LayerAction}), the universal property over that carrier (\leanpath{UnivLift},
  \leanpath{Universal*}, \leanpath{Frobenius*}, \leanpath{FunctorEquivGraded}, \leanpath{Functorial}, \leanpath{StructuralClose},
  \leanpath{ScopeClose}, \leanpath{LightPrimitives}, \leanpath{Leaf2}), and the $J$/$L(U)$ layer over \lean{OrdVModule}
  (\leanpath{Lumax*}, \leanpath{L0Semisimple}, \leanpath{KerJ}, \leanpath{VModQuotGraded}, \leanpath{MuEquivClose}), which cannot
  accommodate the infinite-dimensional regular module. All superseded by the balanced, finrank-free \leanpath{DJCarrier} route.
\item \emph{Alternative or partial proofs of the closure identities} (roughly $60$ files): Li's regular-representation route
  (\leanpath{Li*}); the cross and circle towers (\leanpath{Cross*}, \leanpath{CFirstDiff*}, \leanpath{Hinge*}, \leanpath{Vec*},
  \leanpath{ARD*}, \leanpath{Aout*}, \leanpath{HPoleN}, \leanpath{Term1UCore}, \leanpath{VcoreTwins}, \leanpath{HardMemCFirst2},
  \leanpath{CircDiffARD}, \leanpath{DWeightPoleShift}, \leanpath{SkewDirect}, \leanpath{SkewHardMem}); the collapse
  $O''+O'''\subseteq O'$ (\leanpath{Collapse*}, \leanpath{OnmCollapseZero}); off-diagonal and bimodule-composition descents
  (\leanpath{OffDiagDescent}, \leanpath{UpCrossClosure}, \leanpath{UpRightModule}, \leanpath{BimodComp}, \leanpath{BimodDescent*},
  \leanpath{AnmIdeal}, \leanpath{AnmRightDiff}); alternative proofs of Dong--Jiang lemmas (\leanpath{DJ25}, \leanpath{DJ28Base},
  \leanpath{DJ28Induction}, \leanpath{DJ28IntoFull}, \leanpath{DJ28Term*}, \leanpath{DJ47}, \leanpath{DJ49Iterate},
  \leanpath{DJ49Master}, \leanpath{DefectModeShift}, \leanpath{GateB3}); and the $A_n$ ideal and filtration closures
  (\leanpath{AnIdeal}, \leanpath{ABS}, \leanpath{L1*}, \leanpath{Rgen*} except \leanpath{Rgen}, \leanpath{RightUnit},
  \leanpath{FiltrationFinal}, \leanpath{GradedClose}, \leanpath{GradedSpan}, \leanpath{BottomDefectZero},
  \leanpath{BottomLevelDefect}, \leanpath{BottomOrbitClose}, \leanpath{FreenessCircle}). Some of them (\leanpath{AnIdeal}: \lean{anIdeal_holds}) prove true theorems quoted in this chapter
  that the headline simply does not need.
\item \emph{The descent lemmas of the original $A_n$ bimodule plan} (\leanpath{AnDescent2}--\leanpath{AnDescent26} except the
  three on the path; $23$ files): residue identities for the closure of $O_n$ under the products.
\item \emph{Associativity of $A_n(V)$, $n\ge1$} (\leanpath{Assoc}, \leanpath{AssocCore}): reductions of the open target
  \lean{AssocHom W n}; only $n = 0$ is closed.
\item \emph{Scorecard and interface files built on uninhabited bundles} (\leanpath{InducedAssign}, \leanpath{GenBottom},
  \leanpath{GenFromBottomFinal}, \leanpath{GenFromBottomG}, \leanpath{GameTwo}, \leanpath{CritPath}, \leanpath{Bijection},
  \leanpath{Phase4*}, \leanpath{ScorecardAssemble}, \leanpath{FinalWiring}, \leanpath{Downstream}, \leanpath{Stage7Audit},
  \leanpath{TriageAudit}): bijection theorems taking an \lean{InducedAssignment} (or a structure containing one) as hypothesis. Since
  \lean{InducedAssignment W} is uninhabited (\lean{not_nonempty_inducedAssignment}, Chapter~\ref{ch:L12}), all of them are vacuous;
  see the warning below.
\item \raggedright\emph{Satisfiability and boundary files} (\leanpath{DLM3AdjointBottom}, \leanpath{DLM3BottomProjectionBoundary}): the adjoint
  instance of \lean{BottomProjection} and the equivalence \lean{nonempty_bottomProjection_iff_bottom_injective}
  (Section~\ref{l13:sec:vhat}); quoted in this chapter, not used by the headline.
\end{itemize}

\begin{warning}
\textbf{The file named ``critical path'' is not on the critical path, and its headline is vacuous (serious; documentation only).} The module docstring of
\leanpath{VOA/\allowbreak Modules/\allowbreak ZhuVermaAnArchCritPath.lean} begins ``This file \textbf{audits} the dependency graph of
the final Zhu bijection and banks the single cleanest headline conditional theorem'', and calls its hypothesis bundle ``exactly the data
the unconditional $\Phi(V)\leftrightarrow\mathrm{Irr}A(V)$ bijection consumes''. Its statements are
\begin{leancode}
structure CritPathData (W : VertexOperatorAlgebra 𝒱) where
  /-- The lifting datum `U ↦ M̄(U)` (the generalized-Verma realization; bottom survival +
  functoriality).  Built from `MUVHatData` (`ρ`/`VHatModeData`, `MUGradedAction`, grading,
  bottom residual / freeness); no `A_n`-assoc for `n ≥ 1`. -/
  A : InducedAssignment W
  /-- **Ambient irreducibility scope.** Every bundled lowest-weight module is irreducible — the
  standing restriction of Theorem 5.3.4. -/
  hall : InducedAssignment.AllLWIrreducible W
  /-- **The realizations are irreducible** (`M̄(U) = L(U)`). -/
  hirr : A.RealizeIrreducible
  /-- **Generation-from-bottom** — the universal property of `M(U)` (Frobenius reciprocity for
  irreducibles), the injective-half crux. -/
  hgen : InducedAssignment.GeneratedFromBottom W
theorem zhuBijection_critpath {W : VertexOperatorAlgebra 𝒱}
    (A : InducedAssignment W)
    (hall : InducedAssignment.AllLWIrreducible W)
    (hirr : A.RealizeIrreducible)
    (hgen : InducedAssignment.GeneratedFromBottom W) :
    Function.Bijective (forwardMap W) :=
\end{leancode}
(file \leanpath{VOA/\allowbreak Modules/\allowbreak ZhuVermaAnArchCritPath.lean}). Three problems. (i) The first field is an
\lean{InducedAssignment W}, which is uninhabited for every $V$; so \lean{CritPathData W} is empty and \lean{zhuBijection_critpath} can
never be applied. Checked in Lean against the library (the check is not part of this release):
\lean{critPathData_empty : ¬ Nonempty (CritPathData W)}. (ii) Its
conclusion concerns the \emph{unrestricted} \lean{forwardMap} on all lowest-weight modules, which is not injective in general (a
generalized Verma module and its irreducible quotient have the same bottom), so it is only ``true'' because the hypothesis is
contradictory; its \lean{hall : AllLWIrreducible W} field is false for most $V$, as the docstring of
\leanpath{ZhuVermaAnArchFaithfulBijection.lean} itself states. (iii) None of the file's constants is used by \lean{zhuBijective}; the
actual critical path is the one described above. The same vacuity affects \lean{FaithfulScorecard.bijectiveIrr},
\lean{zhuBijectionIrr_scorecard} and \lean{InducedAssignment.forwardMapIrr_bijective} in
\leanpath{ZhuVermaAnArchFaithfulBijection.lean}, whose docstring presents them as conditional on ``genuine deep inputs''; the library
records this only in \leanpath{ZhuVermaAnArchZhu534.lean} (\lean{not_nonempty_faithfulScorecard}).
\end{warning}

%======================================================================
\section{Recorded traps and dead ends}\label{l13:sec:traps}
%======================================================================

Severities are graded on a four-step scale: serious (a statement that is wrong or vacuous), moderate (weaker or stronger than advertised, or a redundant
hypothesis), open gap (a conditional result whose hypothesis is not discharged, or a missing piece), and documentation
only (docstrings and notes):
\begin{enumerate}
\item \lean{zhuBijective}'s codomain is the image of the forward map, so surjectivity is tautological; the theorem is Zhu's
  uniqueness, not the full correspondence (moderate; Section~\ref{l13:sec:statement}; checked in Lean).
\item \lean{Pr : BottomProjection (parabolicRepOfAMod W (regularAMod W))} is discharged for every $V$ by
  \lean{voaPBWBottomProjection} (from the theorem \lean{VOA.pbw}); this gap, open at the snapshot, is closed
  (Section~\ref{l13:sec:vhat}).
\item ``Bottoms of irreducibles are simple'' (\lean{AllBottomRecovers}) was unproved for general $V$ at the snapshot; it is now the
  theorem \lean{allBottomRecovers} (closed; Theorem~\ref{l13:thm:simplebottom-all}). Combined with existence under $C_2$ it gives
  the genuine bijection \lean{zhuCorrespondence_of_C2Cofinite} (Theorem~\ref{l13:thm:zhucorr}). Without $C_2$, existence of an
  ordinary module over each simple $A(V)$-module remains open in the library (open gap).
\item Release wording (release audit, finding F6): \lean{zhuBijective} and \lean{zhuBijective_of_voaPBW} must be announced as the
  \emph{uniqueness half} of Zhu's correspondence; the bijection to announce is \lean{zhuCorrespondence_of_C2Cofinite}. The
  docstring ``Zhu Theorem 5.3.4 (faithful form)'' of \leanpath{ZhuVermaAnArchLUBijectionPBW.lean} overstates (documentation only).
\item \leanpath{ZhuVermaAnArchCritPath.lean}: misnamed, off the critical path, headline vacuous (serious; documentation only; checked in Lean). Same vacuity for the
  \lean{InducedAssignment}-based theorems of \leanpath{ZhuVermaAnArchFaithfulBijection.lean} and the other scorecard files (serious, already
  proved by the library; docstrings stale, documentation only).
\item \lean{pHL0} had an unused hypothesis \lean{Pr} (moderate; checked in Lean); removed in the library since.
\item Stale trust-boundary docstrings in \leanpath{ZhuVermaAnArchDLM3Survival.lean} and \leanpath{ZhuVermaAnArchDLM3Zero.lean} (``\lean{pbw}
  does not cover this'', ``level independence is the PBW step'') (documentation only; Section~\ref{l13:sec:vhat}).
\item \lean{OnFull}'s docstring claims strict containment \lean{On W n} $\subsetneq$ \lean{OnFull W n} for $n\ge1$ without proof (documentation only).
  The known defect of \lean{On} (Chapter~\ref{ch:L12}) does not reach the headline (Section~\ref{l13:sec:tower}).
\item \lean{AmModHom} is an isomorphism datum (a linear \emph{equivalence}), not a homomorphism (naming only; harmless).
\item Associativity of $A_n(V)$ for $n\ge1$ is open (named target \lean{AssocHom}); no headline depends on it.
\end{enumerate}
Dead ends recorded by the library itself and confirmed here: the \lean{InducedAssignment} interface (uninhabited because it demands a
realization of the zero $A(V)$-module); the \lean{OrdVModule}-based construction of $M(U)$ (impossible for the regular module, whose
bottom is infinite-dimensional); the circle-only quotient $V/O'_{n,m}$ as a bimodule (Dong--Jiang's Remark~2.10 problem); the false
global hypothesis \lean{AllLWIrreducible}.

%======================================================================
\section{Referee's checklist}\label{l13:sec:checklist}
%======================================================================

\begingroup\small
\begin{longtable}{p{4.3cm}p{4.4cm}p{3.3cm}p{1.6cm}}
\toprule
Lean declaration & intended meaning & hypotheses; status & check \\
\midrule
\endhead
\lean{zhuBijective} & irreducibles determined by their bottoms ($\mathrm{fwd}$ bijective onto image) & \lean{Pr}; conditional & Thm.~\ref{l13:thm:headline}\\
\lean{zhuBijective_}\allowbreak\lean{of_}\allowbreak\lean{voaPBW} & same & none; unconditional (via theorem \lean{VOA.pbw}) & Thm.~\ref{l13:thm:headline}\\
\lean{forwardMapIrr}, \lean{IrrLWModuleIsoClass}, \lean{IrrAModIsoClass}, \lean{IsBottomOfIrr} & $[P]\mapsto[P_h]$; $\Phi(V)$; image-defined $\lthirteenIrr A(V)$ & definitions & Def.~\ref{l13:def:irrclasses}\\
\lean{forwardMapIrr_}\allowbreak\lean{surjective_}\allowbreak\lean{uncond} & surjectivity & none; \textbf{tautological} & Thm.~\ref{l13:thm:headline}\\
\lean{forwardMapIrr_}\allowbreak\lean{bijective_}\allowbreak\lean{iff_}\allowbreak\lean{rigidity} & bijective $\iff$ uniqueness & none; unconditional & Thm.~\ref{l13:thm:headline}\\
\lean{isBottomOfIrr_}\allowbreak\lean{iff_}\allowbreak\lean{isSimpleAMod}, \lean{simpleAModsRealized_}\allowbreak\lean{of_}\allowbreak\lean{C2Cofinite} & codomain $=$ simple modules; existence of an ordinary module over each simple $A(V)$-module & \lean{AllBottomRecovers} (a theorem); $C_2$, \lean{Pr} (explicit) & §\ref{l13:sec:scope}\\
\lean{allBottomRecovers}, \lean{isSimpleAMod_}\allowbreak\lean{toAMod_of_irreducible} & bottoms of irreducibles are simple (Zhu part (1)) & none; unconditional & Thm.~\ref{l13:thm:simplebottom-all}\\
\lean{SimpleAMod}, \lean{SimpleAModIsoClass}, \lean{IrrLWModule.}\allowbreak\lean{toSimpleAMod}, \lean{forgetClass} & intrinsic simple $A(V)$-modules up to iso; bottom as a simple module; $\lthirteenIrr\mathrm{LW}/{\cong}\to\Phi$ & definitions & Def.~\ref{l13:def:simpleamodclass}\\
\lean{zhuCorrespondence_}\allowbreak\lean{of_C2Cofinite} & \textbf{Zhu's correspondence as one bijection} $\Phi(V)\simeq\mathrm{Simple}A(V)/{\cong}$, $[P]\mapsto[P_h]$ & $C_2$; conditional; via theorem \lean{VOA.pbw} & Thm.~\ref{l13:thm:zhucorr}\\
\lean{VHat.lieRing} & $\lthirteenVh$ is a Lie algebra & none; unconditional & Def.~\ref{l13:def:vhat}\\
\lean{Pn}, \lean{ParabolicRep}, \lean{parabolicRepOfAMod}, \lean{regularAMod} & $P_0$; level-$0$ datum of an $A(V)$-module & definitions, obligations proved & Defs.~\ref{l13:def:parabolicrep}, \ref{l13:def:regular}\\
\lean{IndMod}, \lean{indBottom}, \lean{levCover} & $\lthirteenUE(\lthirteenVh)\otimes_{\lthirteenUE(P_n)}U$, bottom, levels & definitions & Def.~\ref{l13:def:indmod}\\
\lean{levCover_}\allowbreak\lean{iSupIndep} & levels independent & none; unconditional & Def.~\ref{l13:def:indmod}\\
\lean{BottomProjection}, \lean{PBWFreeness} & DLM's PBW step & definitions; $\iff$ bottom injective & Def.~\ref{l13:def:bottomprojection}\\
\lean{voaPBWBottomProjection} & PBW gives \lean{Pr} & none; unconditional (via theorem \lean{VOA.pbw}) & Thm.~\ref{l13:thm:voapbw}\\
\lean{adjointBottomProjection} & satisfiability (other module) & none; unconditional & Thm.~\ref{l13:thm:adjoint}\\
\lean{circ}, \lean{On}, \lean{An}, \lean{On_zero_eq_O} & library $O_n$, $A_n$; $O_0 = O$ & definitions; unconditional & Def.~\ref{l13:def:circ}, Thm.~\ref{l13:thm:On0}\\
\lean{OnFull}, \lean{leftProd}, \lean{anIdeal_holds} & DLM's $O_n$, $*_n$, two-sided ideal & unconditional & Def.~\ref{l13:def:onfull}\\
\lean{anmGen}, \lean{Onm}, \lean{OnmFull}, \lean{AnmFull}, \lean{DJAnnEqAn} & Dong--Jiang $*^n_{m,p}$, $O_{n,m}$, $A_{n,m}$ & definitions & Defs.~\ref{l13:def:onm}, \ref{l13:def:djann}\\
\lean{DJPiece}, \lean{DJCarrier}, \lean{djModeZero} & $M(U)$ and its modes & definitions; $m=0$ unconditional & Defs.~\ref{l13:def:djcarrier}, \ref{l13:def:djmode}\\
\lean{djCarrierWeak} & $M(U)$ is a weak module & none; unconditional & Thm.~\ref{l13:thm:djweak}\\
\lean{GradedLWWeakModule}, \lean{toLWModule} & finrank-free lowest-weight module & definitions & Def.~\ref{l13:def:gradedlw}\\
\lean{gate_of_}\allowbreak\lean{pbwDatum}, \lean{djCarrierDegZeroEquivOfPbwDatum} & $O_{0,0} = O$; $M(U)(0)\cong U$ & \lean{Pr} & Thm.~\ref{l13:thm:gate}\\
\lean{maxJ}, \lean{maxJ_}\allowbreak\lean{bottom_}\allowbreak\lean{disjoint}, \lean{djLU} & $J$, $L = M/J$ & unconditional / \lean{Pr}, \lean{hU}, \lean{hL0} & Def.~\ref{l13:def:maxJ}\\
\lean{djLU_}\allowbreak\lean{isIrreducible}, \lean{djLU_}\allowbreak\lean{isBottomOfIrr} & $L(U)$ irreducible with bottom $U$ & \lean{Pr}, \lean{hsimple} (+\lean{hfin}) & Thm.~\ref{l13:thm:LUirr}\\
\lean{pPhiBar}, \lean{pPhiBarKer_}\allowbreak\lean{eq_}\allowbreak\lean{djJ} & universal map, kernel $=J$ & \lean{Pr}, irreducible & Def.~\ref{l13:def:phibar}\\
\lean{pHL0}, \lean{pFiniteWeights} & $L_0 = h+n$; finite weights & \lean{Pr} (for \lean{pFiniteWeights}; \lean{pHL0} now has none) & Thm.~\ref{l13:thm:hl0}\\
\lean{pRealizesLWIso} & $P\cong L(\Omega(P))$ & \lean{Pr}, irreducible & Thm.~\ref{l13:thm:realizes}\\
\lean{djLUFunctorLWIso'}, \lean{pH_eq} & $U_1\cong U_2\Rightarrow L(U_1)\cong L(U_2)$ & \lean{Pr} & Thm.~\ref{l13:thm:functorial}\\
\lean{DJRealizationFamily.}\allowbreak\lean{bijective}, \lean{zhuBijective_}\allowbreak\lean{of_}\allowbreak\lean{hHL0} & assembly & bundle / \lean{Pr}, \lean{hHL0} & Thm.~\ref{l13:thm:assembly}\\
\lean{zhuBijection_}\allowbreak\lean{critpath} & (claims) the bijection & \textbf{vacuous} (checked in Lean) & §\ref{l13:sec:files}\\
\bottomrule
\end{longtable}
\endgroup

\paragraph{Machine record.} All \texttt{\#print axioms} results of this chapter were obtained at the snapshot, not re-run for this edition
(when \lean{VOA.pbw} was an axiom; see the re-probe below) on a scratch file
importing \leanpath{VOA/\allowbreak Modules/\allowbreak ZhuVermaAnArchLUBijectionPBW.lean},
\leanpath{VOA/\allowbreak Modules/\allowbreak ZhuVermaAnArchCritPath.lean},
\leanpath{VOA/\allowbreak Modules/\allowbreak ZhuVermaAnArchDLM3AdjointBottom.lean} and
\leanpath{VOA/\allowbreak Modularity/\allowbreak GNProp8Closure.lean}, compiled with \texttt{lake env lean -DmaxSynthPendingDepth=3}
against the existing build (never a bare \texttt{lake env lean}); no full \texttt{lake build} was run for this chapter. Excerpt:
\begin{plaincode}
'VOA.zhuBijective' depends on axioms: [propext, Classical.choice, Quot.sound]
'VOA.zhuBijective_of_voaPBW' depends on axioms: [propext, Classical.choice, Quot.sound, pbw]
'VOA.VHat.DLM3.voaPBWBottomProjection' depends on axioms: [propext, Classical.choice, Quot.sound, pbw]
'VOA.VHat.DLM3.voaPBWBottomInjective' depends on axioms: [propext, Classical.choice, Quot.sound, pbw]
'VOA.VHat.DLM3.levCover_iSupIndep' depends on axioms: [propext, Classical.choice, Quot.sound]
'VOA.VHat.DLM3.adjointBottomProjection' depends on axioms: [propext, Classical.choice, Quot.sound]
'VOA.VHat.lieRing' depends on axioms: [propext, Classical.choice, Quot.sound]
'VOA.gate_of_pbwDatum' depends on axioms: [propext, Classical.choice, Quot.sound]
'VOA.djCarrierWeak' depends on axioms: [propext, Classical.choice, Quot.sound]
'VOA.djLU_isIrreducible' depends on axioms: [propext, Classical.choice, Quot.sound]
'VOA.pRealizesLWIso' depends on axioms: [propext, Classical.choice, Quot.sound]
'VOA.djLUFunctorLWIso'' depends on axioms: [propext, Classical.choice, Quot.sound]
'VOA.pHL0' depends on axioms: [propext, Classical.choice, Quot.sound]
'VOA.forwardMapIrr_surjective_uncond' depends on axioms: [propext, Classical.choice, Quot.sound]
'VOA.forwardMapIrr_bijective_iff_rigidity' depends on axioms: [propext, Classical.choice, Quot.sound]
'VOA.zhuBijection_critpath' depends on axioms: [propext, Classical.choice, Quot.sound]
'VOA.simpleAModsRealized_of_C2Cofinite' depends on axioms: [propext, Classical.choice, Quot.sound]
\end{plaincode}
{\raggedright The same run covered \lean{zhuBijective_of_hHL0}, \lean{voaPBWFreeness} (\lean{pbw}), \lean{BottomProjection}, \lean{PBWFreeness},
\lean{pbwFreeness_iff_bottom_injective}, \lean{bottomProjectionOfInjective}, \lean{parabolicRepOfAMod}, \lean{regularAMod},
\lean{bracketDescends}, \lean{vhatLeibniz_holds}, \lean{indBarWeakModule}, \lean{indBarBottom_injective_of_bottomProjection},
\lean{defectPairsToZero_of_bottomProjection}, \lean{djPieceMEquiv}, \lean{djLU}, \lean{maxJ_bottom_disjoint}, \lean{pFiniteWeights},
\lean{pPhiBar_surjective}, \lean{pPhiBarKer_eq_djJ}, \lean{DJRealizationFamily.bijective}, \lean{forwardMapIrr},
\lean{not_nonempty_inducedAssignment}, \lean{djLU_isBottomOfIrr}, \lean{On_zero_eq_O}, \lean{circ}, \lean{OnmFull}: all report exactly
the standard three axioms, except those marked \lean{pbw} (at the snapshot; none now). Recall that \texttt{\#print axioms} does not report explicit hypotheses such
as \lean{Pr}: ``standard three'' for a conditional theorem means \emph{axiom-clean, conditional on its displayed hypotheses}.\par}
The
Lean check of this chapter against the library (not part of this release) was compiled the same way:
\begin{plaincode}
'L13Cert.zhu_rigidity' depends on axioms: [propext, Classical.choice, Quot.sound]
'L13Cert.zhu_rigidity_voaPBW' depends on axioms: [propext, Classical.choice, Quot.sound, VOA.pbw]
'L13Cert.bijective_of_rigidity_statement' depends on axioms: [propext, Classical.choice, Quot.sound]
'L13Cert.forwardMapIrr_surjective_free' depends on axioms: [propext, Classical.choice, Quot.sound]
'L13Cert.critPathData_empty' depends on axioms: [propext, Classical.choice, Quot.sound]
'L13Cert.pHL0_without_Pr' depends on axioms: [propext, Classical.choice, Quot.sound]
\end{plaincode}
\emph{Re-probed at a later commit}, after \lean{VOA.pbw} became a theorem (not re-run for this edition):
\lean{zhuBijective_of_voaPBW}, \lean{voaPBWBottomInjective}, \lean{voaPBWFreeness}, \lean{voaPBWBottomProjection} and the
theorem \lean{L13Cert.zhu_rigidity_voaPBW} of that check all report \lean{[propext, Classical.choice, Quot.sound]}.

\paragraph{Release audit 2026-10-27: RELEASE\_AUDIT\_ZHU.md.} The independent release audit
(\texttt{Notes/\allowbreak Textbook/\allowbreak referee/\allowbreak RELEASE\_AUDIT\_ZHU.md}, findings F6 and F7) confirmed that
\lean{zhuBijective_of_voaPBW} has the standard three axioms and that its content is injectivity only, and that
\lean{simpleAModsRealized_of_C2Cofinite} takes \lean{Pr} and $C_2$ explicitly. For this revision, at an earlier commit (not re-run for this edition), a
scratch file was run from the repository root with \texttt{lake env lean -DmaxSynthPendingDepth=3} against the oleans built
2026-10-25 (the new file \leanpath{ZhuModularityCanonical.lean} had no olean; its text was included verbatim in the scratch file
and elaborated there):
\begin{plaincode}
'VOA.zhuCorrespondence_of_C2Cofinite' depends on axioms: [propext, Classical.choice, Quot.sound]
'VOA.IrrLWModule.toSimpleAMod' depends on axioms: [propext, Classical.choice, Quot.sound]
'VOA.Huang.LWBottomRecovery.allBottomRecovers' depends on axioms: [propext, Classical.choice, Quot.sound]
'VOA.Huang.LWBottomRecovery.isSimpleAMod_toAMod_of_irreducible' depends on axioms: [propext, Classical.choice, Quot.sound]
'VOA.isBottomOfIrr_iff_isSimpleAMod' depends on axioms: [propext, Classical.choice, Quot.sound]
'VOA.simpleAModsRealized_of_C2Cofinite' depends on axioms: [propext, Classical.choice, Quot.sound]
'VOA.zhuBijective_of_voaPBW' depends on axioms: [propext, Classical.choice, Quot.sound]
\end{plaincode}
No \texttt{lake build} was run for this revision; per the release plan the umbrellas were green when these theorems were
merged, and nothing here counts as released until an aggregate build containing them is green at the release commit.

\paragraph{Second release revision 2026-10-28.} No statement of this chapter changed; only the introduction records the new
consumer of its machinery. The \texttt{\#print axioms} checks of \lean{VOA.Huang.ZhuSemisimple.}\allowbreak\lean{isSemisimpleRing_}\allowbreak\lean{A_of_}\allowbreak\lean{admissiblyRational}
and of the new modularity theorems were run at an earlier commit (standard three axioms; not re-run for this edition) and are recorded in
Chapter~\ref{ch:L15}. No \texttt{lake build} was run.

%======================================================================
\section{Notes and references}\label{l13:sec:notes}
%======================================================================

\textbf{Zhu's algebra and the correspondence.} Y.~Zhu, ``Modular invariance of characters of vertex operator algebras'',
\emph{J. Amer.\ Math.\ Soc.} 9 (1996), 237--302: $A(V)$ (§2.1), the correspondence between irreducible $\N$-gradable modules and
irreducible $A(V)$-modules (Theorems~2.2.1, 2.2.2), proved by a ``generalized Verma module'' constructed via the Lie algebra of modes.
I.~B.~Frenkel and Y.~Zhu, ``Vertex operator algebras associated to representations of affine and Virasoro algebras'',
\emph{Duke Math.\ J.} 66 (1992), 123--168: the bimodules $A(M)$ and the Lie algebra $V\otimes\C[t,t^{-1}]/\operatorname{Im}D$ (also in
R.~E.~Borcherds, ``Vertex algebras, Kac--Moody algebras, and the Monster'', \emph{PNAS} 83 (1986), 3068--3071). T.~Gannon,
\emph{Moonshine beyond the Monster}, Cambridge University Press (2006), §5.3.2, Theorem~5.3.4 (the form cited in the library's
docstrings).

\textbf{Higher Zhu algebras and induced modules.} C.~Dong, H.~Li and G.~Mason, ``Vertex operator algebras and associative algebras'',
\emph{J. Algebra} 206 (1998), 67--96, arXiv:q-alg/9612010: $O_n(V)$ (2.1), $*_n$, Lemma~2.2, Proposition~3.1 (the Lie map
$\lthirteenVh(0)\to A_n(V)$), the induced module $M_n(U)$ (4.1), its PBW step, the radical and Theorem~4.1. C.~Dong, H.~Li and
G.~Mason, ``Twisted representations of vertex operator algebras'', \emph{Math.\ Ann.} 310 (1998), 571--600, arXiv:q-alg/9509005
(§6: the $n=0$ induced module; Proposition~6.1 used for the defect mode shift). C.~Dong and C.~Jiang, ``Bimodules associated to vertex
operator algebras'', \emph{Math.\ Z.} 259 (2008), 799--826, arXiv:math/0601626: $A_{n,m}(V)$, the full ideal $O'+O''+O'''$
(§2, Remark~2.10), the module $M(U)$ (§4), Lemmas~4.1, 4.7, 4.9, Proposition~4.6, Theorems~4.10 and~4.13.

\textbf{Related.} H.~Li, ``Local systems of vertex operators, vertex superalgebras and modules'', \emph{J. Pure Appl.\ Algebra} 109 (1996),
143--195, arXiv:hep-th/9406185, and ``Representation theory and tensor product theory for vertex operator algebras'', PhD thesis,
Rutgers (1994), arXiv:hep-th/9406211 (the Lie algebra of modes and induced modules, §2). J.~Lepowsky and H.~Li, \emph{Introduction to
Vertex Operator Algebras and Their Representations}, Birkh\"auser (2004), Ch.~4--5. M.~R.~Gaberdiel and A.~Neitzke, ``Rationality,
quasirationality and finite W-algebras'', \emph{Comm.\ Math.\ Phys.} 238 (2003), 305--331 (Proposition~8, behind the $C_2$ existence
result of Chapter~\ref{ch:L15}). Poincar\'e--Birkhoff--Witt for arbitrary Lie algebras over a field: N.~Jacobson, \emph{Lie Algebras},
Interscience (1962), Ch.~V; G.~M.~Bergman, ``The diamond lemma for ring theory'', \emph{Adv.\ Math.} 29 (1978), 178--218.

\textbf{Physics.} The bootstrap principle that a sector is determined by its primaries and the role of null vectors: A.~A.~Belavin,
A.~M.~Polyakov and A.~B.~Zamolodchikov, ``Infinite conformal symmetry in two-dimensional quantum field theory'', \emph{Nucl.\ Phys.\
B} 241 (1984), 333--380; P.~Di~Francesco, P.~Mathieu and D.~S\'en\'echal, \emph{Conformal Field Theory}, Springer (1997), Ch.~7--8;
G.~Moore and N.~Seiberg, ``Classical and quantum conformal field theory'', \emph{Comm.\ Math.\ Phys.} 123 (1989), 177--254.

\textbf{Project records.} The trust boundary: \leanpath{VOA/\allowbreak Foundations/\allowbreak AffineLie.lean} and
\leanpath{VOA/\allowbreak Modules/\allowbreak ZhuVermaAnArchDLM3PBWAxiom.lean}; the refutation of \lean{InducedAssignment} and the honest
reading of surjectivity: \leanpath{VOA/\allowbreak Modules/\allowbreak ZhuVermaAnArchZhu534.lean}.

\endgroup
